% QCDM2: current standalone source.
\documentclass[11pt,letterpaper,reqno]{amsart}
\makeatletter
\NeedsTeXFormat{LaTeX2e}[2020/10/01]

\RequirePackage[T1]{fontenc}
\RequirePackage{lmodern}
\RequirePackage{microtype}
\RequirePackage[letterpaper,margin=1in,headheight=15pt]{geometry}
\RequirePackage{amsmath,amsthm,amssymb,mathtools}
\RequirePackage{enumitem}
\RequirePackage{longtable}
\RequirePackage{perpage}
\RequirePackage{xcolor}
\definecolor{manuscriptred}{rgb}{0.6,0.1,0.1}
\definecolor{manuscriptblue}{rgb}{0.1,0.3,0.5}
\definecolor{manuscriptmagenta}{HTML}{60445C}
\RequirePackage[breaklinks=true,colorlinks=true,
  linkcolor=manuscriptred,citecolor=manuscriptblue,urlcolor=manuscriptmagenta,
  filecolor=manuscriptmagenta]{hyperref}
\RequirePackage{cleveref}
\hypersetup{pdfauthor={}}

\setcounter{tocdepth}{1}
\renewcommand{\thefigure}{\Roman{figure}}
\renewcommand{\thetable}{\Roman{table}}
\MakePerPage{footnote}
\renewcommand{\thefootnote}{%
  \ifnum\value{footnote}>9\arabic{footnote}\else\fnsymbol{footnote}\fi}

\theoremstyle{definition}
\newtheorem{ms@theorem}{Theorem}
\renewcommand{\thems@theorem}{\Alph{ms@theorem}}
\newtheorem*{ms@theoremstar}{Theorem}
\newtheorem*{ms@theoremzero}{Theorem 0}
\newtheorem{ms@proposition}{Proposition}[section]
\newtheorem{ms@lemma}[ms@proposition]{Lemma}
\newtheorem{ms@definition}[ms@proposition]{Definition}
\newtheorem{ms@corollary}[ms@proposition]{Corollary}
\newtheorem*{ms@corollarystar}{Corollary}

\newcommand{\ms@majorstart}{%
  \pushQED{\qed}\renewcommand{\qedsymbol}{$\blacksquare\square$}}
\newcommand{\ms@minorstart}{%
  \pushQED{\qed}\renewcommand{\qedsymbol}{$\blacksquare$}}
\newenvironment{theorem}[1][]{\ms@majorstart\ms@theorem[#1]}
  {\popQED\endms@theorem}
\newenvironment{theorem*}[1][]{\ms@majorstart\ms@theoremstar[#1]}
  {\popQED\endms@theoremstar}
\newenvironment{theoremzero}{\ms@majorstart\ms@theoremzero}
  {\popQED\endms@theoremzero}
\newenvironment{proposition}[1][]{\ms@minorstart\ms@proposition[#1]}
  {\popQED\endms@proposition}
\newenvironment{lemma}[1][]{\ms@minorstart\ms@lemma[#1]}
  {\popQED\endms@lemma}
\newenvironment{definition}[1][]{\ms@minorstart\ms@definition[#1]}
  {\popQED\endms@definition}
\newenvironment{corollary}[1][]{\ms@majorstart\ms@corollary[#1]}
  {\popQED\endms@corollary}
\newenvironment{corollary*}[1][]{\ms@majorstart\ms@corollarystar[#1]}
  {\popQED\endms@corollarystar}

\theoremstyle{remark}
\newtheorem{remark}[ms@proposition]{Remark}
\newtheorem*{remark*}{Remark}

\crefname{ms@theorem}{Thm.}{Thms.}
\Crefname{ms@theorem}{Thm.}{Thms.}
\crefname{ms@proposition}{Prop.}{Props.}
\Crefname{ms@proposition}{Prop.}{Props.}
\crefname{ms@lemma}{Lem.}{Lems.}
\Crefname{ms@lemma}{Lem.}{Lems.}
\crefname{ms@definition}{Def.}{Defs.}
\Crefname{ms@definition}{Def.}{Defs.}
\crefname{ms@corollary}{Cor.}{Cors.}
\Crefname{ms@corollary}{Cor.}{Cors.}
\crefname{remark}{Rem.}{Rems.}
\Crefname{remark}{Rem.}{Rems.}

\makeatother


\usepackage{mathrsfs}
\title[Further properties of massive SU(2) theory]{Further properties of three-dimensional SU(2) gauge theory with massive fermions}
\date{}
\hypersetup{pdftitle={Further properties of three-dimensional SU(2) gauge theory with massive fermions},pdfauthor={}}
\begin{document}
\begin{abstract}
We study further properties of the massive fundamental SU(2)
theory in three dimensions. A bound for the mixed dependence of
the vacuum energy on fermion sources gives a uniform variance
estimate for conditional loops at coinciding supports.
It also gives quantitative comparisons between gauge laws
at different masses and mutual absolute continuity of their
simultaneous limits on a fixed physical torus. We enlarge
the scalar meson family to include diquarks, construct its
coordinate-reflection Hilbert space, and prove lower energy
bounds in the charged sectors. We also obtain decay of
parity-odd meson correlations and the short-distance two-point
law for conserved currents. On a fixed physical torus,
subsequences give action-density and meson correlations
holomorphic in coupling and mass away from the total diagonal.
At fixed real coupling, the separated correlations continue
meromorphically in the mass, with a quantitative zero-free
neighbourhood of a longer positive real interval.
The construction paper supplies
the chamber model and its renormalization-group estimates; the
additional source estimates and limit arguments are proved here.
\end{abstract}
\maketitle
\tableofcontents
\section{Introduction}\label{sec:qcdm2-introduction}

We study the massive fundamental SU(2) theory constructed in
\cite{QCDM1}. Its regularization uses the chamber Dirac operator
with Wilson parameter $r=1/2$. Fix $N_f\ge1$ and a physical mass
window $0<m_-\le m_+<1/2$. At spacing $a=L^{-K}$ the bare
parameters are $\beta=\mathsf b/a$ and $m_0=a\mathsf m$.
The gauge measure on a torus is
\[
 \nu_K(dU)=Z_K^{-1}e^{-\beta A(U)}
              |\det(D_+[U]+m_0)|^{2N_f}\,dU,
 \qquad A(U)=\sum_p(1-\tfrac12\operatorname{tr}U_p).
\]
Here $D_+[U]$ is one internal block of the massless chamber
operator for one flavour.

\begin{theoremzero}
For sufficiently weak gauge coupling and masses in a fixed
compact positive interval, conditional fermion loops have
cutoff-independent variance bounds, including at coinciding
supports. On a fixed physical torus, gauge laws at two masses
that converge through the same observables and cutoff sequence
have mutually absolutely continuous limits.

The scalar-bilinear theory extends to diquarks and has lower
energy bounds in its charged sectors. The conserved currents
have a definite short-distance two-point law. On a fixed torus,
subsequences give correlations holomorphic in coupling and mass
off the total diagonal. At fixed real coupling, the separated
correlations continue meromorphically in mass and are holomorphic
on a quantitative neighbourhood of a longer positive real interval.
\end{theoremzero}
The precise statements are \cref{thm:qcdm2-main,C-406-11},
the results of \cref{sec:qcd2-bilinears,sec:qcd2-charged-energy,sec:qcdm2-currents},
and \cref{C-412-55,C-404-60}.

At fixed $K$, an infinite-lattice
state means any weak limit of these measures as the torus
side tends to infinity through the sides permitted by
\cite[Prop. 10.20]{QCDM1}. Such states exist by
\cite[Prop. 2.5]{QCDM1}. For a bounded compactly supported
function $u$ put
\[
 \mathsf G_+=a(D_+[U]+a\mathsf m)^{-1},\qquad
 \mathsf L(u)=\operatorname{Tr}(\mathsf G_+\hat u).
\]
The hat denotes multiplication by $u$ at the site of each
chamber. The trace is over chambers, spin and colour. It is
finite at each cutoff, even on the infinite lattice.

\begin{theorem}[Fluctuations and changes of mass]\label{thm:qcdm2-main}
Choose the odd integer $L\ge13$ and the auxiliary constants
as in \cite[Prop. 10.20]{QCDM1}. There is
$\mathsf b_1<\infty$ such that, for every fixed
$\mathsf b\ge\mathsf b_1$ and every $R<\infty$, there is a
constant $C_R<\infty$ with
\[
 \mathsf E_\nu\bigl|\mathsf L(u)-\mathsf E_\nu\mathsf L(u)\bigr|^2
 \le C_R\|u\|_\infty^2.
\]
This holds for every $K\ge1$, every
$\mathsf m\in[m_-,m_+]$, every infinite-lattice torus-limit
state $\nu$, and every bounded real $u$ whose support lies in
a ball of radius $R$. The constant is independent of the
location of the ball and the chosen state. On a fixed allowed
physical torus the same assertion holds with a constant
depending on its volume.


On a fixed physical torus, fix $0<g\le\gamma$,
$\ell\in\mathfrak L(g)$ and an allowed coupling label
$\mathsf b$, as in \cref{sec:qcdm2-setting}.
Let $f_K$ be Borel maps from the cutoff configuration spaces
to a fixed Polish space. If their laws at two masses
$\mathsf m,\mathsf m'\in[m_-,m_+]$ converge along a common
cutoff sequence to $P,Q$, then $P$ and $Q$ are mutually
absolutely continuous and
\[
 \operatorname{TV}(P,Q)\le C_{g,\ell}|\mathsf m'-\mathsf m|.
\]
The constant is independent of the cutoff, the two masses
and $\mathsf b$. The cutoff measure at each mass uses that
mass in the determinant defining $\nu_K$.
\end{theorem}

The trace need not have a finite limit without subtraction.
The theorem controls its fluctuation at coinciding supports.
Its proof uses a determinant with a complex source in one
internal block and the conjugate source in the other.
The parts of the vacuum energy depending on just one source
are holomorphic. Removing them leaves a mixed energy bounded
uniformly in the cutoff. The first coefficient of the resulting
holomorphic determinant gives the variance bound by Parseval's
identity. \Cref{C-415-17} gives the explicit constants and
proves the variance assertion of \cref{thm:qcdm2-main}.
For a constant source the determinant changes the mass.
The same circle bound controls the density for a short
change of mass. Chaining these bounds across the mass
interval gives the comparison of limiting gauge laws in
\cref{C-406-11}.

This estimate also enters the analysis of parity-odd scalar
bilinears. We treat those bilinears together with mesons,
diquarks and antidiquarks. Their separated correlations have
a common continuum subsequence and a compact symplectic
flavour action by $G=\mathrm{USp}(2N_f)^2$.
Coordinate reflection positivity gives an
enlarged Hilbert space and a time-translation semigroup.
Bounds obtained by counting fermion propagators give a lower
bound $H\ge|q|\mathsf m$ on quark number $q$, and
$H\ge2\mathsf m$ on the complement of the $G$-invariant
vectors. Both sectors $q=2$ and $q=-2$ are nonzero. These claims
and the distinction between coordinate reflection positivity
and full rotation covariance are stated in
\cref{sec:qcd2-bilinears,sec:qcd2-charged-energy}.

We also prove a short-distance two-point law for the conserved
currents. The current is a point-split lattice bilinear. Its
comparison with a site bilinear requires a source that acts
on spin, as well as on internal and flavour indices. The
source construction extends to those matrices; an energy
estimate for the propagator then controls the point splitting.
The result concerns current two-point functions and lattice
reflection-positive current states. It does not assume a
joint reconstruction of all the currents with the scalar
fields of \cite{QCDM1}; see \cref{sec:qcdm2-currents}.

Finally, for the original action-density and meson family
on a fixed physical torus, we allow supports with
empty total intersection. Individual pairs of supports may
overlap. A polymer contributing to a mixed derivative is
either large or far from at least one insertion; in both
cases the local-energy estimate supplies decay. A common
subsequence then gives correlations off the total diagonal,
holomorphic in the coupling and mass near the permitted real
parameter window. The neighbourhood depends on the physical
volume. This is a separate fixed-volume conclusion from the
volume-first construction used for the charged sectors.

At fixed real coupling we continue the separated correlations
meromorphically in the mass. The same subsequence and discrete
exceptional set work for all their orders and test functions.
A quantitative lower bound for the normalized partition
function gives a zero-free neighbourhood of a longer real
mass interval, with width at least
$c/[\ell^3(1+\log\ell)]$. The real-mass cumulant estimates
remain uniform in the volume. These statements, including
the continued short-distance laws, are proved in
\cref{sec:qcdm2-mass-continuation}.

All statements concern the massive chamber model of \cite{QCDM1}.
We use its construction and the estimates explicitly recalled below,
and give the additional determinant, source and continuation
arguments in the same notation.

\section{Model, sources and parameter ranges}\label{sec:qcdm2-setting}

We recall the definitions needed for the observables and the
estimates. The lattice construction is in
\cite[\S2]{QCDM1}. Spatial coordinates in operator formulas
are in lattice units; supports and distances are in physical
units unless stated otherwise.

\subsection{The chamber operator}

Let $W_3$ be the group of signed permutation matrices.
A place is $v=(x,q)$ with $x\in\mathbb Z^3$ and $q\in W_3$.
Put
\[
 (N_0,N_1,N_2,N_3)=(e_1,e_2-e_1,e_3-e_2,-e_3),
 \qquad t_s=3|N_s|,
\]
and let $w_s$ be reflection in $N_s^\perp$. The adjacent
places and the corresponding unit vectors are
\[
 v\cdot s=(x+\delta_{s0}qe_1,qw_s),\qquad
 u_s(v)=qN_s/|N_s|.
\]
Set $U(v,v\cdot s)=1$ for $s=1,2,3$ and
$U(v,v\cdot0)=U_{\langle x,x+qe_1\rangle}$.
The spinor fibre is
\[
 \mathbb C^2_{\rm s}\otimes\mathbb C^2_\tau
       \otimes\mathbb C^2_{\rm c}\otimes\mathbb C^{N_f}.
\]
The factors denote spin, internal index, colour and flavour.
With the Pauli matrices $\sigma_\mu$ and $\tau_\mu$ write
\[
 \gamma_\mu=\sigma_\mu\otimes\tau_3,\qquad
 \gamma_4=1\otimes\tau_1,\qquad e_\pm=(1\pm\tau_3)/2.
\]
The massless operator is
\begin{equation}\label{eq:qcdm2-Dirac}
 (D[U]f)(v)=\frac12\sum_{s=0}^3t_s f(v)
 -\sum_{s=0}^3t_s\bigl(\tfrac12-\gamma(u_s(v))\bigr)
                   U(v,v\cdot s)f(v\cdot s).
\end{equation}
The two internal blocks satisfy $D_-=D_+^*$ and
$\operatorname{Re}(D_\pm+m_0)\ge m_0$. In particular
$\|\mathsf G_\pm\|\le\mathsf m^{-1}$, where
$\mathsf G_\pm=a(D_\pm+m_0)^{-1}$.

On $T_\ell=(\mathbb R/\ell\mathbb Z)^3$ the links are periodic
and spinors are antiperiodic in the three coordinate directions.
The finite-volume functional is
\[
 \Omega_K(F)=\int dU\int d\bar\psi\,d\psi\,
  e^{-\beta A(U)-\bar\psi(D[U]+m_0)\psi}F,
 \qquad \langle F\rangle_K=\Omega_K(F)/\Omega_K(1).
\]
Haar measures have total mass one. The Berezin integral has
the convention giving $\det(D[U]+m_0)$ at $F=1$.
Conditional expectations at fixed $U$ are the Gaussian
Berezin functional with covariance $(D[U]+m_0)^{-1}$.
On the infinite lattice we integrate this functional against
an infinite-lattice torus-limit gauge state. All polynomial
insertions have finite support at fixed $a$.

\subsection{Physical sources}

Choose Hermitian flavour matrices $t_j$ with
$t_1=N_f^{-1/2}1$ and
$\operatorname{tr}(t_jt_l)=\delta_{jl}$. Set
\[
 (c_0,c_1,c_2,c_3)=(1,i,i,1),\quad
 \Gamma_{a,j}=1_{\rm s}\otimes c_a\tau_a\otimes1_{\rm c}\otimes t_j,
 \quad\mathfrak B=\{\Gamma_{a,j}:0\le a\le3,\ 1\le j\le N_f^2\}.
\]
These matrices have operator norm at most one. The physical
action and meson insertions are
\begin{align}
 O_K(w)&=a^{-1}\sum_p w(x(p))
                     (1-\tfrac12\operatorname{tr}U_p),\label{eq:qcdm2-action}\\
 \Phi_K(J)&=a\sum_x\sum_{v:x_v=x}
                     \bar\psi(v)J(ax)\psi(v),\label{eq:qcdm2-meson}
\end{align}
where $J$ takes values in the complex span of $\mathfrak B$.
Here $x(p)$ is the physical plaquette centre, and each positive
coordinate orientation occurs once. The source partition
function is
\[
 Z_K[w,J]=\Omega_K(e^{O_K(w)+\Phi_K(J)})
 =\int e^{-\beta A(U)+O_K(w)}
     \det(D[U]+m_0-a\hat J)\,dU.
\]
For $\alpha=0$ set $\Phi_{\alpha,K}(f)=O_K(f)$, and for
$\alpha=\Gamma\in\mathfrak B$ set
$\Phi_{\alpha,K}(f)=\Phi_K(f\Gamma)$.
We use
\[
 \|f\|_0=\|f\|_\sharp
      =\max\{\|f\|_\infty,40M\operatorname{Lip}f\},
 \qquad\|f\|_\Gamma=\|f\|_\infty.
\]
The constants $M$, $\varrho$, $\varrho_J$ and $r_M$ belong
to the source estimates of \cite[\S10]{QCDM1}.

Joint cumulants are derivatives of the logarithm of the
normalized source partition function at zero. Their centered
moments are the sums over set partitions with no singleton
blocks. We use the same definition for the enlarged bilinear
family introduced below.

\subsection{Permitted parameters and limits}

Fix the mass window and $N_f$. Choose the constants of
\cite[Prop. 10.20]{QCDM1}, and decrease its upper bound
$\gamma$ on $g$ as specified in
\cref{sec:qcdm2-mixed-proof}. These choices depend on the fixed
mass window and $N_f$, and not on $K$ or the volume. Put
\[
 I_g=[g^{-2}+\bar B,L(g^{-2}+B_\sharp)-\bar B],\qquad
 0<g\le\gamma,\quad\mathsf b\in I_g,
 \quad\mathsf m\in[m_-,m_+].
\]
The allowed sides form the cofinal set
$\mathfrak L(g)=\{2L^m:m\ge m_g\}$ of that proposition.
More explicitly, with $M=L^{m'}$,
\[
 \begin{split}
 X_g&=L(g^{-2}+B_\sharp),\qquad
 R(x)=\min\{L^q:q\ge0,\ L^q\ge(\log x)^{r_c}\},\\
 m_g&=1+m'+\max\{\lceil\log_LR(X_g+\bar B)\rceil,
                         \lceil\log_L C_{\rm s}\rceil\}.
 \end{split}
\]
Here $r_c\ge2$ and $C_{\rm s}\ge2$ are the fixed local-chart
constants of the construction. This condition supplies the
room required by the local estimates.
The intervals $I_g$ cover all sufficiently large
$\mathsf b$: for example take $g^{-2}=\mathsf b-\bar B$ and
increase the lower threshold on $\mathsf b$ so that the upper
endpoint also exceeds $\mathsf b$. This explains the parameter
formulation in \cref{thm:qcdm2-main}.

For infinite-volume conclusions we first take torus limits at
fixed $K$, then let $K\to\infty$. Constants may depend on
$g$ unless a stronger uniformity is stated. For the analytic
results in \cref{sec:analytic-labels}, the physical torus is
instead kept fixed throughout the cutoff limit.


\section{Determinant bounds at coinciding supports}
\label{sec:qcdm2-circle}

The estimates for disjoint supports do not bound the variance of a
loop against itself.  We obtain this bound from a different use of
the fermion determinant.  A source in one species gives a
holomorphic determinant; its conjugate source in the other species
gives the squared absolute value.  After removing the two terms of
the vacuum energy which depend on just one species, the remaining
energy is bounded uniformly in the cutoff.  Cauchy's formula then
bounds the coefficients of the determinant, including the loop.

We retain the setting of \cref{sec:qcdm2-setting}, with the
further choices in \cref{sec:qcdm2-mixed-proof}.
Write $\nu_K$ for the
gauge measure, and let $\mathcal G^{\rm tor}_{K,\beta,m_0}$ be the
infinite-lattice states obtained as limits of these torus measures.
All the constants below are uniform in the admitted data except
where dependence on $g$, a torus side, or a support radius is stated.

The two physical resolvents and the simplest loop are
\[
 \mathsf G_\pm=a(D_\pm[U]+m_0)^{-1},\qquad
 \mathsf G_-=\mathsf G_+^*,\qquad
 \mathsf L(f)=\operatorname{Tr}(\mathsf G_+\hat f).
\]
Here $\hat f$ denotes multiplication by $f$, and the trace includes
the position, spin and colour variables but no flavour variable.
When a source acts on flavour, we extend $\mathsf G_\pm$
by the identity on that space.
The resolvent bound is $\|\mathsf G_\pm\|\le\mathsf m^{-1}$.
On the infinite lattice all functions in a determinant or trace
have finite support; consequently the relevant operator has finite
rank.  Set
\begin{equation}\label{eq:qcdm2-Xi}
 \Xi_\nu^f(z)=\mathsf E_\nu
       |\det(1+z\mathsf G_+\hat f)|^{2N_f}.
\end{equation}
We say that the circle bound holds at $(f,\rho,C)$ if there is a
function $h$, holomorphic near $|z|\le\rho$, such that $h(0)=0$ and
\begin{equation}\label{C-414-10}
 \Xi_\nu^f(z)\le C\exp(2N_f\operatorname{Re}h(z)),
 \qquad |z|=\rho.
\end{equation}
The function $h$ absorbs the mean of the loop.  Its size need not
be uniform in $K$.

\begin{theorem}[Circle bound]\label{C-415-17}
There are constants $C_{\rm mix}$ and $C'_{\mathscr C}$, independent
of the admitted data, for which the following holds.  Put
\[
 \rho_*=2\varrho_J\le\tfrac12{m_-} ,\qquad
 \rho_*\le\frac{r_M}{2L},\qquad
 c_V(g)=C+2+C_P\bigl(1+\log(L(g^{-2}+B_\sharp))\bigr).
\]
The constants on the right are those of the partition bounds in
\cite[Props. 10.14 and 10.17]{QCDM1}.
\begin{enumerate}
\item On a torus of side $\ell$, every bounded real $f$ satisfies
the circle bound with
\[
 \rho=\rho_*/\|f\|_\infty,
 \qquad C=\exp((c_V(g)+C_{\rm mix})\ell^3).
\]
The assertion for $f=0$ is understood without the quotient.
One can take
\begin{equation}\label{eq:qcdm2-pure-energy}
 N_fh_K(z)=E_K(0,-e_+zf1)-E_K(0,0),
\end{equation}
which is holomorphic on $|z|\|f\|_\infty<r_M/L$.

\item For $R>0$, let $i(R)$ be the least integer $i\ge1$ with
$L^i\ge2R+3$, and put
\begin{equation}\label{eq:qcdm2-CR}
 C(R)=\exp\left((c_V(g)+C_{\rm mix})L^{3i(R)}
                 +2C'_{\mathscr C}(1+R^3)\right).
\end{equation}
If $f$ is real, $\|f\|_\infty\le1$, and its support lies in a ball
of radius $R$, the circle bound holds at $(f,\rho_*,C(R))$ on
every admitted torus of side $\ell\ge2L^{i(R)}$ and for every
$\nu\in\mathcal G^{\rm tor}_{K,\beta,m_0}$.

\item For any such measure and any real $u$ supported in that ball,
\begin{equation}\label{eq:qcdm2-loop-variance}
 \operatorname{Var}_\nu\mathsf L(u)
 \le \frac{\|u\|_\infty^2}{\rho_*^2}
       \min\left\{C(R)^{1/N_f},\frac{C(R)-1}{N_f^2}\right\}.
\end{equation}
The same assertion on a fixed torus holds with
$C(R)$ replaced by $\exp((c_V(g)+C_{\rm mix})\ell^3)$.
\end{enumerate}
\end{theorem}

The estimate concerns random variables of the gauge field.  It does
not identify the Grassmann bilinear with a random distribution.
The relation between bilinear correlations and conditional loops
is supplied by the conditional Gaussian identities.

\subsection{Coefficients of a determinant}

\begin{proposition}\label{C-414-11}
If the circle bound holds at $(f,\rho,C)$, then
\begin{align}
 \mathsf E_\nu|\mathsf L(f)-h'(0)|^{2N_f}
       &\le \rho^{-2N_f}C,\label{eq:qcdm2-high-moment}\\
 \mathsf E_\nu|\mathsf L(f)-h'(0)|^2
       &\le \frac{C-1}{N_f^2\rho^2}.
       \label{eq:qcdm2-Parseval-variance}
\end{align}
In particular,
\[
 \operatorname{Var}_\nu\mathsf L(f)
 \le\rho^{-2}\min\{C^{1/N_f},(C-1)/N_f^2\}.
\]
More generally, if $Z_U$ is a random holomorphic polynomial with
$Z_U(0)=1$, and $\mathcal F$ is holomorphic near the closed disc
with $\mathcal F(0)=0$, write
$e^{-\mathcal F(z)}Z_U(z)=\sum_{n\ge0}d_n(U)z^n$.  Then
\begin{equation}\label{eq:qcdm2-Parseval}
 \sum_{n\ge0}\rho^{2n}\mathsf E_\nu|d_n|^2
 =\frac1{2\pi}\int_0^{2\pi}
       e^{-2\operatorname{Re}\mathcal F(\rho e^{i\theta})}
       \mathsf E_\nu|Z_U(\rho e^{i\theta})|^2\,d\theta.
\end{equation}
\end{proposition}

\begin{proof}
Parseval's identity holds at each $U$; Tonelli's theorem permits
integration of its nonnegative terms.  Apply it first to
$Z_U(z)=\det(1+z\mathsf G_+\hat f)^{N_f}$ and
$\mathcal F=N_fh$.  The first two coefficients are
$d_0=1$ and $d_1=N_f(\mathsf L(f)-h'(0))$.  This gives
\eqref{eq:qcdm2-Parseval-variance}.

For the higher moment, use the first coefficient of
$e^{-h(z)}\det(1+z\mathsf G_+\hat f)$ instead.  Cauchy's formula
and Jensen's inequality give
\[
 |\mathsf L(f)-h'(0)|^{2N_f}
 \le\rho^{-2N_f}\frac1{2\pi}\int_0^{2\pi}
 |e^{-h(z)}\det(1+z\mathsf G_+\hat f)|^{2N_f}\,d\theta,
 \qquad z=\rho e^{i\theta}.
\]
Integrate over $\nu$.  Finally the mean minimizes the squared
distance from a constant, and the $L^{2N_f}$ bound implies the
$L^2$ bound with constant $C^{1/N_f}$.
\end{proof}

For the passage to an infinite-lattice state it is preferable to
avoid extracting a convergent family of the functions $h$.
The next observation expresses their existence in terms of the
determinant alone.

\begin{lemma}\label{C-415-16}
Let $\rho\|f\|_\infty<\mathsf m$.  The function
$\log\Xi_\nu^f$ is subharmonic on
$|z|\|f\|_\infty<\mathsf m$, and the circle bound holds at
$(f,\rho,C)$ if and only if
\begin{equation}\label{eq:qcdm2-log-circle}
 \frac1{2\pi}\int_0^{2\pi}
       \log\Xi_\nu^f(\rho e^{i\theta})\,d\theta\le\log C.
\end{equation}
\end{lemma}

\begin{proof}
Put $Z_U(z)=\det(1+z\mathsf G_+\hat f)^{N_f}$.
The resolvent bound makes this determinant nonzero in the stated
disc.  Tilting $\nu$ by $|Z_U(z)|^2/\Xi_\nu^f(z)$ gives a probability
measure $\nu_z$, and differentiation gives
\[
 \partial_z\partial_{\bar z}\log\Xi_\nu^f(z)
 =\mathsf E_{\nu_z}|\partial_z\log Z_U(z)|^2
    -|\mathsf E_{\nu_z}\partial_z\log Z_U(z)|^2\ge0.
\]
The circle bound implies \eqref{eq:qcdm2-log-circle} by the mean
value property for $\operatorname{Re}h$.  Conversely the boundary
function $\log\Xi_\nu^f$ is real analytic.  Let $H$ be holomorphic
near the closed disc, with real part equal to this boundary
function and $H(0)$ equal to its real circle mean $M$.  Such an
$H$ is obtained by its Fourier series.  Then
$h=(H-M)/(2N_f)$ gives the circle bound with $C=e^M$.
\end{proof}

\subsection{The two parts of the vacuum energy}

Write $e_\pm=(1\pm\tau_3)/2$.  A split source is a pair
$J_1=e_+J_1e_+$, $J_2=e_-J_2e_-$.  For a function $F$ of the
source set
\begin{equation}\label{C-415-2}
 D^{\rm mix}F(J_1,J_2)
 =F(J_1+J_2)-F(J_1)-F(J_2)+F(0).
\end{equation}
The vacuum energy $E_K(0,J)$ is the energy defined by the source
construction in \cite[Prop. 10.6]{QCDM1}; put
$E_K^{\rm mix}=D^{\rm mix}E_K(0,\cdot)$.

\begin{proposition}[Mixed vacuum energy]\label{C-415-5}
For split sources with $\|J_b\|_\infty\le\rho_*$,
\begin{equation}\label{eq:qcdm2-mixed-energy}
 |E_K^{\rm mix}(J_1,J_2)|\le C_{\rm mix}\ell^3.
\end{equation}
More precisely, set
\[
 \lambda_j=L^{j-K},\qquad n_j=\ell^3\lambda_j^{-3},\qquad
 \mathfrak l_k=\mathsf t_k^{-1/2}(\log\mathsf t_k)^{-p_0},
 \qquad \mathsf t_k=L^{K-k}\mathsf b,
\]
and retain $\phi_k=\phi(x_k)$ from the source construction.
The energy contribution at level $j$ satisfies, for
$\lambda_j\|J_b\|_\infty<r_M$,
\begin{align}
 \left|\sum_XD^{\rm mix}e_j(X;0,\cdot)(J_1,J_2)\right|
 &\le C_{\rm gm3}M^{-3}n_j
 \left(\phi_{j-1}^3
       +\phi_{j-1}^2e^{-c_{\rm t}\ell/(2\lambda_j)}
       +\mathfrak l_{j-1}^2\right)
 \frac{\lambda_j^2\|J_1\|_\infty\|J_2\|_\infty}{r_M^2}.
 \label{eq:qcdm2-mixed-step}
\end{align}
The bound is on the sum of the energy contributions at a level.
\end{proposition}

Its proof occupies \cref{sec:qcdm2-mixed-proof}.  The three terms
have different origins.  Independent colour averages remove the
product of two linear gauge vertices.  The remaining terms with
fermion legs have three powers of $\phi$, apart from a term caused
by polymers crossing a torus seam.  Two scalar vertices need not
cancel; their contribution is controlled by $\mathfrak l^2$.

\begin{proof}[Proof of \cref{C-415-17}]
For $J=e_+A+e_-B$, block diagonalization gives
\begin{equation}\label{C-415-13}
 \frac{\det(D+m_0+a\hat J)}{\det(D+m_0)}
       =\det(1+\mathsf G_+\hat A)
        \det(1+\mathsf G_-\hat B).
\end{equation}
At $A=zf1$, $B=\bar zf1$, integration gives
\[
 \Xi_{\nu_K}^f(z)
  =\frac{Z[0,-e_+zf1-e_-\bar zf1]}{Z[0,0]}.
\]
The partition bounds in \cite[Prop. 10.14]{QCDM1} bound this quotient
by the exponential of
\[
 \operatorname{Re}\bigl(E_K(0,-e_+zf1-e_-\bar zf1)-E_K(0,0)\bigr)
       +c_V(g)\ell^3.
\]
The two pure-source energies are complex conjugates.  Their sum
is $2N_f\operatorname{Re}h_K(z)$, with $h_K$ as in
\eqref{eq:qcdm2-pure-energy}; the mixed part is bounded by
\cref{C-415-5}.  This proves the fixed-torus assertion.

For the local assertion place the support in a cube of side
$L^{i(R)}$, at distance at least one from its boundary.  The
chessboard estimate in \cite[proof of Prop. 10.17]{QCDM1} reduces the
source quotient to the $N$th root of the quotient for the reflected
source on the whole board, where $\ell^3=NL^{3i(R)}$.
Reflection alternates the source $z$ with its conjugate.  Thus the
upper species source has the form
$zf_{\rm even}+\bar zf_{\rm odd}$, and the lower source is its
conjugate.  Its mixed energy is bounded by $C_{\rm mix}\ell^3$.
The comparison of the energy of separated source copies with the
sum of their energies, by
\cite[Prop. 10.16 and equation (46)]{QCDM1} at separation
one, costs at most
$2NC'_{\mathscr C}(1+R^3)$.  The pure-source terms, after conjugating
the odd copies, give twice the real part of a holomorphic function
which vanishes at zero.  Taking the $N$th root gives
\eqref{eq:qcdm2-CR}.

At fixed $K$ and fixed support the determinants are polynomials of
a common finite degree.  Torus convergence of their coefficients
follows from the convergence of the local Gaussian functionals
in \cite[Prop. 2.5]{QCDM1}; their uniform resolvent bounds also
give uniform convergence on the circle.  Since
$\rho_*\le{m_-}/2$, the determinants are nonzero there
and their squared absolute values have positive expectations.
The logarithms of these expectations therefore
converge uniformly.  Pass to the limit in
\eqref{eq:qcdm2-log-circle} and apply \cref{C-415-16}.
The variance assertion now follows from \cref{C-414-11}, after
dividing $u$ by its supremum norm.
\end{proof}

\subsection{Loops with two or more propagators}

Some longer loops require less information about the energy.
A flavour cycle forces the first mixed source monomial to have
degree at least four.  The ordinary energy estimate is already
summable at this degree.

\begin{theorem}\label{C-415-3}
Suppose $2\le k\le N_f$, and let $f_1,\ldots,f_k$ be real bounded
functions on the torus.  Set
\[
 J_+=\sum_{i=1}^k f_iE_{i,i+1},\qquad E_{k,k+1}=E_{k,1},
 \qquad
 \lambda_{\mathsf q}=\operatorname{Tr}
       (\mathsf G_+\hat f_1\cdots\mathsf G_+\hat f_k).
\]
If $\rho\max_i\|f_i\|_\infty\le\rho_*$, there is a holomorphic
$\mathcal F$, with $\mathcal F(0)=0$, such that
\[
 \mathsf E_{\nu_K}|\det(1+z\mathsf G_+\hat J_+)|^2
 \le e^{(c_V(g)+C_4)\ell^3}|e^{\mathcal F(z)}|^2,
 \qquad |z|\le\rho.
\]
Consequently
\[
 \operatorname{Var}_{\nu_K}\lambda_{\mathsf q}
       \le\rho^{-2k}e^{(c_V(g)+C_4)\ell^3}.
\]
If $N_f\ge2$, $f,h$ are real and bounded, and
$\rho\|f\|_\infty\|h\|_\infty\le\rho_*^2$, the same conclusion
holds for
\[
 \Upsilon(\zeta)=\mathsf E_{\nu_K}
 |\det(1-\zeta\mathsf G_+\hat f\mathsf G_-\hat h)|^2,
 \quad |\zeta|\le\rho,
\]
with a holomorphic factor, and
\[
 \operatorname{Var}_{\nu_K}\operatorname{Tr}
       (\mathsf G_+\hat f\mathsf G_-\hat h)
       \le\rho^{-2}e^{(c_V(g)+C_4)\ell^3}.
\]
One can take
\[
 C_4=C_EC_{\mathfrak E}M^{-3}(2L)^{-4}
       \frac L{L-1}\bigl(1-(2L)^{-2}\bigr)^{-2}.
\]
\end{theorem}

\begin{proof}
Give the $k$ upper-species edges separate parameters $z_i$, and
the lower-species edges separate parameters $w_i$.  Diagonal
flavour rotations leave the determinants unchanged and multiply
the edge parameters by phases $u_i$ with $\prod_i u_i=1$.
Average the energy bound over this phase group in each species.
The surviving monomials are powers of
$\zeta=\prod_i z_i$ and $\xi=\prod_i w_i$.
The terms depending only on $\zeta$ or only on $\xi$ form the
holomorphic factor.  Every remaining monomial has a positive
power of each and hence at least $2k\ge4$ source factors.

The energy coefficients in \cite[Prop. 10.6]{QCDM1} are
holomorphic for $\lambda_j\|J\|<r_M$, with polymer bound
$C_{\mathfrak E}e^{-\kappa d_j(X)}$, with
$\sum_Xe^{-\kappa d_j(X)}\le C_En_jM^{-3}$.
Cauchy's estimate, at source size
$\rho_*\le r_M/(2L)$, bounds the mixed sum at level $j$ by
\[
 C_EC_{\mathfrak E}M^{-3}n_j
       (2L)^{-4}\lambda_j^4
       \bigl(1-(2L)^{-2}\bigr)^{-2}.
\]
Since $\sum_{j=1}^Kn_j\lambda_j^4
\le\ell^3L/(L-1)$, its sum is at most $C_4\ell^3$.
The geometric mean of the source bounds gives the asserted
determinant bound.  The flavour trace vanishes until the cycle
closes, so
\[
 \det(1+z\mathsf G_+\hat J_+)
       =1+(-1)^{k-1}z^k\lambda_{\mathsf q}+O(z^{2k}).
\]
Use the coefficient of $z^k$ in \eqref{eq:qcdm2-Parseval};
centering can only decrease its squared norm.

For the second assertion use the two-flavour source
\[
 (zf\tau_++z'h\tau_-)\otimes E_{11}
 +(wf\tau_++w'h\tau_-)\otimes E_{22}.
\]
The Schur determinant identity gives factors
$\det(1-zz'\mathsf G_+\hat f\mathsf G_-\hat h)$ and its
conjugate when $(w,w')=(\bar z,\bar z')$.  Phase averaging leaves
only powers of $zz'$ and $ww'$; a mixed term again has degree at
least four.  For $\zeta=zz'$ choose
$|z|\|f\|_\infty=|z'|\|h\|_\infty
=\sqrt{|\zeta|\|f\|_\infty\|h\|_\infty}$, which stays in the
source domain.  The same energy sum applies.  The coefficient of
$\zeta$ is the negative of the stated trace, so
\eqref{eq:qcdm2-Parseval} gives its variance bound.
\end{proof}

\subsection{Comparing two masses}\label{sec:qcdm2-mass-comparison}

A constant source changes the mass at every site. The circle
bound therefore compares the gauge measures at two masses.
We first control the derivative of the determinant and the
density for a short change of mass. Iterating this comparison
gives bounds throughout the permitted mass interval, uniformly
in the cutoff. If the same observables converge at both masses
along one cutoff sequence, their limiting laws are mutually
absolutely continuous; see \cref{C-406-11}.

Keep the chamber parameter $r=1/2$ and the constants of
\cref{sec:qcdm2-setting,sec:qcdm2-mixed-proof}. Fix
$0<g\le\gamma$, $\ell\in\mathfrak L(g)$ and
$\mathsf b\in I_g$. The physical torus $T_\ell$ remains fixed.
At cutoff $K$, write $a=L^{-K}$ and, for $s>0$, put
\begin{align}
 F_s(U)&=\det(D_+[U]+as),\qquad
 Z_K(s)=\int e^{-(\mathsf b/a)A(U)}|F_s(U)|^{2N_f}\,dU,
 \label{eq:qcdm2-mass-partition}\\
 \nu_{K,s}(dU)&=Z_K(s)^{-1}
       e^{-(\mathsf b/a)A(U)}|F_s(U)|^{2N_f}\,dU,
 \label{eq:qcdm2-mass-law}\\
 \mathfrak l_s(U)&=\partial_s\log|F_s(U)|^2
       =2a\operatorname{Re}\operatorname{Tr}(D_+[U]+as)^{-1}.
 \label{eq:qcdm2-mass-loop}
\end{align}
The determinant and trace include chambers, spin and colour,
but no flavour variable. We suppress $K$ in $F_s$ and
$\mathfrak l_s$, and write $\mathsf E_s$ for expectation under
$\nu_{K,s}$. The resolvent bound makes $F_s$ nonzero. At fixed
$K$ these densities are strictly positive and continuous on
the compact configuration space, and depend smoothly on
$s>0$.

With $\rho_*$ and $c_V(g)$ from \cref{C-415-17}, set
\begin{equation}\label{eq:qcdm2-mass-constants}
 \begin{split}
 M_g&=\exp\bigl((c_V(g)+C_{\rm mix})\ell^3\bigr),\\
 \mathfrak X_\star&=\frac{64M_g}{e^2N_f^2\rho_*^2},\qquad
 \mathfrak X_n^\sharp=4^n(1+M_g)^3\mathfrak X_\star
                       \quad(n\ge1).
 \end{split}
\end{equation}
These constants are independent of $K$, of
$\mathsf b\in I_g$ and of the masses in $[m_-,m_+]$.
Their dependence on the physical volume enters through $M_g$.

\begin{theorem}[Moments and changes of mass]\label{C-406-55}
For every cutoff $K\ge1$ the following assertions hold.
\begin{enumerate}[label=(\roman*)]
\item Let $t\in[m_-,m_+]$, $s\in\mathbb R$, $|s-t|\le r<\rho_*$ and
$p\ge1/2$. There is a real constant $c$, independent of $U$
and $p$, such that
\begin{equation}\label{eq:qcdm2-mass-moments}
 \mathsf E_t|\mathfrak l_s-c|^{2p}
 \le\left(\frac2{N_f}\right)^{2p}
      \left(\frac{4\rho_*}{(\rho_*-r)^2}\right)^{2p}
      \left(\frac pe\right)^{2p}M_g.
\end{equation}
In particular,
\begin{equation}\label{eq:qcdm2-score-variance}
 \operatorname{Var}_{\nu_{K,t}}\mathfrak l_t
       \le\mathfrak X_\star.
\end{equation}

\item If $t\in[m_-,m_+]$, $s\in\mathbb R$ and $|s-t|\le\rho_*/9$, then
the density $\rho_{s|t}=d\nu_{K,s}/d\nu_{K,t}$ satisfies
\begin{equation}\label{eq:qcdm2-nearby-density}
 \mathsf E_t\rho_{s|t}^2\le(1+M_g)^3,
 \qquad
 \mathsf E_t\rho_{s|t}^{-2}\le(1+M_g)^3.
\end{equation}

\item Fix $m_-\le\mathsf m<\mathsf m'\le m_+$ and put
\begin{equation}\label{eq:qcdm2-mass-chain-length}
 n_{\rm ch}=\left\lceil\frac{9(\mathsf m'-\mathsf m)}{\rho_*}
                 \right\rceil.
\end{equation}
For every $s,t\in[\mathsf m,\mathsf m']$,
\begin{align}
 \operatorname{Var}_{\nu_{K,t}}\mathfrak l_s
       &\le\mathfrak X_{n_{\rm ch}}^\sharp,
       \label{eq:qcdm2-cross-mass-variance}\\
 |\operatorname{Cov}_{\nu_{K,t}}(\mathfrak l_s,\mathfrak l_t)|
       &\le\mathfrak X_{n_{\rm ch}}^\sharp.
       \label{eq:qcdm2-cross-mass-covariance}
\end{align}
For every Borel set $E$ of configurations,
\begin{equation}\label{eq:qcdm2-mass-event-comparison}
 \nu_{K,t}(E)\le(1+M_g)^3
                       \nu_{K,s}(E)^{2^{-n_{\rm ch}}}.
\end{equation}
\end{enumerate}
\end{theorem}

\begin{proof}
Fix $K$ and $t\in[m_-,m_+]$. Apply \cref{C-415-17} on
$T_\ell$ with $f=1$, and denote its holomorphic function by
$h_t$. Since $\rho_*\le t/2$ and $\rho_*\le r_M/(2L)$,
\[
 \Omega_U(z)=N_f\left[
  \operatorname{Tr}\log\bigl(1+za(D_+[U]+at)^{-1}\bigr)-h_t(z)
                   \right]
\]
is holomorphic near $|z|\le\rho_*$ and vanishes at zero.
The matrix logarithm is the power series at the identity.
The exponential of its trace is $F_{t+z}/F_t$. The circle bound says
\begin{equation}\label{eq:qcdm2-mass-log-moment}
 \mathsf E_t e^{2\operatorname{Re}\Omega_U(z)}\le M_g,
                   \qquad |z|\le\rho_*.
\end{equation}
For interior points this follows from the maximum principle
applied to $\log\Xi_{\nu_{K,t}}^1-2N_f\operatorname{Re}h_t$,
which is subharmonic by \cref{C-415-16}.
Differentiation at real $z$ gives
\begin{equation}\label{eq:qcdm2-mass-loop-derivative}
 \mathfrak l_{t+z}
   =\frac2{N_f}\operatorname{Re}\Omega_U'(z)
                         +2\operatorname{Re}h_t'(z).
\end{equation}

Write
\[
 Y_U^+=\frac1{2\pi}\int_0^{2\pi}
       \bigl(\operatorname{Re}\Omega_U(\rho_*e^{i\theta})\bigr)^+
                                                    \,d\theta.
\]
The real part of $\Omega_U$ has mean zero on the circle, so
the mean of its absolute value is $2Y_U^+$. Schwarz's integral
formula, with $w=\rho_*e^{i\theta}$, therefore gives
\begin{align*}
 \Omega_U(z)&=\frac1{2\pi}\int_0^{2\pi}
       \operatorname{Re}\Omega_U(w)\frac{2z}{w-z}\,d\theta,\\
 \Omega_U'(z)&=\frac1{2\pi}\int_0^{2\pi}
       \operatorname{Re}\Omega_U(w)\frac{2w}{(w-z)^2}\,d\theta.
\end{align*}
Consequently, for $|z|\le r<\rho_*$,
\begin{equation}\label{eq:qcdm2-mass-Schwarz}
 |\Omega_U(z)|\le\frac{4r}{\rho_*-r}Y_U^+,
 \qquad
 |\Omega_U'(z)|\le\frac{4\rho_*}{(\rho_*-r)^2}Y_U^+.
\end{equation}
For $p\ge1/2$, Jensen's inequality and
$(y^+)^{2p}\le(p/e)^{2p}e^{2y}$ give
\[
 (Y_U^+)^{2p}\le\left(\frac pe\right)^{2p}
       \frac1{2\pi}\int_0^{2\pi}
             e^{2\operatorname{Re}\Omega_U(\rho_*e^{i\theta})}
                                                        \,d\theta.
\]
Similarly $e^{2y^+}\le1+e^{2y}$ and Jensen's inequality imply
\[
 e^{2Y_U^+}\le1+\frac1{2\pi}\int_0^{2\pi}
             e^{2\operatorname{Re}\Omega_U(\rho_*e^{i\theta})}
                                                        \,d\theta.
\]
Integrating these nonnegative bounds and using
\eqref{eq:qcdm2-mass-log-moment}, we obtain
\begin{equation}\label{eq:qcdm2-positive-boundary-moments}
 \mathsf E_t(Y_U^+)^{2p}\le(p/e)^{2p}M_g,
 \qquad \mathsf E_t e^{2Y_U^+}\le1+M_g.
\end{equation}
Take $c=2\operatorname{Re}h_t'(s-t)$ in
\eqref{eq:qcdm2-mass-loop-derivative}. Equations
\eqref{eq:qcdm2-mass-Schwarz} and
\eqref{eq:qcdm2-positive-boundary-moments} prove
\eqref{eq:qcdm2-mass-moments}. At $s=t$, take $r=0$ and
$p=1$ and use the minimizing property of the mean to get
\eqref{eq:qcdm2-score-variance}.

For (ii), put $z=s-t$. The factor $e^{2N_f\operatorname{Re}h_t(z)}$
is independent of $U$ and cancels from the normalized density:
\begin{equation}\label{eq:qcdm2-mass-density-identity}
 \rho_{s|t}(U)=\frac{|F_s(U)/F_t(U)|^{2N_f}}
                        {\mathsf E_t|F_s/F_t|^{2N_f}}
       =\frac{e^{2\operatorname{Re}\Omega_U(z)}}
                    {\mathsf E_t e^{2\operatorname{Re}\Omega_U(z)}}.
\end{equation}
If $|z|\le\rho_*/9$, then
$4|z|/(\rho_*-|z|)\le1/2$, and hence
$4|\Omega_U(z)|\le2Y_U^+$. In particular,
\[
 \mathsf E_t e^{4|\Omega_U(z)|}\le1+M_g,
 \qquad
 \mathsf E_t e^{2|\Omega_U(z)|}\le1+M_g.
\]
Cauchy--Schwarz, applied to $e^{|\Omega_U(z)|}$ and its
reciprocal, gives
\[
 (1+M_g)^{-1}
 \le\mathsf E_t e^{-2|\Omega_U(z)|}
 \le\mathsf E_t e^{2\operatorname{Re}\Omega_U(z)}
 \le M_g.
\]
Substitution in \eqref{eq:qcdm2-mass-density-identity} yields
\[
 \mathsf E_t\rho_{s|t}^2\le(1+M_g)^3,
 \qquad
 \mathsf E_t\rho_{s|t}^{-2}\le(1+M_g)M_g^2
                                      \le(1+M_g)^3.
\]

For (iii), first let $s\ne t$, set
$n'=\lceil9|t-s|/\rho_*\rceil\le n_{\rm ch}$ and divide
the segment from $s$ to $t$ into $n'$ equal parts, with
endpoints $u_i=s+i(t-s)/n'$. All these masses lie in the
admitted interval. For any nonnegative Borel function $Y$,
\eqref{eq:qcdm2-nearby-density} gives
\[
 \mathsf E_{u_{i+1}}Y
 =\mathsf E_{u_i}(Y\rho_{u_{i+1}|u_i})
 \le(1+M_g)^{3/2}(\mathsf E_{u_i}Y^2)^{1/2}.
\]
Iterating, the exponents $1/2,1/4,\ldots,2^{-n'}$ sum
to $1-2^{-n'}$, so
\begin{equation}\label{eq:qcdm2-mass-chain}
 \mathsf E_tY\le(1+M_g)^{3(1-2^{-n'})}
                        (\mathsf E_sY^{2^{n'}})^{2^{-n'}}.
\end{equation}
With $Y=1_E$ this proves \eqref{eq:qcdm2-mass-event-comparison}.
For the variance take $Y=(\mathfrak l_s-c_s)^2$, where
$c_s$ is the constant in (i) at the centre $s$, and use
(i) with $r=0$ and $p=2^{n'}$. Then
\[
 \begin{split}
 \operatorname{Var}_{\nu_{K,t}}\mathfrak l_s
 &\le\mathsf E_tY\\
 &\le(1+M_g)^3
       \frac{64\,4^{n'}}{e^2N_f^2\rho_*^2}M_g^{2^{-n'}}
 \le\mathfrak X_{n_{\rm ch}}^\sharp.
 \end{split}
\]
The last inequality uses $M_g\ge1$. Cauchy--Schwarz and
\eqref{eq:qcdm2-score-variance} give
\eqref{eq:qcdm2-cross-mass-covariance}. If $s=t$, the
variance estimate is already (i), and the event bound follows
from $\nu_{K,s}(E)\le\nu_{K,s}(E)^{2^{-n_{\rm ch}}}$.
\end{proof}

We now pass from the cutoff measures to any pair of limits
formed with the same observables. For probability measures
$P,Q$ on one measurable space, write
\[
 \operatorname{TV}(P,Q)=\sup_E|P(E)-Q(E)|,
 \qquad
 \operatorname{Aff}(P,Q)=\int\sqrt{pq}\,dR,
\]
where $R$ dominates $P,Q$ and $p=dP/dR$, $q=dQ/dR$.
The affinity is independent of $R$. We use
\[
 \operatorname{Ent}(P\mid Q)
   =\int\log\frac{dP}{dQ}\,dP
\]
when $P\ll Q$, and assign the value $+\infty$ otherwise.

\begin{proposition}[Limit laws at two masses]\label{C-406-11}
Retain the fixed $g,\ell,\mathsf b$ above, and choose
$m_-\le\mathsf m<\mathsf m'\le m_+$. Let $Y$ be a Polish
space and let $f_K$ be Borel maps from the cutoff
configuration spaces to $Y$. Suppose that, along one sequence
$K_j\to\infty$,
\[
 (f_{K_j})_*\nu_{K_j,\mathsf m}\ \Longrightarrow\ P,
 \qquad
 (f_{K_j})_*\nu_{K_j,\mathsf m'}\ \Longrightarrow\ Q.
\]
Put $\Delta=\mathsf m'-\mathsf m$, let $n_{\rm ch}$ be
\eqref{eq:qcdm2-mass-chain-length}, and set
\[
 x=\frac{N_f}{2}\sqrt{\mathfrak X_\star}\,\Delta
     =\frac{4\sqrt{M_g}}{e\rho_*}\,\Delta,
 \qquad
 H=\frac{N_f^2}{2}\mathfrak X_{n_{\rm ch}}^\sharp\Delta^2.
\]
Then $P$ and $Q$ are mutually absolutely continuous, and
for every Borel set $E\subset Y$,
\begin{equation}\label{eq:qcdm2-limit-mass-comparison}
 Q(E)\le(1+M_g)^3P(E)^{2^{-n_{\rm ch}}},
 \qquad
 P(E)\le(1+M_g)^3Q(E)^{2^{-n_{\rm ch}}}.
\end{equation}
Moreover,
\begin{align}
 \operatorname{TV}(P,Q)&\le x,
   \label{eq:qcdm2-limit-mass-TV}\\
 \operatorname{Ent}(P\mid Q),\ 
 \operatorname{Ent}(Q\mid P)&\le H,
   \label{eq:qcdm2-limit-mass-entropy}\\
 \operatorname{Aff}(P,Q)&\ge e^{-H/2},
   \label{eq:qcdm2-limit-mass-affinity}
\end{align}
and $\operatorname{Aff}(P,Q)\ge\cos x$ when $x\le\pi/2$.
\end{proposition}

\begin{proof}
We first prove the estimates before taking the limit. At
fixed $K$, let $p_t=d\nu_{K,t}/dU$. Differentiation of
\eqref{eq:qcdm2-mass-law} gives its logarithmic derivative,
\begin{equation}\label{eq:qcdm2-mass-score}
 \partial_t\log p_t
    =N_f(\mathfrak l_t-\mathsf E_t\mathfrak l_t),
 \qquad
 \partial_t\mathsf E_tF
    =N_f\operatorname{Cov}_{\nu_{K,t}}(F,\mathfrak l_t)
\end{equation}
for a bounded $F$ independent of $t$. Every differentiation
here is justified on a compact positive mass interval by the
resolvent bound and compactness of the configuration space.

The functions $p_t^{1/2}$ form a continuously differentiable
curve on the unit sphere of $L^2(dU)$, with speed
\[
 \|\partial_t p_t^{1/2}\|_2
    =\frac{N_f}{2}
           \sqrt{\operatorname{Var}_{\nu_{K,t}}\mathfrak l_t}.
\]
The angle between its endpoints is at most its length.
Their inner product is the affinity, so
\begin{equation}\label{eq:qcdm2-cutoff-mass-angle}
 \arccos\operatorname{Aff}(\nu_{K,\mathsf m},\nu_{K,\mathsf m'})
 \le\frac{N_f}{2}\int_{\mathsf m}^{\mathsf m'}
        \sqrt{\operatorname{Var}_{\nu_{K,t}}\mathfrak l_t}\,dt
 \le x.
\end{equation}
For any two probability densities, Cauchy--Schwarz gives
\[
 \operatorname{TV}(P,Q)
 =\frac12\int|\sqrt p-\sqrt q|(\sqrt p+\sqrt q)\,dR
 \le\sqrt{1-\operatorname{Aff}(P,Q)^2}
 \le\arccos\operatorname{Aff}(P,Q).
\]
Thus the total variation at the cutoff is at most $x$ as well.

For entropy, integration of the first identity in
\eqref{eq:qcdm2-mass-score} gives
\[
 \log p_{\mathsf m}-\log p_{\mathsf m'}
 =N_f\int_{\mathsf m}^{\mathsf m'}
                (\mathsf E_s\mathfrak l_s-\mathfrak l_s)\,ds.
\]
Taking its expectation at $\mathsf m$ and applying the
second identity in \eqref{eq:qcdm2-mass-score}, with $s$
held fixed, yields
\begin{align}
 \operatorname{Ent}(\nu_{K,\mathsf m}\mid\nu_{K,\mathsf m'})
 &=N_f^2\int_{\mathsf m}^{\mathsf m'}\!ds
        \int_{\mathsf m}^{s}\!dt\,
           \operatorname{Cov}_{\nu_{K,t}}(\mathfrak l_s,\mathfrak l_t),
       \label{eq:qcdm2-mass-entropy-triangle}\\
 \operatorname{Ent}(\nu_{K,\mathsf m'}\mid\nu_{K,\mathsf m})
 &=N_f^2\int_{\mathsf m}^{\mathsf m'}\!ds
        \int_s^{\mathsf m'}\!dt\,
           \operatorname{Cov}_{\nu_{K,t}}(\mathfrak l_s,\mathfrak l_t).
       \label{eq:qcdm2-mass-reverse-entropy-triangle}
\end{align}
Each triangle has area $\Delta^2/2$.
Equation \eqref{eq:qcdm2-cross-mass-covariance} therefore
bounds both entropies by $H$.

We recall the measure-theoretic facts needed to preserve
these estimates. Pushforward by a Borel map decreases total
variation and entropy and increases affinity. For total
variation the supremum is over fewer sets. For affinity,
take $R=(P+Q)/2$ and condition the densities $p,q$ on the
map: conditional Cauchy--Schwarz gives
\[
 \mathsf E_R(\sqrt{pq}\mid f)
    \le\sqrt{\mathsf E_R(p\mid f)\mathsf E_R(q\mid f)}.
\]
For entropy use the variational formula
\begin{equation}\label{eq:qcdm2-entropy-variational}
 \operatorname{Ent}(P\mid Q)
  =\sup_v\left(\int v\,dP-\log\int e^v\,dQ\right).
\end{equation}
On a Polish space the supremum may be taken either over
bounded Borel functions or over bounded continuous functions;
see \cite[equation (1.1) and the following paragraph]{DM22}.
Using $v\circ f$ in the Borel version proves the pushforward
inequality. The continuous version makes entropy lower
semicontinuous under simultaneous weak convergence.
The same is true for total variation, since
\[
 \operatorname{TV}(P,Q)
   =\frac12\sup_{\substack{v\in C_b(Y)\\\|v\|_\infty\le1}}
                   \left|\int v\,dP-\int v\,dQ\right|.
\]

Affinity is upper semicontinuous under that convergence.
Indeed,
\begin{equation}\label{eq:qcdm2-affinity-variational}
 \operatorname{Aff}(P,Q)
 =\inf_{\substack{v\in C_b(Y)\\\inf v>0}}
          \left(\int v\,dP\right)^{1/2}
          \left(\int v^{-1}\,dQ\right)^{1/2}.
\end{equation}
Cauchy--Schwarz proves one inequality. For the reverse,
truncate $\sqrt{q/p}$ to $[\epsilon,\epsilon^{-1}]$,
taking the value $\epsilon^{-1}$ where $p=0$. If the result
is $v_\epsilon$, then
\[
 \int v_\epsilon\,dP\le\operatorname{Aff}(P,Q)+2\epsilon,
 \qquad
 \int v_\epsilon^{-1}\,dQ\le\operatorname{Aff}(P,Q)+2\epsilon.
\]
Approximate $v_\epsilon$ in $L^1(P+Q)$ by bounded continuous
functions with values in the same interval, using Lusin's
theorem and extension from compact sets. Their reciprocals
converge in $L^1(P+Q)$ too. Letting $\epsilon$ decrease to
zero proves \eqref{eq:qcdm2-affinity-variational}. Its
right-hand side is an infimum of continuous functions of
$(P,Q)$, as required.

Apply these pushforward and semicontinuity statements to
\[
 P_j=(f_{K_j})_*\nu_{K_j,\mathsf m},\qquad
 Q_j=(f_{K_j})_*\nu_{K_j,\mathsf m'}.
\]
The cutoff estimates
give \eqref{eq:qcdm2-limit-mass-TV},
\eqref{eq:qcdm2-limit-mass-entropy}, and the cosine lower
bound. Finite entropy in both orders gives mutual absolute
continuity. Jensen's inequality then gives
\[
 -\log\operatorname{Aff}(P,Q)
 =-\log\int\left(\frac{dQ}{dP}\right)^{1/2}\,dP
 \le\frac12\operatorname{Ent}(P\mid Q)\le H/2,
\]
which proves \eqref{eq:qcdm2-limit-mass-affinity}.

It remains to retain the event comparison. By
\eqref{eq:qcdm2-mass-event-comparison}, with preimages under
$f_{K_j}$, every Borel $E\subset Y$ satisfies
\[
 Q_j(E)\le(1+M_g)^3P_j(E)^\alpha,
 \qquad \alpha=2^{-n_{\rm ch}}.
\]
For a nonempty closed set $F\subset Y$, choose a compatible
metric and put $G_\delta=\{y:\operatorname{dist}(y,F)<\delta\}$
and $F_\delta=\{y:\operatorname{dist}(y,F)\le\delta\}$.
The weak-convergence inequalities for open and closed sets give
\[
 \begin{split}
 Q(F)&\le Q(G_\delta)
      \le\liminf_j Q_j(G_\delta)\\
 &\le(1+M_g)^3\limsup_j P_j(F_\delta)^\alpha
      \le(1+M_g)^3P(F_\delta)^\alpha.
 \end{split}
\]
As $\delta\downarrow0$, $P(F_\delta)\downarrow P(F)$.
Inner regularity extends the inequality from closed sets
to all Borel sets. Interchanging the masses proves the
other inequality in \eqref{eq:qcdm2-limit-mass-comparison}.
\end{proof}


\section{The mixed source energy}
\label{sec:qcdm2-mixed-proof}

The ordinary source estimate gives one small factor for each
source.  Applied to the mixed energy, these two factors leave a
nonsummable power of the scale.  We need additional smallness in
the coefficient of the two sources.  We prove it in two stages.
First, at zero gauge source, the activities with fermion legs are
quadratic in $\phi_k$.  Second, at the trivial gauge configuration,
independent colour averages in the two species remove the mixed
terms of weighted degree two which have fermion legs.  The scalar
term left at this degree is bounded by $\mathfrak l_k^2$.

We recall the scale quantities of \cite[\S4]{QCDM1}:
\[
 \begin{gathered}
 \lambda_j=L^{j-K},\qquad \mathsf t_j=\lambda_j^{-1}\mathsf b,
 \qquad x_0=L^K\mathsf b,\quad
 x_{j+1}=L^{-1}x_j-b_{j+1},\quad |b_{j+1}|\le B_\sharp,\\
 \phi(t)=t^{-1/2}(1+\log t)^{p_0},\qquad
 \phi_j=\phi(x_j),\qquad
 \delta_j=A_1\mathsf t_j^{-1/2}(\log\mathsf t_j)^{p_0}.
 \end{gathered}
\]
The fixed exponent is $p_0\ge6r_c+1$.
The running and calibration coefficients satisfy
$|x_j-\mathsf t_j|\le\bar B$, and
$\phi_i/\phi_j\le(2L^{-1/2})^{j-i}$ for $i\le j$.
A level-$j$ polymer is a connected union of blocks of physical
side $M\lambda_j$; $d_j(X)$ is the least number of edges in
a block tree meeting every block of $X$. In the formulas
below $\delta_k$ is the fluctuation radius, $h_F$ is the
fixed Grassmann norm parameter, and $c_{\rm cl}$ is the
radius on which the cluster expansion converges.

We use the norms, local extraction and source coordinates of
\cite[Prop. 10.1 and its proof]{QCDM1}.  The level-$j$ norm
of a polymer family is the supremum, over a level-$j$ cube, of
the sum of the Grassmann norms of all its summands containing
that cube, weighted by $e^{\kappa d_j(X)}$.  Summands include
the local quadratic forms introduced by extraction.  A pointed
family is obtained by retaining a specified fermion leg or
specified scalar bond in the localization; its norm uses the
same weights.  The constant $K_\mu^{\rm s}$ bounds the increase
of this norm under local mass extraction.

A regular pair $(\mathbf U,\mathbf J)$ consists of the complex
connection and gauge-current coordinate in the analytic core of
\cite[\S4]{QCDM1}. This domain includes the bounds for the
local gauge potentials and the reduced block towers, as well
as those for curvature and current. We use radius $\alpha_0$
at every level. The gauge-current coordinate $\mathbf J$ is
distinct from the physical fermion source $J$ below; the
trivial pair is $(1,0)$.
There is no gauge source.  The doubled coordinate domain is
\begin{equation}\label{eq:qcdm2-zero-source-domain}
 2\mathscr D_k^\circ
 =\{\mathfrak c=(0,\mathsf n,(\sigma_i)_{i\le k}):
       L^{k+1}\|\mathsf n\|<2r_M,
       \|\sigma_i\|<2B_\sigma\phi_i\}.
\end{equation}
Its inner half is $\mathscr D_k^\circ$.
The bare coordinate is $\mathsf n=-aJ$.
The matrices $\sigma_i$ record the changes in the extracted
local mass at the preceding levels. At zero gauge source
the coordinate distance appearing below is
\[
 \|\mathfrak c-\mathfrak c'\|_{k,y}
 =\sup_x e^{-\kappa_\sigma|x-y|/L^{k+1}}
 \max\left\{\frac{L^{k+1}|\mathsf n-\mathsf n'|(x)}{r_M},
       \max_{i\le k}\frac{|\sigma_i-\sigma_i'|(x)}{B_\sigma\phi_i}\right\},
\]
where $x,y$ are in fine lattice units.
The matrices $\sigma_i$ remain free in the first-order domain;
the stronger estimate proved below is a conclusion, not an
assumption on these coordinates.

The role of these coordinates can be seen in one step of
\cite[Prop. 10.1]{QCDM1}. At level $k$, the remaining
fields have a quadratic fermion form and local interaction
activities. Completing the quadratic form after insertion
of the block-average kernel gives, on the variables being
integrated,
\[
 \Psi_k=H_k\Psi_{k+1}+W,\qquad
 \bar\Psi_k=\bar\Psi_{k+1}\bar H_k+\bar W.
\]
Here $H_k$ extends the prescribed coarse field to its fermion
background, $\bar H_k$ is its row twin, and $W,\bar W$ are
integrated with covariance $\Gamma_k^{-1}$. The row twin
need not be the adjoint. On a history, this notation
suppresses the stored variables, which remain unchanged.
Substituting these expressions into the old activities
expresses their fermion legs in terms of the new background
and the fluctuation. The gauge fluctuation is denoted by
$B$. The background pair $(1,0)$ fixes the background
connection and current; setting $B=0$ is a further
evaluation, not part of the definition of that pair.

The connected cluster expansion for the joined gauge and
Berezin integration produces activities
$\widetilde E^{(k+1)}(X;\mathfrak c)$ on level-$(k+1)$
polymers. Each is an even polynomial in the remaining
fermion fields, with coefficients depending on the
background pair and the coordinates $\mathfrak c$.
Its terms of degree zero contribute to the scalar energy.
More precisely, let $\mathfrak E_{k+1}(X;\mathfrak c)$
be their value at the trivial pair, including the explicit
normalizer and Gaussian terms. The source-dependent
constant removed from the activity is
\[
 e_{k+1}(X;\mathfrak c)
 =\mathfrak E_{k+1}(X;\mathfrak c)
                    -\mathfrak E_{k+1}(X;0).
\]
Its sum contributes to the source energy. Terms of degree
zero may still depend on the gauge background after this
constant has been removed.

The quadratic part has a matrix kernel. Its symmetric
zeroth moment determines the local mass increment
$\nu_{k+1}(\mathfrak c)$, with
$\mu_{k+1}=\nu_{k+1}(0)$ the zero-source increment.
The extraction rewrites this local quadratic contribution
as a mass term and retains the remaining kernel and its
localized correction forms in the activities. This is
why the activity norms below also include those forms.
For a physical source $J$, the mass change relative to
zero source at level $k$ is $-\lambda_kJ+\sigma_k(J)$.
The corrections are determined recursively by
\[
 \sigma_0(J)=0,\qquad
 \sigma_{k+1}(J)=L\sigma_k(J)
      +\nu_{k+1}(\mathfrak c_k(J))-\mu_{k+1},\qquad
 \mathfrak c_k(J)=(0,-aJ,(\sigma_i(J))_{i\le k}).
\]
The factor $L$ rescales the preceding mass correction.
The two extractions therefore give the outputs needed
here: a scalar source energy and a new local mass
coordinate, together with the activities carried to
the next step.

We first estimate the step with $\mathsf n$ and the
$\sigma_i$ independent in
\eqref{eq:qcdm2-zero-source-domain}. We then substitute
$\mathfrak c_k(J)$, as in \cite[Lem. 10.4]{QCDM1}.
For split sources, restrict an expression by zeroing
one species' entries while holding the other entries
fixed; for activities with fermion legs, zero that
species' generators as well. The mixed part is the
full expression minus its two one-species restrictions,
plus its restriction with both species removed.
Turning off one
physical source instead recomputes the entire chain:
the other species' corrections $\sigma_i$ can change
as well. The later comparison of source chains estimates
this additional change. Thus a bound for the mixed
coordinate expression is an input to the bound for
the mixed physical-source energy.

\subsection{Quadratic bounds at zero gauge source}

\begin{proposition}\label{C-407-42}
There are constants $C_2^{\rm s}$, $B_\nu^{(2)}$ and
$B_\sigma^{(2)}$, chosen independently of the admitted data,
such that the following holds after decreasing $\gamma$.
For $1\le j\le K$, the new activity with fermion legs, including
its extracted summands and every pointed family, has norm at
most $C_2^{\rm s}\phi_{j-1}^2$ throughout
$2\mathscr D_{j-1}^\circ$.  The increment satisfies
\begin{equation}\label{eq:qcdm2-quadratic-increment}
 |\nu_{k+1}(y;\mathfrak c)-\nu_{k+1}(y;\mathfrak c')|
 \le B_\nu^{(2)}\phi_k^2
        \|\mathfrak c-\mathfrak c'\|_{k,y},
\end{equation}
for $\mathfrak c\in\mathscr D_k^\circ$ and
$\mathfrak c'\in2\mathscr D_k^\circ$.
The norm on the right is the weighted coordinate norm of
\cite[Prop. 10.1]{QCDM1}.  Along the source chain,
\begin{equation}\label{eq:qcdm2-quadratic-chain}
 \|\sigma_j(0,J)\|\le
 B_\sigma^{(2)}\phi_{j-1}^2
          \frac{\lambda_j\|J\|}{r_M},
 \qquad
 B_\sigma^{(2)}=\frac{B_\nu^{(2)}}{1-4/L}.
\end{equation}
\end{proposition}

\begin{proof}
We give the additional estimates needed in the proof of
\cite[Prop. 10.1]{QCDM1}.  Write the potential as the sum
of the scalar term $V_0$, the gauge--fermion vertex $V_F$,
the old activity $V_{\rm old}$ with fermion legs, and the
quadratic source term $V_s$.  On the field disc,
\[
 \|V_F\|\le C_Vh_F^2\mathsf C_{\mathcal A}\delta_k,
 \quad \|V_0\|\le K_0^{\rm s}\delta_k,
 \quad \|V_s\|\le c'_\psi c_{\mathsf S}h_F^2
                  (\bar B_M+2B_\sigma)\phi_k,
 \quad \delta_k\le\sqrt2A_1\phi_k.
\]
Here one can take
\[
 K_0^{\rm s}=7K_P+c_{\mathsf s}K_{\rm old}^{\rm s}
       +\frac{13}{10}\mathsf C_{\mathcal A}K_\Lambda^{\rm s}
                         K_Mc_{\rm cl}/\alpha_0,
\]
using the scalar localization constants at radius $\alpha_0$.
The old-activity argument, including its
cell sums, gives
\[
 \|V_{\rm old}\|\le c'_\psi\Sigma_k,\qquad
 \Sigma_k=(1+c_{\rm cs})\sum_{i\le k}
             \omega_{k+1-i}\epsilon_i
             +C_{0,\sharp}^{\rm cs}h_F^2\phi_k^2,
\]
where $\epsilon_i$ includes the extraction factor
$K_\mu^{\rm s}$ and the zero-source bound $2C_2\phi_{i-1}^2$.
The additional term in
$C_{0,\sharp}^{\rm cs}=C_0^{\rm cs}
 +3L^{1/2}C_{\rm pc}S_{1/4}c_{\mathsf S}\bar B_M$
uses $|\mu_i|\le3\bar B_M\phi_i^2$ in the cell calculation.
Here $\mu_i$ is the zero-source mass increment.  Its local
mass forms cost $c_{\mathsf S}h_F^2|\mu_i|$ before they are
read at a later level.  Their sum is bounded by
\[
 \sum_{i\le k}\omega_{k+1-i}c_{\mathsf S}h_F^2|\mu_i|
 \le3c_{\mathsf S}\bar B_Mh_F^2\phi_k
                     \sum_{i\le k}\omega_{k+1-i}\phi_i
 \le3L^{1/2}C_{\rm pc}S_{1/4}
                 c_{\mathsf S}\bar B_Mh_F^2\phi_k^2.
\]
The change of a cell moment under the bare matrix mass has
an additional factor $L^i\|\mathsf n\|/r_M$.
Reading that cell at level $k$ costs $L^{k-i}$, and
\[
 L^{k-i}\frac{L^i\|\mathsf n\|}{r_M}<\frac2L
 \quad\hbox{on }2\mathscr D_k^\circ.
\]
Thus its sum is controlled by $h_F^2\sum_{i\le k}\phi_i^2$,
which is at most $h_F^2\phi_k^2/(1-2L^{-1/2})$.
This is the term included in $C_0^{\rm cs}$.
The other cell terms are linear in the old activity norms
$\epsilon_i$ and give the factor $1+c_{\rm cs}$.
These are the reasons the old-activity estimate has the
quadratic remainder $C_{0,\sharp}^{\rm cs}h_F^2\phi_k^2$.

Introduce the four-variable potential
\[
 V(z)=z_3V_0(B)+V_F(z_1B)
          +z_2V_{\rm old}(z_1B)+z_4V_s(z_1B).
\]
Take the radii
\[
 \begin{split}
 r_1&=\frac{\min\{c_{\rm cl}/(8C_Vh_F^2),\beta_0/2\}}
                 {\mathsf C_{\mathcal A}\delta_k},\\
 r_2&=\frac{c_{\rm cl}}{8c'_\psi\Sigma_k},\qquad
 r_3=\frac{c_{\rm cl}}{8K_0^{\rm s}\delta_k},\qquad
 r_4=\frac{c_{\rm cl}}
 {8c'_\psi c_{\mathsf S}h_F^2(\bar B_M+2B_\sigma)\phi_k}.
 \end{split}
\]
Each of the four terms has norm at most $c_{\rm cl}/8$
on this polydisc.  In particular, scaling $B$ by $z_1$
keeps the fine field in its analytic domain; the extra
bound $\beta_0/2$ in $r_1$ also keeps the changes of the
background legs within their Cauchy radius.  The cluster
expansion of \cite[Lem. 7.5]{QCDM1} therefore gives a
holomorphic family
\[
 \widetilde E_{\ge2}(X;z)=\sum_{n\in\mathbb N^4}a_n(X)z^n,
 \qquad
 \|a_n(X)\|\le K'_Mc_{\rm cl}
              \prod_{i=1}^4r_i^{-n_i}e^{-\kappa_+d_{k+1}(X)}.
\]
The coefficients take values in a finite Grassmann algebra,
so Cauchy's estimate applies in its norm.  After summing
polymers, $K'_M$ is replaced by $K_M$.

Assign weights $1,2,1,1$ to $z_1,z_2,z_3,z_4$.
We examine the coefficients before estimating their sum.
When $z_1=z_2=z_4=0$ the potential is scalar, so every
coefficient involving only $z_3$ has Grassmann degree zero.
When $z_2=z_3=z_4=0$, the substitution $B\mapsto-B$
replaces $z_1$ by $-z_1$.  The zero-source gauge Gaussian,
its tests and its local quadratic term are even in $B$;
the decoupled fermion covariance is independent of $B$.
The cluster construction respects this symmetry polymer
by polymer.  Its coefficient of $z_1$ is therefore zero.

For the coefficient linear in $z_4$, put $z_1=z_2=0$.
Let $\psi^{\rm bg}=H_{0,k+1}\Psi_{k+1}$ and
$\mathsf N=s_k+\sigma_k$, where $s_k$ is the zero-source
mass coordinate.  At $B=0$ the source vertex is
\[
 V_s(0)=-L^{-k}\left[
  \mathsf S_1\bigl(\mathsf N(\psi^{\rm bg}+H_{0,k}W)\bigr)
                -\mathsf S_1(\mathsf N\psi^{\rm bg})\right].
\]
Here $\mathsf S_1(\mathsf N\psi)$ denotes the quadratic
mass form with coefficient $\mathsf N$.  Expanding this
difference gives two terms linear in the fluctuation
generators and one term quadratic in those generators.
The linear terms have zero Berezin mean; the quadratic
term is contracted to a scalar.  At any power of $z_3$,
a coefficient linear in $z_4$ has exactly one such
vertex, independent of $B$, beside scalar vertices.
Its refined Gaussian integral thus has degree zero,
and the subsequent cluster operations preserve that
degree.  This proves that every coefficient of weighted
degree at most one has no part with fermion legs.

The coefficient of $z_2$ has weight two and must be
kept separately: it is the first variation of the old
family at $B=0$.  For an even potential $V$ with legs,
the cluster activity at $vV$ is holomorphic for
$|v|<c_{\rm cl}/\|V\|$ and has bound
$(17/16)K'_Mc_{\rm cl}e^{-\kappa_+d_{k+1}(X)}$.
Cauchy's estimate at $v=0$ bounds its first variation
by $(17/16)K'_M\|V\|e^{-\kappa_+d_{k+1}(X)}$.
At $V=V_{\rm old}(0)$ this yields
$(17/16)K_Mc'_\psi\Sigma_k$ in the polymer norm.

Choose $C_r^{\rm s}$ so that
$r_1^{-1},r_3^{-1},r_4^{-1}\le x=C_r^{\rm s}\phi_k$, and put
$y=r_2^{-1}=8c'_\psi\Sigma_k/c_{\rm cl}$.
For $n_2=0$, all remaining coefficients have
$n_1+n_3+n_4\ge2$.  Their Cauchy factors sum to
\[
 (1-x)^{-3}-1-3x\le22x^2,\qquad 0\le x\le\tfrac12.
\]
The quotient by $x^2$ has nonnegative coefficients and
takes the value 22 at $x=1/2$.  For $n_2\ge1$, remove
the first-variation coefficient already treated.
The remaining sum is at most
\[
 y\bigl((1-x)^{-3}(1-y)^{-1}-1\bigr)
       \le y(28x+2y),\qquad 0\le x,y\le\tfrac12,
\]
using $(1-x)^{-3}-1\le14x$ and $(1-y)^{-1}\le2$.
Together these give
\begin{equation}\label{eq:qcdm2-second-recurrence}
 \|\widetilde E_{\ge2}^{(k+1)}\|_{k+1}
 \le\tfrac{17}{16}K_Mc'_\psi\Sigma_k
       +K_Mc_{\rm cl}(22x^2+28xy+2y^2).
\end{equation}
For a pointed family, the insertion $(V_\psi)_z$
multiplies each monomial by a number in $[0,1]$ and
therefore preserves its degree and does not increase
its norm.  Replace the part with legs by
$tV_\psi+\zeta(V_\psi)_z$, for $|t|<9/8$ and
$|\zeta|<1$.  The pointed activity is
\[
 \int_0^1\left.\partial_\zeta
    \widetilde E_{\ge2}[z_3V_0,tV_\psi+\zeta(V_\psi)_z]
                                      \right|_{\zeta=0}\,dt.
\]
The scalar part costs $c_{\rm cl}/8$, and the part with
legs costs at most $(9/8+1)3c_{\rm cl}/8$; their sum is
$59c_{\rm cl}/64$.  Cauchy's estimate in $\zeta$ gives
the same Taylor coefficient bound as above.  The three
vanishing-coefficient arguments still hold because
pointing retains the monomials and commutes with scaling
$B$.  Thus \eqref{eq:qcdm2-second-recurrence} holds at
every pointer, with the same constants.

For completeness, the constants can be chosen in the following
order.  Put
\[
 R_a=8c'_\psi(1+c_{\rm cs})C_{\rm pc}S_{1/4}/c_{\rm cl},
 \qquad R_b=8c'_\psi C_{0,\sharp}^{\rm cs}h_F^2/c_{\rm cl},
\]
\[
 C_{(2)}^{\rm s}=\theta_E^{\rm s}C_{0,\sharp}^{\rm cs}h_F^2
  +K_\mu^{\rm s}K_Mc_{\rm cl}(22(C_r^{\rm s})^2+15R_b),
 \qquad C_2^{\rm s}=\max\{2C_{(2)}^{\rm s},2C_2\}.
\]
The source norm parameter $h_F$ is fixed with
$4\theta_E^{\rm s}C_{\rm pc}S_{1/4}(1+c_{\rm cs})\le1/2$,
as permitted in \cite[Prop. 10.20]{QCDM1}.  Decrease
$\gamma$ so that, with $\phi_* =\phi(\gamma^{-2}/2)$,
\[
 C_r^{\rm s}\phi_*\le\tfrac12,\qquad
 (R_aC_2^{\rm s}+R_b)\phi_*\le C_r^{\rm s},\qquad
 240K_\mu^{\rm s}K_Mc_{\rm cl}C_r^{\rm s}R_a\phi_*\le1.
\]
To close the induction, suppose
$\epsilon_i\le C_2^{\rm s}\phi_{i-1}^2$ for $i\le k$.
The old-family sum then satisfies
\[
 \sum_{i\le k}\omega_{k+1-i}\epsilon_i
 \le C_2^{\rm s}\phi_k
                  \sum_{i\le k}\omega_{k+1-i}\phi_{i-1}
 \le C_{\rm pc}S_{1/4}C_2^{\rm s}\phi_k^2.
\]
Consequently $y\le(R_aC_2^{\rm s}+R_b)\phi_k^2\le x\le1/2$.
After extraction the linear old term is bounded by
\[
 \tfrac{17}{16}K_\mu^{\rm s}K_Mc'_\psi\Sigma_k
 \le\tfrac18C_2^{\rm s}\phi_k^2
                 +\theta_E^{\rm s}C_{0,\sharp}^{\rm cs}h_F^2\phi_k^2.
\]
For the terms involving $y$, use $28x+2y\le30C_r^{\rm s}\phi_k$.
The third smallness condition and $C_r^{\rm s}\phi_k\le1/2$ give
\[
 K_\mu^{\rm s}K_Mc_{\rm cl}y(28x+2y)
 \le\tfrac18C_2^{\rm s}\phi_k^2
                      +15K_\mu^{\rm s}K_Mc_{\rm cl}R_b\phi_k^2.
\]
Adding the $22x^2$ term leaves
$(C_2^{\rm s}/4+C_{(2)}^{\rm s})\phi_k^2
\le3C_2^{\rm s}\phi_k^2/4$.
The other term in the definition of $\epsilon_{k+1}$ is
$2C_2\phi_k^2\le C_2^{\rm s}\phi_k^2$.
This closes the induction for the activity and every
pointed family at once.

The increment is the local zeroth moment of the quadratic
kernel.  We describe how its coordinate estimate uses the
improved norm.  A term $f$ in the moment at $y$ comes
from a polymer $X$ and a walk $\omega$.  It is holomorphic
on $2\mathscr D_k^\circ$ and bounded by $M_0$, the
product of its polymer bound and the row-sum bound of
the walk.  Restrict the bare-coordinate difference to
$X\cup\omega$ and each $\sigma_i$ difference to $X$;
call the resulting coordinate difference $\Delta$.
Let $T$ be the maximum of the normalized absolute
differences on these sets.  The value of $f$ depends
only on these entries, so
$f(\mathfrak c')=f(\mathfrak c+\Delta)$.
Since $\mathfrak c$ is in the inner domain, the line
$\mathfrak c+\zeta\Delta/T$, $|\zeta|<1$, lies in the
doubled domain.  Schwarz's lemma gives
\[
 |f(\mathfrak c')-f(\mathfrak c)|\le2M_0T;
\]
for $T>1$ this follows directly from the two bounds by
$M_0$.  Summing these inequalities is exactly the
kernel and walk sum in \cite[Prop. 10.1]{QCDM1}:
the reserved factor $e^{-\kappa d_{k+1}(X)/4}$ absorbs
the polynomial $1+d_{k+1}(X)$, and $e^{-|\omega|/4}$
absorbs the distance in the weighted coordinate norm.
Replacing the polymer size there by $C_2^{\rm s}\phi_k^2$
proves \eqref{eq:qcdm2-quadratic-increment} with
$B_\nu^{(2)}=2c_{\rm inc}h_F^{-2}C_2^{\rm s}$.

To evaluate this estimate on the source chain, we first
check its domain by induction using the first-order
increment estimate of \cite[Prop. 10.1 and Lem. 10.4]{QCDM1}.
The inductive bound is
\[
 \|\sigma_i(0,J)\|\le
 B_\nu S_\uparrow\phi_i\frac{\lambda_i\|J\|}{r_M},
 \qquad B_\nu S_\uparrow\le LB_\sigma,
\]
where $S_\uparrow$ bounds $\sum_{i\le j}\phi_i/\phi_j$.
It holds at level zero, where $\sigma_0=0$.
On $\lambda_j\|J\|<2r_M$, every earlier entry has
$\lambda_i\|J\|<2r_M/L$ and is therefore inside
the doubled coordinate domain.  Moreover, the normalized
coordinate distance at step $j$ is at most
$\lambda_j\|J\|/r_M$: this is its bare-coordinate
entry, and for $i<j$,
\[
 \frac{\|\sigma_i(0,J)\|}{B_\sigma\phi_i}
 \le L\frac{\lambda_i\|J\|}{r_M}
 \le\frac{\lambda_j\|J\|}{r_M}.
\]
The ordinary increment estimate is applicable with
one point $0\in\mathscr D_{j-1}^\circ$ and the other
$\mathfrak c_{j-1}(0,J)\in2\mathscr D_{j-1}^\circ$.
It gives
\[
 \|\nu_j(\mathfrak c_{j-1}(0,J))-\mu_j\|
       \le B_\nu\phi_j\frac{\lambda_j\|J\|}{r_M}.
\]
Summing this in the recurrence preserves the inductive bound because
$L^{j-i}\lambda_i=\lambda_j$ and
$\sum_{i\le j}\phi_i\le S_\uparrow\phi_j$.
The next increment is therefore defined, and
holomorphy follows by composition.  The stronger
estimate is applied to a chain whose existence and
domain have already been established.
Since $\nu_j(0)=\mu_j$, it gives
\[
 \|\nu_j(\mathfrak c_{j-1}(0,J))-\mu_j\|
       \le B_\nu^{(2)}\phi_{j-1}^2
                           \frac{\lambda_j\|J\|}{r_M}.
\]
Finally, iterate
$\sigma_j=L\sigma_{j-1}+\nu_j-\mu_j$ from $\sigma_0=0$.
For $1\le i\le j$,
\[
 L^{j-i}\lambda_i\phi_{i-1}^2
 =\lambda_j\phi_{i-1}^2
 \le\lambda_j\phi_{j-1}^2(4/L)^{j-i}.
\]
It follows that
\[
 \|\sigma_j(0,J)\|
 \le B_\nu^{(2)}\frac{\lambda_j\|J\|}{r_M}
                  \sum_{i=1}^j\phi_{i-1}^2
 \le\frac{B_\nu^{(2)}}{1-4/L}\phi_{j-1}^2
                           \frac{\lambda_j\|J\|}{r_M},
\]
which is \eqref{eq:qcdm2-quadratic-chain}.
\end{proof}

\subsection{Separation of the species}

Restrict to split coordinates in \eqref{eq:qcdm2-zero-source-domain}.
Let $R_b$ set both the coordinates and the Grassmann generators
of species $b$ to zero.  Define
\[
 R^{(0)}=R_1R_2,\quad R^{(1)}=(1-R_1)R_2,\quad
 R^{(2)}=R_1(1-R_2),\quad R^{\rm mix}=(1-R_1)(1-R_2).
\]
Let $P_{\rm col}$ average independent constant
$\mathrm{SU}(2)$ colour rotations in the two species at the
trivial pair.  These operations are contractions except for
the explicit factors $2$ in $R^{(b)}$ and $4$ in $R^{\rm mix}$.

\begin{lemma}\label{C-415-6}
At split coordinates the covariances, background maps and their
walk terms are block diagonal, with each block depending only
on its own coordinates.  The normalizers and Gaussian logarithms
are sums of two separate block contributions.  The source chain
stays split.  The new activities are neutral under the independent
phase rotations of the two species.  At the trivial pair
$P_{\rm col}$ commutes with reading at later levels and is the
identity on the quadratic part.
\end{lemma}

\begin{proof}
Each assertion about operators follows by taking the two blocks
in their defining inverse, average or walk formula.  A logarithmic
trace also separates into the two blocks.  The extraction map
respects independent phase rotations; hence its quadratic
kernel and its source increment are block diagonal.  This proves
the assertion about the chain by induction.  At the trivial pair
the background maps are scalar in colour.  A quadratic monomial
which is neutral in each species lies wholly in one block; its
colour-invariant kernel is a multiple of the identity in colour,
and independent colour averaging fixes it.
\end{proof}

\begin{lemma}\label{C-415-7}
Throughout the regular pair domain, the mixed scalar fluctuation
family and its pointed families satisfy
\[
 \|R^{\rm mix}F_0^{(k+1)}\|_{k+1}\le C_{\rm gm}\phi_k^2.
\]
The activities with fermion legs have mixed norm at most
$4C_2^{\rm s}\phi_k^2$.
\end{lemma}

\begin{proof}
The latter assertion follows from \cref{C-407-42}.  Write
$V=R^{(0)}V+R^{(1)}V+R^{(2)}V+R^{\rm mix}V$.
The preceding bounds give $\|R^{(b)}V\|\le V_1\phi_k$ and
\[
 \|R^{\rm mix}V\|
 \le4c'_\psi\bar\Sigma\phi_k^2+C_\Lambda C_{\rm gm}\phi_k^3,
 \qquad
 \bar\Sigma=(1+c_{\rm cs})C_{\rm pc}S_{1/4}C_2^{\rm s}
                 +C_{0,\sharp}^{\rm cs}h_F^2.
\]
Here $C_\Lambda=(13/12)K_\Lambda^{\rm s}
\mathsf C_{\mathcal A}\sqrt2A_1/\alpha_0$ is the scalar
localization constant in the proof of
\cite[Prop. 10.1]{QCDM1}, and
\[
 V_1=\sqrt2A_1(C_Vh_F^2\mathsf C_{\mathcal A}+2K_0^{\rm s})
       +c'_\psi c_{\mathsf S}h_F^2(\bar B_M+2B_\sigma)
       +2c'_\psi\bar\Sigma.
\]
Use parameters $t_1,t_2,u$ multiplying the last three parts,
of radii $c_{\rm cl}/(8V_1\phi_k)$ and
$c_{\rm cl}/(8\|R^{\rm mix}V\|)$, respectively.
At $u=0$ the mixed activity vanishes if either $t_b=0$:
the remaining Berezin covariance and all vertices then read
only the other species.  Schwarz's lemma in $t_1,t_2$ costs
$O(\phi_k^2)$; the difference between $u=1$ and $u=0$ costs
$O(\|R^{\rm mix}V\|)$.

Scalar subtraction compares the actual pair with the trivial
pair and costs a factor two.  The same argument applies to
pointed terms by the auxiliary differentiation used in the
previous proof, and to scalar bond pointers by the localization
bound \cite[Lem. 7.4]{QCDM1}.  In its notation one may take
\[
 C_{\rm gm}=17K_M^\bullet
       (64V_1^2/c_{\rm cl}+64c'_\psi\bar\Sigma).
\]
Decreasing $\gamma$ makes the last, cubic term at most half
the proposed quadratic bound and all parameter radii at least
two.  The induction starts with no old mixed family.
\end{proof}

\begin{lemma}\label{C-415-8}
At the trivial pair put $S_b=R^{(b)}V_0$.  Then
\begin{equation}\label{eq:qcdm2-scalar-quadratic}
 |S_b(Y;B)|\le v_{\mathsf S}\mathfrak l_k
       \left(M^{-3}\sum_{b'\subset Y}\mathsf t_k|B(b')|^2\right)
           e^{-\lambda_Vd_k(Y)},
 \qquad
 v_{\mathsf S}=\frac{2K_0^{\rm s}M^3}
                    {\mathsf C_{\mathcal A}^2A_1}.
\end{equation}
\end{lemma}

\begin{proof}
At $k=0$ there is no old scalar part, so $S_b=0$.
For $k\ge1$, scalar localization gives $S_b(Y;0)=0$ and the bound
$2K_0^{\rm s}\delta_ke^{-\lambda_Vd_k(Y)}$ on
$|B(b')|<\mathsf C_{\mathcal A}\delta_k$.  Its derivative at
zero is invariant under $B\mapsto\operatorname{Ad}(u)B$.
The Haar average of the nontrivial adjoint representation of
$\mathrm{SU}(2)$ is zero, so this derivative vanishes.
Apply Schwarz's lemma of order two to $S_b(Y;\zeta B)$,
using $\max_{b'\subset Y}|B(b')|^2
\le\sum_{b'\subset Y}|B(b')|^2$ and
$\delta_k\mathsf t_k=A_1\mathsf t_k^{1/2}
(\log\mathsf t_k)^{p_0}$.
\end{proof}

\subsection{A representative of the total mixed activity}

We must use cancellations in the sum of the activities at one
level.  They need not hold in each polymer of its chosen
localization.  We therefore propagate a representative of that
sum.  The representative is given only at the trivial pair;
no holomorphic extension in the pair is claimed or needed.
To express a level-$j$ family at level $k$, we substitute the
composed fermion background maps in its legs. At the next step
the substitution is $\Psi_k=H_k\Psi_{k+1}+W$, where $H_k$ is
the fermion background extension and $W$ is the Gaussian
Berezin fluctuation. Its fluctuation difference is obtained
by subtracting the same expression with $W=0$.
The sequence in \cite[Lem. 7.9]{QCDM1} is
$\omega_n=C_\omega L^{n/4}$.

The bound below separates the quadratic terms that contain a
difference of background fields from those that contain only their
values at a cell centre.  The former have the gain supplied by the
regularity of the background maps.  The latter retain the cell's
zeroth moment and must be estimated separately.  Equality of two
total activities does not provide this separation: the polymer norm
takes absolute values before summing over polymers.  We therefore
specify both the local kernels of a representative and the grouped
cell moments which will enter its bound.

\begin{lemma}\label{C-415-9}
Let $1\le j\le k<K$ and let $\mathsf n$ satisfy
$L^k\|\mathsf n\|<2r_M/L$.
Let $F$ be an even family at the trivial pair, with no
scalar part and with legs of levels at most $j$. Write
$F(X)=\sum_\sigma F_\sigma(X)$, assigning each summand to a
cell of $X$, and put
\[
 \epsilon_F=\sup_\square\sum_{X\supset\square}
 e^{\kappa d_j(X)}\sum_\sigma\|F_\sigma(X)\|_{h^{(j)}}.
\]
Express its quadratic part in the level-$j$ field $\Phi$ as
\[
 \sum_u\sum_{y,y'}\bar\Phi(y)\beta_u(y,y')\Phi(y').
\]
Each kernel $\beta_u$ is assigned to a cell $c(u)$, with
centre $z_{c(u)}$. Assume, with a fixed constant $C$, that
\begin{equation}\label{eq:qcdm2-cell-unit-bound}
 \sum_{u:c(u)=c}\sum_{y,y'}\|\beta_u(y,y')\|
 \bigl(1+|y-z_c|_j+|y'-z_c|_j\bigr)
 \le Ch_F^{-2}\epsilon_F
\end{equation}
for every cell $c$. Here $|\cdot|_j$ and
$\operatorname{dist}_j$ use level-$j$ lattice units.
The Grassmann weights $h^{(j)}$ and the background-map
constants $A_\kappa,C_H,\alpha_D$ are those of
\cite[Lem. 7.9]{QCDM1}.
The units are obtained by moving
diagonal terms $\bar\Phi(y)a_X(y)\Phi(y)$ from the cell
assigned to $X$ to the cell containing $y$. Assume
\begin{equation}\label{eq:qcdm2-moved-diagonal-bound}
 \|a_X(y)\|\le Ch_F^{-2}
 \bigl(\|F(X)\|+\|F_1(X)\|+\|G(X)\|\bigr)
 e^{-\alpha_D\operatorname{dist}_j(y,X)/2},
\end{equation}
where $F=F_1+F_2$ and $\|F_1\|_j,\|G\|_j\le\epsilon_F$.
For sums of summands, replace each norm in this inequality
by the sum of the corresponding norms. Each $a_X(y)$ is
also required to be a fixed linear combination, with sum
of absolute coefficients at most $2$, of row and column
sums at $y$ of the quadratic kernels of $F_1(X)$ and
$G(X)$. Their legs are expressed in $\Phi$ by the background
maps. The families $F_1$ and $G$ have the same even local
structure as $F$; $G$ may use a second mass matrix, with
the same background-map bounds. All these matrices are
in the admitted mass domain of
\cite[Lem. 7.9]{QCDM1}.

Let $m_c^{\rm tor}$ be the sum of the zeroth moments of
these units at $c$. On the torus the columns are transported
to $z_c$ along the lifted paths, including their
antiperiodic signs. For the resulting decomposition
$m_c^{\rm tor}=\sum_p m_{c,p}$ into terms supported on
polymers $Y_p$, use
\[
 \|m_c^{\rm tor}\|_{\rm nat}
 :=\sum_p\|m_{c,p}\|
 e^{\kappa L^{-(k-j)}d_j(Y_p)}(1+d_j(Y_p))^3.
\]
The summands here are the grouped coefficients obtained
after the cell cancellations.
Let $n_{\rm lag}$ be
the least nonnegative integer such that
$2(2A_\kappa C_H)^2L^{-n_{\rm lag}}\le1$.
When $F$ is expressed at level $k$ by the composed
background maps, at the trivial pair and with no shift,
its fluctuation difference satisfies
\begin{equation}\label{eq:qcdm2-old-representative}
 \|V_F\|\le c_\psi\left(
 \omega_{k+1-j}\epsilon_F+
 C_{\rm hl}L^{k-j}h_F^2
 \sup_c\|m_c^{\rm tor}\|_{\rm nat}\right)
\end{equation}
when $k-j\ge n_{\rm lag}$. For $k-j<n_{\rm lag}$ the first
term alone suffices. The constants depend on the fixed RG
parameters and the constant $C$ in the two hypotheses.
The total of $V_F$ is
$P(H_k\Psi_{k+1}+W)-P(H_k\Psi_{k+1})$, where $P$ is the
total of the read family.  The first variation $G_F$ of the
cluster activities at potential $vV_F$ has norm at most
$(17/16)K_M\|V_F\|$, and its total is the Berezin expectation
of the total of $V_F$.
\end{lemma}

\begin{proof}Put $n=k-j$.  We use the leg and walk estimates in the proof of
\cite[Lem. 7.9]{QCDM1}, at the trivial pair and with no shift.
The comparison there between the coefficient at a variable pair
and its value at the trivial pair is identically zero here.
The remaining estimates concern the Grassmann coefficients,
the background maps and the cell units in the hypotheses.

Suppose first that $n\ge n_{\rm lag}$.  Express each old leg in
the level-$j$ field $\Phi$, and then substitute the composed
background map from level $k$ into $\Phi$.  Each leg contributes
$L^{-n}$ relative to its old weight.  Its two expansions into
local and walk terms cost at most $(2A_\kappa C_H)^2$.
The definition of $n_{\rm lag}$ makes their product at most
$1/2$.  For Grassmann degree $2p\ge4$, the factors beyond
the fourth leg are therefore bounded by $2^{-(2p-4)}$;
this also pays the polynomials in the number of legs in the
walk count.  The four remaining legs give $L^{-4n}$,
while a level-$k$ cube contains $L^{3n}$ level-$j$ cells.
Thus the part of degree at least four is bounded by
$C\epsilon_FL^{-n}$ per level-$k$ cube, with the same
localization estimates as in that lemma.

For the quadratic part, fix a unit $u$ assigned to $c$ and
put $z=z_c$.  Carry its fields and kernel to $z$ along the
prescribed lifted paths.  In the next display the transported
quantities are written without additional notation; the
transports include the antiperiodic signs.  Adding and
subtracting the reference fields gives the exact identity
\begin{align*}
 \sum_{y,y'}\bar\Phi(y)\beta_u(y,y')\Phi(y')
 ={}&\bar\Phi(z)m_u\Phi(z)\\
 &+\sum_{y,y'}[\bar\Phi(y)-\bar\Phi(z)]
                       \beta_u(y,y')\Phi(y')\\
 &+\sum_{y,y'}\bar\Phi(z)\beta_u(y,y')
                              [\Phi(y')-\Phi(z)],
 \qquad m_u=\sum_{y,y'}\beta_u(y,y').
\end{align*}
After summing the units assigned to $c$, the first term is
$\bar\Phi(z_c)m_c^{\rm tor}\Phi(z_c)$.  The other two
terms each contain a difference.  This is the point at
which the grouped zeroth moment enters: the norm of this
zeroth-moment term is estimated using $m_c^{\rm tor}$ after
the cancellations in the cell have been made.

The background estimate with H\"older exponent $3/4$
bounds a transported difference along a path of length
$s$ in level-$j$ units by $Ch_FsL^{-7n/4}$.
The other leg costs $Ch_FL^{-n}$.  For a moved unit,
the path lies in its assigned cell and its length is
paid by \eqref{eq:qcdm2-cell-unit-bound}.  For a polymer
unit, the paths run through its polymer and the walk
polymers of its two legs.  Their lengths are bounded by
a fixed multiple of $M(1+d_j(X))$ plus the walk lengths,
where $X$ is the original polymer. The unused polymer
decay and the walk weights
pay these factors, as in the proof of
\cite[Lem. 7.9]{QCDM1}.  We obtain
\[
 \text{sum of norms of the difference terms at one cell}
       \le C\epsilon_FL^{-11n/4}.
\]
Summing over the $L^{3n}$ old cells gives
$C\epsilon_FL^{n/4}$.  The estimate uses the actual
lifted paths; no wrapping path is replaced by a shorter
distance on the torus.

We also keep track of the diagonal terms that were moved
between cells.  A term
$\bar\Phi(y)a_X(y)\Phi(y)$ is subtracted from the unit
assigned to $X$ and added to the unit assigned to the
cell of $y$.  These two copies cancel in the total.
For the local estimate, split $a_X(y)=\sum_pa_{X,p}(y)$
according to the local or walk term in each of the two
legs of the kernels from which it was formed.  Let
$Z_p$ be their supporting polymer.  The row-and-column
hypothesis ensures that this is the same decomposition
used for the kernels of $F_1$ and $G$, with uniformly
bounded coefficients, and that $a_{X,p}(y)=0$ unless
$y\in Z_p$.  The localized row-sum estimate gives
\[
 \sum_{p,y}\|a_{X,p}(y)\|
             e^{\kappa(d_j(Z_p)-d_j(X))}
 \le CA_\kappa^2h_F^{-2}
          \bigl(\|F_1(X)\|+\|G(X)\|\bigr).
\]
For families with several summands, the right-hand side
is the sum of their norms.  The same estimate holds
for $G$ at its second mass matrix, since its background
maps have the same constants.  This is the estimate for
moved pieces in \cite[Prop. 10.1]{QCDM1}.
It controls their actual supports as well as their
coefficients.

Changing the cell assignment changes a level-$k$ cube
sum only through terms whose two assignments lie on
opposite sides of its boundary.  The count in
\cite[Lem. 7.9]{QCDM1} bounds these terms by
$C\epsilon_F$: the boundary contains $O(L^{2n})$
old cells, the two legs give $L^{-2n}$, and the
exponential in \eqref{eq:qcdm2-moved-diagonal-bound}
pays the distance of a moved site from its polymer.
The polymer and walk weights pay the remaining
diameter factors.  This contribution is included in
$\omega_{k+1-j}\epsilon_F$.

To localize the difference terms, take the level-$k$
hull of each supporting polymer and enlarge it by the
walk polymers of its two final legs.  The hull estimate
\[
 d_k(Z^{\rm up})\le L^{-n}d_j(Z)
\]
and the weighted walk sums give the local polymer norm.
The spare rate $\kappa(1-L^{-n})d_j(Z)$ pays the
polynomials left by the path and cell counts.  Together
with the preceding degree and boundary estimates,
these are the terms bounded by
$\omega_{k+1-j}\epsilon_F$.

It remains to localize the zeroth-moment terms.  Use the
decomposition $m_c^{\rm tor}=\sum_pm_{c,p}$ from the
statement, after its cell cancellations.  Assign
\[
 \bar\Phi(z_c)m_{c,p}\Phi(z_c)
\]
to the level-$k$ hull of $Y_p$, enlarged by the walk
polymers of the two legs.  Those legs give
$Ch_F^2L^{-2n}$, and the count of old cells gives
$L^{3n}$.  The hull estimate accounts for the exponential
$e^{\kappa L^{-n}d_j(Y_p)}$ in the norm of the moment;
the factor $(1+d_j(Y_p))^3$ pays the number of cells
from which its hull can reach a fixed output cube.
The hull and walk count of \cite[Prop. 10.1]{QCDM1}
therefore gives
\[
 C_{\rm hl}L^nh_F^2
                  \sup_c\|m_c^{\rm tor}\|_{\rm nat}.
\]
This is the second term of
\eqref{eq:qcdm2-old-representative}.  It retains the
torus moment itself, including any contribution of
wrapping pieces.

For $n<n_{\rm lag}$, use the original summands without
expanding their legs about cell centres.  The count is
at most $L^{3n_{\rm lag}}$, a fixed factor already
included in $C_\omega$.  This proves the asserted bound
without a separate moment term at these lags.

Finally substitute $\Psi_k=H_k\Psi_{k+1}+W$ and
subtract the value at $W=0$.  Every surviving monomial
contains a fluctuation leg.  The factor $c_\psi$ follows
from the fluctuation weights, rather than from counting
these monomials individually.  Let $h$ be the original
leg weight, $h_a$ a bound for its background image,
and $h_b$ the weight of its fluctuation image.
Choose a positive allowance $h_b'$ with
$h_a+h_b'\le h$ and $h_b\le h_b'$. For a homogeneous
product of degree $q$,
\[
 (h_a+h_b)^q-h_a^q
 \le\frac{h_b}{h_b'}
            \bigl((h_a+h_b')^q-h_a^q\bigr).
\]
The background image uses at most half the leg weight;
take the other half for $h_b'$. The localized $W$ bound
then gives $h_b/h_b'\le c_\psi$, as in
the old-potential estimate of \cite[\S7]{QCDM1}.
Applying this inequality to the localized terms above
proves \eqref{eq:qcdm2-old-representative}, including
its short-lag assertion.

All decompositions used here preserve the sum: the
moved diagonal terms were added and subtracted, the
cell expansion is an identity, and the walk series
converge absolutely.  Thus the total of $V_F$ is
$P(H_k\Psi_{k+1}+W)-P(H_k\Psi_{k+1})$.
For the first variation, the cluster activity at
potential $vV_F$ is holomorphic on
$|v|<c_{\rm cl}/\|V_F\|$ and has bound
\[
 \frac{17}{16}K_M'c_{\rm cl}
                         e^{-\kappa_+d_{k+1}(X)}.
\]
Cauchy's estimate, followed by the polymer sum, yields
$\|G_F\|_{k+1}\le(17/16)K_M\|V_F\|$.
Since $V_F$ is independent of $B$, the joined integral
factors into its gauge integral and
$\mu_{\Gamma_k^{-1}}(e^{v\sum_XV_F(X)})$.
Differentiating its logarithm at $v=0$, and using
$\mu_{\Gamma_k^{-1}}(1)=1$, identifies the total of
$G_F$ with $\mu_{\Gamma_k^{-1}}(\sum_XV_F(X))$.
\end{proof}

Put
\[
 q_k=\phi_k^3+\phi_k^2e^{-c_{\rm t}\ell/(2\lambda_{k+1})},
 \qquad \bar q_k=q_k+\mathfrak l_k^2.
\]
The following estimate contains the cancellation needed for
\cref{C-415-5}.

\begin{lemma}\label{C-415-10}
There are constants $E_3,N_3,C_{\mathfrak d}$, independent
of the admitted data, with these properties.
At level $j$, the total of
$P_{\rm col}R^{\rm mix}E_{\ge2}^{(j)}$ at the trivial pair
has an even representative $F^j$, neutral in each species
and invariant under $P_{\rm col}$, with no scalar part and
\begin{equation}\label{eq:qcdm2-representative-size}
 \epsilon_j^{(3)}\le E_3q_{j-1}.
\end{equation}
Its cell sums satisfy
\begin{align}
 \|m_c^{\rm tor}\|_{\rm nat}
 &\le C_{\rm pw}h_F^{-2}\epsilon_j^{(3)}
 \left[\eta_{\min}^{-2}
   \left(C_{\rm bg}^{\rm w}\frac{L^j\|\mathsf n\|}{r_M}
         +C_Ae^{-4c_{\rm t}\ell/\lambda_j}\right)
            +e^{-c_{\rm t}\ell/\lambda_j}\right]
       +C_{\mathfrak d}h_F^{-2}\phi_{j-1}^2
                       e^{-c_{\rm t}\ell/\lambda_j}.
 \label{eq:qcdm2-representative-cells}
\end{align}
If $f_j(\mathfrak c)$ is the total scalar output of step $j$
at the trivial pair, including its normalizers and Gaussian
terms, then
\begin{equation}\label{eq:qcdm2-mixed-total}
 |R^{\rm mix}f_j(\mathfrak c)|
       \le E_3M^{-3}n_j\bar q_{j-1}.
\end{equation}
For split $\mathfrak c\in2\mathscr D_{j-1}^\circ$, write
$\nu_{j,1}=e_+\nu_je_+$ and $\nu_{j,2}=e_-\nu_je_-$.
If $b'\ne b$, then
\begin{equation}\label{eq:qcdm2-cross-increment}
 \sup_y\|\nu_{j,b}(\mathfrak c)(y)
             -\nu_{j,b}(R_{b'}\mathfrak c)(y)\|
 \le N_3q_{j-1}.
\end{equation}
Thus the coordinates of species $b$ are held fixed in this
comparison. No source factor is included in this bound.
\end{lemma}

\begin{proof}
The induction starts with no old mixed activity.  Fix a split
coordinate $\mathfrak c\in2\mathscr D_k^\circ$ and suppose the
three assertions hold through level $k$.  All quantities in this
proof are evaluated at the trivial pair.  We first identify the
terms whose totals must be cancelled, then construct a local
family representing the remaining total.

\emph{The deformed potential.}
Let $V_g=R^{(0)}V_0$ be the zero-source scalar potential.
For species $b$, let $V_{F,b}$ and $V_{s,b}$ be the
gauge--fermion vertex and mass vertex restricted to its
generators; let $V_{{\rm old},b}=R^{(b)}V_{{\rm old},\ge2}$;
and let $S_b=R^{(b)}V_0$ as in \cref{C-415-8}.
The old mixed family is split into its value at $B=0$
and its change from that value.  With ten parameters
$p=(z_1,z_2,x_1,x_2,y_1,y_2,d_1,d_2,u,u')$, put
\begin{align*}
 V_g+\sum_{b=1}^2\bigl(
  V_{F,b}(z_bB)+x_bV_{s,b}(z_bB)
       +y_bV_{{\rm old},b}(z_bB)+d_bS_b(B)\bigr)
       +uV_{{\rm mix},0}+u'V_{{\rm mix},1}(B).
\end{align*}
Here the last two terms mean
\[
 V_{{\rm mix},0}=R^{\rm mix}V_{{\rm old},\ge2}(0),\qquad
 V_{{\rm mix},1}(B)
 =R^{\rm mix}\bigl(V_{{\rm old},\ge2}(B)
                       -V_{{\rm old},\ge2}(0)\bigr)
                         +R^{\rm mix}V_0(B).
\]
In a term evaluated at $z_bB$, all its background maps
are evaluated at the fine connection reconstructed from
$z_bB$.  At $p=(1,\ldots,1)$ this is the original potential.
The parameter $u$ isolates the part which can be propagated
by the old representative; $u'$ groups terms which already
have three small factors.  Indeed, the bound
$\|R^{\rm mix}V_{{\rm old},\ge2}\|
\le4c'_\psi\bar\Sigma\phi_k^2$, followed by Schwarz's lemma
in the scaling of $B$, gives
\[
 \|R^{\rm mix}(V_{{\rm old},\ge2}(B)
                       -V_{{\rm old},\ge2}(0))\|
       \le8c'_\psi\bar\Sigma\phi_k^2r_z^{-1}.
\]
The scalar localization estimate in \cref{C-415-7}
supplies the other term.  Thus
\[
 \|V_{{\rm mix},1}\|
 \le \left(8c'_\psi\bar\Sigma C'_r
             +C_\Lambda C_{\rm gm}\right)\phi_k^3.
\]
Here
$C'_r=\sqrt2A_1\mathsf C_{\mathcal A}
\max\{16C_Vh_F^2/c_{\rm cl},2/\beta_0\}$.
The first term has acquired an extra factor $\phi_k$
from $r_z^{-1}\le C'_r\phi_k$; the second has that factor
from scalar localization.

Choose parameter radii
\begin{align*}
 r_z&=\frac{\min\{c_{\rm cl}/(16C_Vh_F^2),\beta_0/2\}}
                  {\mathsf C_{\mathcal A}\delta_k},&
 r_x&=\frac{c_{\rm cl}}
 {16c'_\psi c_{\mathsf S}h_F^2(\bar B_M+2B_\sigma)\phi_k},\\
 r_y&=\frac{c_{\rm cl}}{32c'_\psi\bar\Sigma\phi_k^2},&
 r_d&=\frac{c_{\rm cl}}{16v_{\mathsf S}\mathfrak l_k},\\
 r_u&=\frac{c_{\rm cl}}{64c'_\psi\bar\Sigma\phi_k^2},&
 r_{u'}&=\frac{c_{\rm cl}}
 {16(8c'_\psi\bar\Sigma C'_r+C_\Lambda C_{\rm gm})\phi_k^3}.
\end{align*}
The eleven displayed parts of the potential have total norm
at most $11c_{\rm cl}/16$; the cluster tail costs at most
$c_{\rm cl}/16$.  Assign weight one to $z_b,x_b,d_b$, weight
two to $y_b,u$, and weight three to $u'$.
Since $\mathfrak l_k\le\sqrt2\phi_k$, there is a fixed
$C_r^{\rm b}\ge2$ such that the inverse radii are bounded by
$(C_r^{\rm b}\phi_k)$ to their respective weights.  Decrease
$\gamma$ until $C_r^{\rm b}\phi_*\le1/2$.
For a multi-index $v\in\mathbb N^{10}$ denote its weighted
degree by $|v|_w$.  Cauchy's estimate gives
\begin{equation}\label{eq:qcdm2-ten-parameter-coefficients}
 \widetilde E(X;p)=\sum_v a_v(X)p^v,\qquad
 \|a_v(X)\|\le K'_Mc_{\rm cl}r^{-v}
                        e^{-\kappa_+d_{k+1}(X)},\qquad
 r^{-v}\le(C_r^{\rm b}\phi_k)^{|v|_w}.
\end{equation}
Every radius is at least two, so these series converge
at $p=(1,\ldots,1)$.

To identify a coefficient of the total, let $\mu_B$
denote the normalized zero-source gauge integral of this
step, including its scalar potential $V_g$, and let
$\mu_1,\mu_2$ be the decoupled Berezin means of the two
species.  The remaining background generators are not
integrated.  The total of the cluster activities is
the logarithm of the joined integral.  Hence the
coefficient of two distinct linear parameters is
their connected expectation under
\[
 \langle A\rangle_*:=\mu_B\mu_1\mu_2(A),\qquad
 \kappa_*(A,B)=\langle AB\rangle_*
                         -\langle A\rangle_*\langle B\rangle_*.
\]
This identity holds in the even algebra of the retained
background generators.  It describes the sum over $X$
of a coefficient in \eqref{eq:qcdm2-ten-parameter-coefficients},
not its separate polymer terms.

\emph{The coefficients of weight two.}
Call an index mixed if it contains a parameter of each
species, or contains $u$ or $u'$.  Every other coefficient
is killed by $R^{\rm mix}$ polymer by polymer.  For example,
with $u=u'=0$ and all parameters of species two zero,
the potential has no generator of that species and
does not depend on its coordinates.  Its Berezin
integrals use only $\mu_1$, the unused normalized
mean contributing one.  Thus its activities also
have no dependence on species two.

There are precisely ten mixed indices of weight at most
two: the index $u$, and the nine products formed by one
of $z_1,x_1,d_1$ and one of $z_2,x_2,d_2$.
The vertices in these products are
\[
 l_b(B)=\left.\partial_{z_b}V_{F,b}(z_bB)\right|_{z_b=0},
 \qquad M_b=V_{s,b}(0),\qquad S_b(B).
\]
The first is linear in $B$ and quadratic in the fields
of species $b$; $M_b$ is independent of $B$ and has only
that species' generators; $S_b$ is scalar.

The five products containing a mass vertex have zero
total.  For example, the independence of $\mu_1$ and
$\mu_2$ gives
\[
 \langle M_1M_2\rangle_*
       =\mu_1(M_1)\mu_2(M_2)
       =\langle M_1\rangle_*\langle M_2\rangle_*.
\]
Since $M_2$ is independent of $B$, the same factorization
holds for $\langle l_1(B)M_2\rangle_*$ and
$\langle S_1(B)M_2\rangle_*$.  Exchanging the species
proves the cancellations for
$x_1x_2,z_1x_2,x_1z_2,x_1d_2,d_1x_2$.

For the other three products containing a gauge vertex,
integrate the fermion fluctuations first.  Put
$L_b(B)=\mu_b(l_b(B))$.  The corresponding totals are
\[
 \kappa_B(L_1,L_2),\qquad
 \kappa_B(L_1,S_2),\qquad
 \kappa_B(S_1,L_2),
\]
where $\kappa_B(A,B)=\mu_B(AB)-\mu_B(A)\mu_B(B)$.
These need not vanish before colour averaging.
Let $\alpha_g^{(b)}$ rotate the background generators
of species $b$ by the constant colour matrix $g$.
At the trivial pair the vertex and covariance satisfy
\[
 L_b(\operatorname{Ad}(g)B)\circ\alpha_g^{(b)}=L_b(B).
\]
The covariance is scalar in colour, which permits
commuting the rotation through $\mu_b$.
Since $L_b$ is linear in $B$, its average in the
background colour variables is
\[
 P^{(b)}L_b(B)
 =L_b\left(\int_{\mathrm{SU}(2)}
                         \operatorname{Ad}(g^{-1})B\,dg\right)=0.
\]
The adjoint representation has no invariant vector.
The averages of the two species are independent, so
averaging the first species kills $\kappa_B(L_1,L_2)$
and $\kappa_B(L_1,S_2)$, and averaging the second kills
$\kappa_B(S_1,L_2)$.  In Grassmann degree zero the
colour average is the identity: the scalar part of
$L_b$ is itself an invariant linear form in $B$ and
vanishes.  Thus these cancellations hold in every
degree after $P_{\rm col}R^{\rm mix}$.

The ninth product, $d_1d_2$, is different.  Its total
is the gauge covariance $\kappa_B(S_1,S_2)$, so it
has degree zero.  Both vertices vanish when their
own species coordinate is zero; therefore $R^{\rm mix}$
fixes this coefficient.  There is no cancellation
asserted for it.  Its Cauchy estimate is
\begin{equation}\label{eq:qcdm2-two-scalar-coefficient}
 \sum_X|a_{d_1d_2}(X)|
 \le K_Mc_{\rm cl}r_d^{-2}n_{k+1}M^{-3}
 =\frac{256K_Mv_{\mathsf S}^2}{c_{\rm cl}}
                  \mathfrak l_k^2n_{k+1}M^{-3}.
\end{equation}

\emph{Propagation of the old mixed total.}
The remaining coefficient of weight two is $u$.
Because $V_{{\rm mix},0}$ is independent of $B$, its
total is $\mu_1\mu_2$ of the total of that vertex.
If $P_i$ is the total of the level-$i$ family after
it has been read at level $k$, the fluctuation vertex
has total
\[
 R^{\rm mix}\sum_{i\le k}
       \bigl(P_i(H_k\Psi_{k+1}+W)-P_i(H_k\Psi_{k+1})\bigr).
\]
The projection $P_{\rm col}$ commutes with these
background maps, with $R^{\rm mix}$ and with $\mu_1\mu_2$.
Thus the induction hypothesis permits replacing each
total by that of $F^i$.  Apply \cref{C-415-9} to each
$F^i$ and add the resulting families.  If $V_F$ denotes
this sum, the coefficient of $u$, after colour averaging,
has the total of its first-variation family $G_F$.
This replacement uses equality of totals, not equality
of the chosen localizations.

In the cell bound, terms containing
$L^i\|\mathsf n\|/r_M$ acquire the factor $L^{k-i}$
from reading at level $k$; their product is less than
$2/L$.  The exponential terms have uniformly bounded
products with $L^{k-i}$.  These factors are included
in $c_{\rm cs}$, exactly as in the ordinary old-family
count.  The last term of
\eqref{eq:qcdm2-representative-cells} must instead be
summed separately.  We obtain
\begin{align*}
 \|V_F\|\le c_\psi\biggl[
 &(1+c_{\rm cs})\sum_{i\le k}\omega_{k+1-i}\epsilon_i^{(3)}\\
 &+C_{\rm hl}C_{\mathfrak d}\sum_{i\le k}
       L^{k-i}\phi_{i-1}^2e^{-c_{\rm t}\ell/\lambda_i}\biggr].
\end{align*}
Both sums are bounded by a multiple of $q_k$.
For the first, monotonicity of $\phi_i$ and of the
torus exponential gives
\[
 q_{i-1}\le\frac{\phi_{i-1}}{\phi_k}q_k,
 \qquad
 \sum_{i\le k}\omega_{k+1-i}\phi_{i-1}
                         \le C_{\rm pc}S_{1/4}\phi_k.
\]
For the second, write $v=k-i$ and divide by the
torus term $\phi_k^2e^{-c_{\rm t}\ell/(2\lambda_{k+1})}$
in $q_k$.  The quotient is at most
\[
 \sum_{v\ge0}L^v
   \exp\left[-\frac{c_{\rm t}\ell}{\lambda_{k+1}}
                          (L^{v+1}-\tfrac12)\right]
 \le\sum_{v\ge0}L^ve^{-2c_{\rm t}(L^{v+1}-1/2)},
\]
because $\ell/\lambda_{k+1}\ge2$.
Consequently
\begin{equation}\label{eq:qcdm2-old-cubic}
  \|V_F\|\le c_\psi\left((1+c_{\rm cs})C_{\rm pc}S_{1/4}E_3
         +C_{\rm hl}C_{\mathfrak d}S_{\rm tor}\right)q_k,
 \qquad
 S_{\rm tor}=\sum_{v\ge0}L^v
       e^{-2c_{\rm t}(L^{v+1}-1/2)}<\infty.
\end{equation}
The factor from the first variation is $(17/16)K_M$.

\emph{The remaining coefficients.}
For a multi-index of weighted degree at least three the
Cauchy product has at least three powers of $\phi_k$.
Summing these bounds uses
\[
 N_{10}=\sum_{m\ge3}\binom{m+10}{10}2^{-(m-3)}<\infty.
\]
It bounds their total norm by
$4K_Mc_{\rm cl}N_{10}(C_r^{\rm b})^3\phi_k^3$ after
$P_{\rm col}R^{\rm mix}$.
Adding their degree-at-least-two parts to the first
variation gives a pre-extraction representative $R$ with
\[
 \|R\|_{k+1}\le
 \left[\tfrac{17}{4}K_Mc_\psi
       \bigl((1+c_{\rm cs})C_{\rm pc}S_{1/4}E_3
                    +C_{\rm hl}C_{\mathfrak d}S_{\rm tor}\bigr)
       +4K_Mc_{\rm cl}N_{10}(C_r^{\rm b})^3\right]q_k.
\]
For clarity, the pre-extraction representative in this
display is the explicitly defined family
\[
 R(X)=\bigl(R^{\rm mix}G_F(X)\bigr)_{\ge2}
   +\sum_{\substack{v\ \text{mixed}\\|v|_w\ge3}}
          P_{\rm col}R^{\rm mix}(a_v(X))_{\ge2}.
\]
It has the total of the mixed activity with legs,
because all omitted low-weight totals were just shown
to vanish or to be scalar.  Each retained term is
neutral in both species.  The first-variation family
is formed from colour-invariant representatives, so
it is also invariant under $P_{\rm col}$.

The scalar estimate can already be closed.  The
normalizers and Gaussian logarithms have no mixed
part.  In degree zero the colour projection is the
identity, and the coefficient bound for $R^{\rm mix}$
costs at most four.  The scalar output is therefore
bounded by the first-variation contribution, the
sum of coefficients of weight at least three, and
\eqref{eq:qcdm2-two-scalar-coefficient}.  Their sum is
\[
 M^{-3}n_{k+1}\left[
       (E_3/2+A_3)q_k
       +\frac{256K_Mv_{\mathsf S}^2}{c_{\rm cl}}\mathfrak l_k^2
                       \right]
       \le E_3M^{-3}n_{k+1}\bar q_k,
\]
with the constants chosen below.  This proves
\eqref{eq:qcdm2-mixed-total}; it also explains why the
scalar term $\mathfrak l_k^2$ need not enter the size
of the representative with fermion legs.

\emph{Extraction and cell sums.}
In the following estimates put $j=k+1$.
Write $K^R_X=\sum_pK^R_{X,p}$ for the quadratic kernel on
$X$, decomposed by expanding the
two background legs into their local and walk terms, and
let $Z_p$ be their supporting polymers. For such a kernel
define
\[
 M_y[K]=\frac12\sum_{y'}\bigl(K(y,y')+K(y',y)\bigr),
\]
with columns transported to $y$ inside the fundamental
cell centered at $y$.  This is a linear function of $K$.
The moved
coefficients are $a_{X,p}(y)=M_y[K^R_{X,p}]$.
The row-sum bound from the proof of
\cite[Prop. 10.1]{QCDM1} is
\[
 \sum_{p,y}\|a_{X,p}(y)\|
 e^{\kappa[d_j(Z_p)-d_j(X)]}
 \le CA_\kappa^2h_F^{-2}\|R(X)\|.
\]
After summing over $X$, it gives, for every cell $c$,
every cube $\square$, and $n\ge1$,
\begin{align}
 \sum_{y\in c,X,p}\|a_{X,p}(y)\|
 e^{\kappa L^{-n}d_j(Z_p)}(1+d_j(Z_p))^3
 &\le\tfrac12C_{\rm pw}h_F^{-2}\|R\|_j,
 \label{eq:qcdm2-moved-cell-sum}\\
 \sum_{Z_p\supset\square,y}\|a_{X,p}(y)\|
 e^{\kappa d_j(Z_p)}
 &\le\tfrac12C_{\rm pw}h_F^{-2}\|R\|_j.
 \label{eq:qcdm2-moved-polymer-sum}
\end{align}
The factor $1/2$ appears because $C_{\rm pw}$ in that proof
bounds the two families at the source and at zero together;
here there is one family. At lag zero use
\eqref{eq:qcdm2-moved-polymer-sum}.
Equality of the totals
identifies the sum of these moments with the moment of
the actual total.  This does not identify the moments
polymer by polymer.  The source increment uses only
polymers of diameter less than one quarter of the torus
side and reads their moments along lifted paths.  Denote
the actual mixed quadratic kernels by $K_{X,p}$.
By \cite[Prop. 10.1, equation (35)]{QCDM1},
\[
 -\mathfrak s_jR^{\rm mix}\nu_j(y)
   =\sum_{\operatorname{diam}X<\ell/4,p}
                  \operatorname{ord}_0[K_{X,p}](y),
\]
where $\mathfrak s_j\ge\eta_{\min}^2$ is the zero-source
normalization of the extraction map.  On a nonwrapping
$Z_p$, the moment $\operatorname{ord}_0$ read along the
lifted paths agrees with $M_y$.
Equality of the total kernels therefore gives the exact
identity
\begin{equation}\label{eq:qcdm2-extraction-discrepancy}
 -\mathfrak s_jR^{\rm mix}\nu_j(y)
       =\sum_{X,p}M_y[K^R_{X,p}]-d(y),
\end{equation}
where
\[
 d(y)=\sum_{X,p}'M_y[K_{X,p}]
                 -\sum_{X,p}''\operatorname{ord}_0[K_{X,p}](y).
\]
The first sum contains actual pieces for which either
$\operatorname{diam}X\ge\ell/4$ or $Z_p$ wraps around the
torus.  The second contains pieces with
$\operatorname{diam}X<\ell/4$ whose $Z_p$ wraps around
the torus.  This formula accounts for both the omitted
large polymers and the difference between the two
moment conventions on wrapping paths.
Every supporting polymer in this difference satisfies
\[
 d_{k+1}(Z_p)\ge
 \frac{\ell}{4c_{\mathcal D}M\lambda_{k+1}}-1.
\]
Here $c_{\mathcal D}$ is the fixed diameter constant,
$\operatorname{diam}Y\le c_{\mathcal D}M\lambda_j(1+d_j(Y))$.
Indeed, either the original polymer has diameter at least
one quarter of the torus side, or $Z_p$ contains a path
wrapping around the torus. The tail of
\eqref{eq:qcdm2-moved-cell-sum} therefore gains
$e^{-c_{\rm t}\ell/\lambda_{k+1}}$, by the choice of
$c_{\rm t}$ in \cite[Prop. 10.1]{QCDM1}.
Apply this estimate to the actual mixed quadratic activity,
whose norm is at most $4C_2^{\rm s}\phi_k^2$ by
\cref{C-407-42}.
Consequently
\begin{equation}\label{eq:qcdm2-seam-discrepancy}
 \sum_{y\in c}|d(y)|
 \le4C_{\rm pw}h_F^{-2}C_2^{\rm s}\phi_k^2
                       e^{-c_{\rm t}\ell/\lambda_{k+1}}.
\end{equation}
We now perform the extraction on the representative.
With $\psi^{\rm bg}=H_{0,j}\Psi$, define the quadratic
forms
\begin{align*}
 P^R_{z;X,p}
  &=-\mathfrak s_j^{-1}L^{-j}
          \mathsf S_1\bigl(M_z[K^R_{X,p}]1_z\psi^{\rm bg}\bigr),\\
 P^d_z
  &=\mathfrak s_j^{-1}L^{-j}
          \mathsf S_1\bigl(d(z)1_z\psi^{\rm bg}\bigr).
\end{align*}
Here $1_z$ restricts the coefficient to $z$ and
$\mathsf S_1$ is the mass bilinear used in the source
extraction.  Assign both forms to the cell containing
$z$; their supporting polymers are $Z_p$ and the cube
of $z$, respectively.  By
\eqref{eq:qcdm2-extraction-discrepancy},
\[
 \sum_{z;X,p}P^R_{z;X,p}+\sum_zP^d_z
       =L^{-j}\mathsf S_1(R^{\rm mix}\nu_j\psi^{\rm bg}).
\]
This is precisely the local form that the actual
extraction adds to the mixed activity.  Thus
\[
 F_{\mathfrak c}^{(j)}
       =R+\sum_{z;X,p}P^R_{z;X,p}+\sum_zP^d_z
\]
represents the required extracted total.  Its forms
are neutral in both species and invariant under
$P_{\rm col}$, since their coefficients are split and
the reconstruction maps at the trivial pair act
trivially on colour.
The first forms add at most
$\eta_{\min}^{-2}c_{\mathsf S}M^{-3}C_{\rm pw}\|R\|_{k+1}$
to the polymer norm by
\eqref{eq:qcdm2-moved-polymer-sum}.
The correction forms add at most
\[
 4\eta_{\min}^{-2}c_{\mathsf S}M^{-3}C_{\rm pw}
 C_2^{\rm s}\phi_k^2e^{-c_{\rm t}\ell/\lambda_{k+1}}.
\]
Thus a common bound for the norm and the cell-unit
hypotheses is
\begin{equation}\label{eq:qcdm2-extraction-size}
 \epsilon_{k+1}^{(3)}
 =K_\mu^{\rm s}\|R\|_{k+1}
 +4(1+\eta_{\min}^{-2}c_{\mathsf S}M^{-3}C_{\rm pw})
 C_2^{\rm s}\phi_k^2e^{-c_{\rm t}\ell/\lambda_{k+1}}.
\end{equation}
For the unit hypotheses, split the quadratic kernel into
$K^R_X-\delta_{yy'}M_y[K^R_X]$, assigned to the cell of
$X$, and the moved diagonal pieces, assigned to the cell
of $y$. Their variation and first moments obey
\eqref{eq:qcdm2-cell-unit-bound} by the kernel bound and
\eqref{eq:qcdm2-moved-cell-sum}.
The coefficients of the two types of extraction forms,
summed over a cell, are bounded respectively by
\[
 \tfrac12\eta_{\min}^{-2}C_{\rm pw}h_F^{-2}\|R\|_{k+1},
 \qquad
 4\eta_{\min}^{-2}C_{\rm pw}h_F^{-2}C_2^{\rm s}
 \phi_k^2e^{-c_{\rm t}\ell/\lambda_{k+1}}.
\]
Each is at most
$\eta_{\min}^{-2}C_{\rm pw}h_F^{-2}\epsilon_{k+1}^{(3)}$.
The moved coefficient is exactly one half of a row sum
plus one half of a column sum of $K^R_{X,p}$, as required
in \cref{C-415-9}; its second family $G$ is zero.

For a nonwrapping $Z_p$, the zeroth moment of
$K^R_{X,p}-\delta_{yy'}M_y[K^R_{X,p}]$ is zero. The
two moment conventions agree on this support, so the
subtracted diagonal has exactly the same total moment
as the original kernel.  A wrapping piece contributes
at most twice the double sum of its kernel.  Assigning
it to the hull of $Z_p$ and using the tail of
\eqref{eq:qcdm2-moved-cell-sum} bounds these contributions
in the cell norm by
\[
 C_{\rm pw}h_F^{-2}\|R\|_j
                      e^{-c_{\rm t}\ell/\lambda_j}.
\]
A moved piece and its extraction form have joint
coefficient
\[
 \left(1-\mathfrak s_{k+1}^{-1}
       \sum_yA_{k+1}^{\mathsf n}(y,z)\right)
 M_z[K^R_{X,p}],
\]
where $A_j^{\mathsf n}$ is the zeroth-moment map of the
background mass form.  Its source variation and its
torus comparison from \cite[proof of Prop. 10.1]{QCDM1}
give the bound
\[
 \eta_{\min}^{-2}
    \left(C_{\rm bg}^{\rm w}\frac{L^j\|\mathsf n\|}{r_M}
                 +C_Ae^{-4c_{\rm t}\ell/\lambda_j}\right)
       \frac12C_{\rm pw}h_F^{-2}\|R\|_j
\]
for the grouped moved pieces and their forms.  Keeping
them grouped is what produces the small coefficient
$1-\mathfrak s_j^{-1}\sum_yA_j^{\mathsf n}(y,z)$.

The correction form $P^d_z$ contributes
$\mathfrak s_j^{-1}\sum_yA_j^{\mathsf n}(y,z)d(z)$.
At zero source,
\[
 \left|\sum_yA_j^0(y,z)-\mathfrak s_j\right|
                    \le C_Ae^{-4c_{\rm t}\ell/\lambda_j}.
\]
The change with $\mathsf n$, expanded in the local and
walk terms of the two background legs, is at most
$C_{\rm bg}^{\rm w}L^j\|\mathsf n\|/r_M$ in the same
weighted cell norm.  Since $L^j\|\mathsf n\|/r_M<2$,
\eqref{eq:qcdm2-seam-discrepancy} bounds these correction
forms by
\[
 C_{\mathfrak d}h_F^{-2}\phi_k^2
                     e^{-c_{\rm t}\ell/\lambda_j},\qquad
 C_{\mathfrak d}
 =4\bigl(1+\eta_{\min}^{-2}(C_A+2C_{\rm bg}^{\rm w})\bigr)
                           C_{\rm pw}C_2^{\rm s}.
\]
Combining the three types of terms proves
\eqref{eq:qcdm2-representative-cells}.

Choose
\[
 \begin{split}
 C_{\mathfrak d}
 &=4(1+\eta_{\min}^{-2}(C_A+2C_{\rm bg}^{\rm w}))
                                      C_{\rm pw}C_2^{\rm s},\\
 A_3&=K_\mu^{\rm s}\left(
       \tfrac{17}{4}K_Mc_\psi C_{\rm hl}C_{\mathfrak d}S_{\rm tor}
           +4K_Mc_{\rm cl}N_{10}(C_r^{\rm b})^3\right)\\
 &\hspace{1em}
       +4(1+\eta_{\min}^{-2}c_{\mathsf S}M^{-3}C_{\rm pw})C_2^{\rm s},\\
 E_3&=2A_3+256K_Mv_{\mathsf S}^2/c_{\rm cl}.
 \end{split}
\]
After extraction the coefficient of $E_3q_k$ is at most
$4\theta_E^{\rm s}(1+c_{\rm cs})C_{\rm pc}S_{1/4}\le1/2$.
Thus the norm is at most $(E_3/2+A_3)q_k$.
These choices close both \eqref{eq:qcdm2-representative-size} and
\eqref{eq:qcdm2-mixed-total}.  Finally,
\eqref{eq:qcdm2-extraction-discrepancy}, the first
moment bound and \eqref{eq:qcdm2-seam-discrepancy} give
\begin{align*}
 |R^{\rm mix}\nu_j(y)|
 &\le\eta_{\min}^{-2}C_{\rm pw}h_F^{-2}
     \left(\tfrac12\|R\|_j
           +4C_2^{\rm s}\phi_k^2
                         e^{-c_{\rm t}\ell/\lambda_j}\right)\\
 &\le N_3q_k,
 \qquad
 N_3=\eta_{\min}^{-2}C_{\rm pw}h_F^{-2}(2E_3+4C_2^{\rm s}).
\end{align*}
To read this as an estimate for one species, apply
$R^{\rm mix}$ to its local quadratic term.  For
$b'\ne b$ the result is
\[
 R^{\rm mix}(\bar\psi_b\nu_{j,b}(\mathfrak c)\psi_b)
   =\bar\psi_b\bigl(\nu_{j,b}(\mathfrak c)
                    -\nu_{j,b}(R_{b'}\mathfrak c)\bigr)\psi_b.
\]
The preceding kernel estimate is therefore
\eqref{eq:qcdm2-cross-increment}.
\end{proof}

\subsection{The source chain and the sum over scales}

The preceding lemma concerns free split coordinates.  The
vacuum energy is evaluated along the actual source chain.
We next control the change in one chain when the source in
the other species is added.

\begin{lemma}\label{C-415-11}
For split sources and $\lambda_i\|J_b\|<r_M$,
\begin{equation}\label{eq:qcdm2-chain-feedback}
 \|\sigma_{i,1}(J_1,J_2)-\sigma_{i,1}(J_1,0)\|
 \le C_{\rm s3}q_{i-1}\frac{\lambda_i\|J_2\|}{r_M},
 \qquad C_{\rm s3}=\frac{3N_3}{1-4/L},
\end{equation}
and the analogous assertion holds for species two.
If two inner coordinates have the same $\mathsf n$, and
\[
 T=\max_{i\le j-1}
       \frac{\|\sigma_i-\sigma_i'\|}{B_\sigma\phi_i}\le1,
\]
then
\begin{equation}\label{eq:qcdm2-scalar-Lipschitz}
 |f_j(\mathfrak c)-f_j(\mathfrak c')|
 \le8K_MM^{-3}n_jC_{\rm Lip}\phi_{j-1}T,
\end{equation}
where
$C_{\rm Lip}=c'_\psi c_{\mathsf S}h_F^2B_\sigma
 +2c'_\psi\bar\Sigma+2\sqrt2A_1K_0^{\rm s}$.
\end{lemma}

\begin{proof}
The chains exist and are holomorphic on
$\lambda_j\|J_b\|<2r_M$ by \cref{C-407-42}.
On the smaller domain in the statement their coordinates
are split and lie in the inner coordinate domain.
Write
\[
 \mathfrak c_i^+=\mathfrak c_i(J_1,J_2),\qquad
 \mathfrak c_i^0=\mathfrak c_i(J_1,0),\qquad
 \Delta_l=\sigma_{l,1}(J_1,J_2)-\sigma_{l,1}(J_1,0).
\]
The bare first-species coordinate is $-aJ_1$ in both
chains.  It cancels from their difference, leaving
\begin{equation}\label{eq:qcdm2-feedback-recurrence}
 \Delta_l=\sum_{i=1}^l L^{l-i}D_i,\qquad
 D_i=\nu_{i,1}(\mathfrak c_{i-1}^+)
                  -\nu_{i,1}(\mathfrak c_{i-1}^0).
\end{equation}
Insert the coordinates with their second block set to
zero into each increment:
\begin{align*}
 D_i={}&
  [\nu_{i,1}(\mathfrak c_{i-1}^+)
       -\nu_{i,1}(R_2\mathfrak c_{i-1}^+)]\\
 &+[\nu_{i,1}(R_2\mathfrak c_{i-1}^+)
       -\nu_{i,1}(R_2\mathfrak c_{i-1}^0)]\\
 &+[\nu_{i,1}(R_2\mathfrak c_{i-1}^0)
       -\nu_{i,1}(\mathfrak c_{i-1}^0)].
\end{align*}
The first and third brackets are each bounded by
$N_3q_{i-1}$, by \eqref{eq:qcdm2-cross-increment}.
The third bracket must be retained: even when the bare
source $J_2$ is zero, the chain can have nonzero
second-species mass corrections.
The two coordinates in the middle bracket have the same
bare coordinate, and their earlier first-species entries
differ by $\Delta_l$, $l<i$.  Hence
\eqref{eq:qcdm2-quadratic-increment} gives
\begin{equation}\label{eq:qcdm2-feedback-increment}
 \|D_i\|\le2N_3q_{i-1}
   +B_\nu^{(2)}\phi_{i-1}^2
       \max_{l<i}\frac{\|\Delta_l\|}{B_\sigma\phi_l}.
\end{equation}
The maximum is zero at $i=1$.

Let $d_i$ be the supremum of $\|D_i\|$ on
$\lambda_i\|J_2\|<r_M$, with the admissible $J_1$ fixed.
The increment difference is holomorphic in $J_2$ and
vanishes at $J_2=0$.  Schwarz's lemma on each complex
line through the origin, followed by
\eqref{eq:qcdm2-feedback-recurrence}, gives
\begin{equation}\label{eq:qcdm2-feedback-Schwarz}
 \|D_i\|\le d_i\frac{\lambda_i\|J_2\|}{r_M},\qquad
 \|\Delta_l\|\le
      \frac{\lambda_l\|J_2\|}{r_M}\sum_{i=1}^l d_i.
\end{equation}
Here the identity $L^{l-i}\lambda_i=\lambda_l$ accounts
for the scale factor in the second inequality.

We prove $d_i\le3N_3q_{i-1}$ by induction.
The scale comparison from \cite[Prop. 10.1]{QCDM1}
gives $\phi_{i'-1}\le q_\phi^{l-i'}\phi_{l-1}$
with $q_\phi^2\le4/L$.  The exponential in $q_{i'-1}$
is at most the corresponding exponential in $q_{l-1}$.
It follows that
\begin{equation}\label{eq:qcdm2-feedback-scale-sum}
 \sum_{i'=1}^l q_{i'-1}\le\frac{q_{l-1}}{1-4/L}.
\end{equation}
When $\lambda_i\|J_2\|<r_M$ and $l<i$, the earlier
scale satisfies $\lambda_l\|J_2\|<r_M/L$.
The induction hypothesis and
\eqref{eq:qcdm2-feedback-Schwarz} therefore imply
\[
 \|\Delta_l\|\le
       \frac{3N_3}{L(1-4/L)}q_{l-1}.
\]
Moreover,
\[
 \frac{q_{l-1}}{\phi_l}
 \le\phi_{i-1}
       \left(\phi_{i-1}
             +e^{-c_{\rm t}\ell/(2\lambda_i)}\right),
 \qquad l<i.
\]
Thus the last term of
\eqref{eq:qcdm2-feedback-increment} is at most
\[
 \frac{3B_\nu^{(2)}\phi_{i-1}N_3}
      {L(1-4/L)B_\sigma}\,q_{i-1}
 \le N_3q_{i-1}.
\]
This proves the induction.  Substitution into
\eqref{eq:qcdm2-feedback-Schwarz} and
\eqref{eq:qcdm2-feedback-scale-sum} proves
\eqref{eq:qcdm2-chain-feedback}; the other species is
treated by interchanging the two blocks.
The smallness conditions used here and below are
\[
 3B_\nu^{(2)}\phi_*\le L(1-4/L)B_\sigma,
 \qquad 2C_{\rm s3}\phi_*\le B_\sigma.
\]

For the scalar estimate put
$\Delta V=V(\mathfrak c')-V(\mathfrak c)$.
The two coordinates have the same $\mathsf n$, so their
fermion covariances, reconstruction maps, Gaussian terms
and normalizing factors coincide.  The potential can
change only through the mass member and the old
activities.  The mass difference is bounded by
\[
 c'_\psi c_{\mathsf S}h_F^2
        \|\sigma_{j-1}'-\sigma_{j-1}\|
 \le c'_\psi c_{\mathsf S}h_F^2B_\sigma\phi_{j-1}T.
\]
If $T=0$ there is nothing to prove.  Otherwise, the line
\[
 \mathfrak c+\zeta(\mathfrak c'-\mathfrak c)/T,
 \qquad |\zeta|<1,
\]
lies in the doubled coordinate domain.  On this line
the old members with fermion variables have norm at most
$c'_\psi\bar\Sigma\phi_{j-1}^2$, while their scalar
members have norm at most $K_0^{\rm s}\delta_{j-1}$.
Schwarz's lemma applied to each difference from its value
at $\zeta=0$ supplies the factor $2T$.  Consequently
\begin{align*}
 \|\Delta V\|
 &\le\left(c'_\psi c_{\mathsf S}h_F^2B_\sigma\phi_{j-1}
       +2c'_\psi\bar\Sigma\phi_{j-1}^2
       +2K_0^{\rm s}\delta_{j-1}\right)T\\
 &\le C_{\rm Lip}\phi_{j-1}T,
\end{align*}
where we used $\phi_{j-1}\le1$ and
$\delta_{j-1}\le\sqrt2 A_1\phi_{j-1}$.

At the trivial pair the potential $V(\mathfrak c)$ has
norm at most $11c_{\rm cl}/16$.  On the disc
\[
 |v|<\frac{c_{\rm cl}}{4\|\Delta V\|},
\]
the potential $V(\mathfrak c)+v\Delta V$ has norm at
most $15c_{\rm cl}/16$.  The tail has norm at most
$c_{\rm cl}/16$, so the cluster theorem applies throughout
this disc.  The scalar total is holomorphic there and
bounded by $K_Mc_{\rm cl}n_jM^{-3}$: the local cluster
bound is summed over the $n_jM^{-3}$ cubes.
Schwarz's lemma for its difference from the value at
$v=0$ gives
\[
 |f_j(\mathfrak c')-f_j(\mathfrak c)|
 \le2K_Mc_{\rm cl}n_jM^{-3}
                \frac{4\|\Delta V\|}{c_{\rm cl}}.
\]
If the displayed radius is at most one, the same estimate
follows from the individual bounds at the two admissible
endpoints.  Substituting the bound for $\|\Delta V\|$ proves
\eqref{eq:qcdm2-scalar-Lipschitz}.
\end{proof}

\begin{proof}[Proof of \cref{C-415-5}]
Let $\mathfrak c_{j-1}(J_1,J_2)$ be the source-chain
coordinate, and write
\[
 G_j(J_1,J_2)
   =D^{\rm mix}[f_j\circ\mathfrak c_{j-1}](J_1,J_2).
\]
This is the sum over polymers of the mixed step energy.
The bound for $R^{\rm mix}f_j$ does not apply to $G_j$
directly: a single-source part of $f_j$ can still depend
on the other bare source through the mass corrections
in its own species.  We compare with a coordinate in
which the two species depend on separate sources.

Define $\mathfrak c^\flat(J_1,J_2)$ by taking every
first-species entry from $\mathfrak c_{j-1}(J_1,0)$
and every second-species entry from
$\mathfrak c_{j-1}(0,J_2)$.  It is in the inner domain,
and its bare coordinate agrees with that of
$\mathfrak c_{j-1}(J_1,J_2)$.  In the decomposition
\[
 f_j=f_j(0)+R^{(1)}f_j+R^{(2)}f_j+R^{\rm mix}f_j,
\]
the composition of $R^{(b)}f_j$ with $\mathfrak c^\flat$
depends only on $J_b$.  Further,
$R^{\rm mix}f_j(\mathfrak c^\flat)$ vanishes on each
source axis, because the corresponding block of
$\mathfrak c^\flat$ is taken from the zero-source chain.
Consequently
\[
 D^{\rm mix}[f_j(\mathfrak c^\flat)]
            =R^{\rm mix}f_j(\mathfrak c^\flat).
\]

At each corner of the mixed difference the original and
separated coordinates have the same bare coordinate.
At $(J_1,J_2)$ their two mass corrections differ by
\[
 \sigma_{i,1}(J_1,J_2)-\sigma_{i,1}(J_1,0),\qquad
 \sigma_{i,2}(J_1,J_2)-\sigma_{i,2}(0,J_2).
\]
At $(J_1,0)$ only the second block differs, by
$\sigma_{i,2}(J_1,0)$; at $(0,J_2)$ only the first
block differs, by $\sigma_{i,1}(0,J_2)$.
At $(0,0)$ there is no difference.  The axis differences
are also bounded by
\eqref{eq:qcdm2-chain-feedback}, with the other source
set to zero.  Thus at each of the three nonzero corners
the normalized coordinate difference satisfies
\[
 T\le\frac{C_{\rm s3}}{B_\sigma}
            \max_{i\le j-1}\frac{q_{i-1}}{\phi_i},
 \qquad
 \phi_{j-1}T\le\frac{C_{\rm s3}}{B_\sigma}q_{j-1}.
\]
For the second inequality, use
\[
 \frac{q_{i-1}}{\phi_i}
 \le\phi_{i-1}
       \left(\phi_{i-1}+e^{-c_{\rm t}\ell/(2\lambda_i)}\right)
 \le\phi_{j-1}
       \left(\phi_{j-1}+e^{-c_{\rm t}\ell/(2\lambda_j)}\right).
\]
The condition $2C_{\rm s3}\phi_*\le B_\sigma$ gives
$T\le1$.  Apply \eqref{eq:qcdm2-scalar-Lipschitz}
at the three corners and
\eqref{eq:qcdm2-mixed-total} to $\mathfrak c^\flat$.
Since $q_{j-1}\le\bar q_{j-1}$, we obtain
\begin{align*}
 |G_j(J_1,J_2)|
 &\le E_3M^{-3}n_j\bar q_{j-1}
       +3\cdot8K_MM^{-3}n_jC_{\rm Lip}
                    \frac{C_{\rm s3}}{B_\sigma}q_{j-1}\\
 &\le C_{\rm gm3}M^{-3}n_j\bar q_{j-1},
\end{align*}
where
\[
 C_{\rm gm3}=E_3+24K_MC_{\rm Lip}C_{\rm s3}/B_\sigma.
\]
This estimate holds throughout $\lambda_j\|J_b\|<r_M$.
For nonzero $J_1,J_2$, set
\[
 G(\zeta_1,\zeta_2)=G_j(\zeta_1J_1,\zeta_2J_2),\qquad
 R_b=\frac{r_M}{\lambda_j\|J_b\|}.
\]
The function $G$ is holomorphic on $|\zeta_b|<R_b$
and zero on both coordinate axes.  First apply Schwarz's
lemma in $\zeta_1$.  The quotient $G/\zeta_1$ has a
holomorphic extension through $\zeta_1=0$, is bounded
by $\sup|G|/R_1$, and is zero when $\zeta_2=0$.
A second application in $\zeta_2$ gives
\[
 |G_j(J_1,J_2)|\le C_{\rm gm3}M^{-3}n_j\bar q_{j-1}
      \frac{\lambda_j\|J_1\|}{r_M}
      \frac{\lambda_j\|J_2\|}{r_M}.
\]
If either source is zero, this inequality is immediate.
Since the energy contribution is
$f_j(\mathfrak c_{j-1}(J))-f_j(0)$, the inequality is
\eqref{eq:qcdm2-mixed-step}.

For the sum over scales, the two source factors supply
$\lambda_j^2$, whereas the number of sites contributes
$n_j=\ell^3\lambda_j^{-3}$.  Their product leaves one
inverse length:
\[
 n_j\lambda_j^2=\ell^3\lambda_j^{-1}
                      =\ell^3L^{K-j}.
\]
Since $\rho_*/r_M\le(2L)^{-1}$, it follows that
\begin{equation}\label{eq:qcdm2-energy-before-sum}
 |E_K^{\rm mix}(J_1,J_2)|
 \le C_{\rm gm3}M^{-3}(2L)^{-2}\ell^3
                       \sum_{j=1}^K L^{K-j}\bar q_{j-1}.
\end{equation}
The three terms in $\bar q_{j-1}$ are summed separately.
Put $n=K-j$.  Then
$\mathsf t_{j-1}=L^{n+1}\mathsf b$ and
\[
 x_{j-1}\ge\mathsf t_{j-1}-\bar B
              \ge\tfrac12 L^{n+1}\mathsf b,
\]
because $\mathsf b>\gamma^{-2}\ge L\bar B$.
The function $\phi(t)=t^{-1/2}(1+\log t)^{p_0}$ is
decreasing on $[e^{2p_0},\infty)$, which contains
$[\frac12\gamma^{-2},\infty)$ by the choice of $\gamma$.
Consequently
\begin{align*}
 \phi_{j-1}
 &\le\phi(\tfrac12L^{n+1}\mathsf b)\\
 &\le\sqrt2L^{-(n+1)/2}\phi(\mathsf b)
                       (1+(n+1)\log L)^{p_0}.
\end{align*}
For the second inequality use
$1+a+b\le(1+a)(1+b)$ for nonnegative $a,b$ on the
logarithmic factor.  Since $\phi(\mathsf b)\le\phi_*\le1$,
\[
 L^n\phi_{j-1}^3
 \le2^{3/2}L^{-3/2}L^{-n/2}
                       (1+(n+1)\log L)^{3p_0}.
\]
The geometric factor $L^{-n/2}$ makes this series
summable despite its logarithmic polynomial.

For the scalar covariance left at weight two, the
definition of $\mathfrak l_{j-1}$ gives the exact identity
\[
 L^n\mathfrak l_{j-1}^2
  =\frac{1}{L\mathsf b}
       (\log\mathsf b+(n+1)\log L)^{-2p_0}
  \le L^{-1}((n+1)\log L)^{-2p_0}.
\]
The scale powers cancel. Since $2p_0>1$, the upper bounds sum to
\[
 \frac{\zeta(2p_0)}{L(\log L)^{2p_0}}.
\]
Finally, $\ell\ge2$ and $\phi_{j-1}\le1$ imply
\[
 L^n\phi_{j-1}^2e^{-c_{\rm t}\ell/(2\lambda_j)}
   =L^n\phi_{j-1}^2e^{-c_{\rm t}\ell L^n/2}
   \le L^ne^{-c_{\rm t}L^n}.
\]
This series is summable because its exponential contains
$L^n$.  Thus all three contributions in
\eqref{eq:qcdm2-energy-before-sum} are bounded uniformly
in $K$ and $\ell$.  One may take
\begin{align*}
 S_{\rm mix}
 &=2^{3/2}L^{-3/2}\sum_{n\ge0}
        L^{-n/2}(1+(n+1)\log L)^{3p_0}
       +\frac{\zeta(2p_0)}{L(\log L)^{2p_0}}
       +\sum_{n\ge0}L^ne^{-c_{\rm t}L^n},\\
 C_{\rm mix}&=C_{\rm gm3}M^{-3}(2L)^{-2}S_{\rm mix}.
\end{align*}
This proves \eqref{eq:qcdm2-mixed-energy}.
\end{proof}

\paragraph{Order of the constants.}\label{C-415-12}
The mass window remains fixed. The new constants are chosen
after the mass-independent source norm and localization
constants of \cite[Prop. 10.20]{QCDM1}, and before
the final choice of $\gamma$. First choose
$K_0^{\rm s},C_r^{\rm s},C_2^{\rm s},B_\nu^{(2)}$;
then $\bar\Sigma,V_1,C_{\rm gm},v_{\mathsf S},C_r^{\rm b}$;
then $C_{\mathfrak d},A_3,E_3,N_3,C_{\rm s3},C_{\rm Lip}$
and $C_{\rm mix}$.  All are functions of constants already
chosen.  In addition to the three smallness conditions in
the proof of \cref{C-407-42}, impose
\[
 C_r^{\rm b}\phi_*\le\tfrac12,\qquad
 \tfrac{17}{2}\,16K_M^\bullet C_\Lambda\phi_*\le\tfrac12,
 \qquad
 3B_\nu^{(2)}\phi_*\le L(1-4/L)B_\sigma,
 \qquad2C_{\rm s3}\phi_*\le B_\sigma.
\]
They hold after decreasing $\gamma$, since $\phi_*\to0$.
They do not depend on the cutoff, volume, source or scale.

\section{Scalar mesons and diquarks}
\label{sec:qcd2-bilinears}

We retain the setting of \cref{sec:qcdm2-setting} and write
$D_U=D[U]+m_0$ for the massive operator. Constants may depend
on the fixed coupling interval and mass window, but are
independent of the cutoff, volume and infinite-volume state.
In this section and the next two, sites are labelled by
$x\in a\mathbb Z^3$. Thus $(x,q)$ corresponds to the place
$(a^{-1}x,q)$ in \eqref{eq:qcdm2-Dirac}, and $x_v$ denotes
the physical site of $v$.

The mesons $\bar\psi X\psi$ do not form an invariant family under
all internal symmetries of the model. Pseudoreality of the colour
representation permits a symmetry to replace a quark by a
charge-conjugate antiquark. We first include the resulting
diquarks among the fields, and then prove estimates for arbitrary
joint correlations of this larger family.

\subsection{The labels and the exact symmetries}

Put $\epsilon=i\sigma_2$ on colour and
$C_0=\sigma_2\otimes\epsilon$ on spin and colour. The identities
\[
 C_0^{\mathsf T}=C_0,\qquad C_0^{-1}=-C_0,
 \qquad C_0^\dagger=-C_0,\qquad
 \sigma_\mu^{\mathsf T}C_0^{-1}=-C_0^{-1}\sigma_\mu
\]
will be used below. With $\chi=C_0\bar\psi^{\mathsf T}$, write
$\widehat\Psi=(\psi,\chi)$. The internal label space of this
column is
\[
 W=\mathbb C^2_\tau\otimes\mathbb C^2_{\mathrm d}
                         \otimes\mathbb C^{N_f}.
\]
The second factor distinguishes $\psi$ from $\chi$.
Let $J=\left(\begin{smallmatrix}0&-1\\1&0\end{smallmatrix}\right)$
on this factor, let $\mathcal J=1\otimes J\otimes1$, and let
$P_\pm=(1\pm\tau_3)/2$. The complex label space and its
reflection-real form are
\begin{equation}\label{C-401-2}
 \mathscr S=\{M\in\operatorname{End}W:
                 \mathcal JM^{\mathsf T}\mathcal J^{-1}=M\},
 \qquad M^*=\tau_3M^\dagger\tau_3,
 \qquad\mathscr S_{\mathbb R}=\{M\in\mathscr S:M^*=M\}.
\end{equation}
Thus $\mathcal JM$ is antisymmetric and
$\dim_{\mathbb C}\mathscr S=2N_f(4N_f-1)$.
For $X$ on $\mathbb C^2_\tau\otimes\mathbb C^{N_f}$, let
$\widehat X$ act as $X$ on the $\psi$ components and as
$X^{\mathsf T}$ on the $\chi$ components. It belongs to
$\mathscr S$. When this notation is applied to a meson label
$\Gamma$, we suppress its identity factors on spin and colour.

For a chamber $v$ over the physical site $x_v$, define
\[
 \phi[M](v)=\frac12\widehat\Psi(v)^{\mathsf T}
                (\mathcal JM\otimes C_0^{-1})\widehat\Psi(v),
 \qquad
 \Phi[M](f)=a\sum_x f(x)\sum_{v:x_v=x}\phi[M](v).
\]
In particular $\phi[\widehat X]=\bar\psi X\psi$.
The blocks of $\mathcal JM$ with two $\psi$ indices give the
diquarks; those with two $\chi$ indices give the antidiquarks.
The label $0$ will continue to denote the action-density field
$O(f)$ and is not the zero matrix in $\mathscr S$.

Let
\begin{equation}\label{C-401-3}
 G_{\mathbb C}=\mathrm{Sp}(2N_f,\mathbb C)^2,
 \qquad G=\mathrm{USp}(2N_f)^2.
\end{equation}
Each factor acts on $P_\pm W$ with its symplectic form
$J\otimes1$. We regard $(S_+,S_-)$ as the block diagonal
matrix $S=P_+S_++P_-S_-$ on $W$, and define the algebra
automorphism $\alpha_S$ by
$\alpha_S\widehat\Psi(v)=S^{-1}\widehat\Psi(v)$.
It acts trivially on functions of the links.

\begin{proposition}[The doubled field and its symmetries]
\label{C-401-16}\label{C-401-17}\label{C-401-18}
At every cutoff the following identities hold on a torus and
for every infinite-volume state used in \cite{QCDM1}.
For $S\in G_{\mathbb C}$,
\[
 \alpha_S\phi[M]=\phi[SMS^{-1}],\qquad
 \omega\circ\alpha_S=\omega.
\]
At a lattice reflection with mirror $P$,
\[
 \Theta_P\Phi[M](f)=\Phi[M^*](\overline f\circ\theta_P).
\]
Every $\alpha_S$ with $S\in G$ commutes with $\Theta_P$.
All these fields are gauge invariant and even.

On a finite torus, let $\widehat D_U=\sum_{h=+,-}P_h\otimes D_{U,h}$ be the massive
operator on the doubled space and let $M$ be a function of the
chambers with values in $\mathscr S$. Its Berezin integral
\[
 Z_U[M]=\int d\bar\psi\,d\psi\,
       \exp\left(-\bar\psi D_U\psi-\sum_v\phi[M(v)](v)\right)
\]
is a polynomial satisfying
\[
 Z_U[M]^2=\det(\widehat D_U+M),\qquad Z_U[0]=\det D_U.
\]
These identities also hold for complex link variables whenever
the displayed expressions are defined.
\end{proposition}

\begin{proof}
The relation $\epsilon u^{\mathsf T}\epsilon^{-1}=u^{-1}$
for $u\in\mathrm{SL}(2,\mathbb C)$ and the Pauli identities
give
$C_0D_{U,h}^{\mathsf T}C_0^{-1}=D_{U,h}$, with the transpose
also interchanging the chamber arguments. Consequently
$K_h=C_0^{-1}D_{U,h}$ is symmetric. Since odd columns satisfy
$b^{\mathsf T}Aa=-a^{\mathsf T}A^{\mathsf T}b$,
\[
 \bar\psi_hD_{U,h}\psi_h
 =\frac12\widehat\Psi_h^{\mathsf T}(J\otimes K_h)
                                      \widehat\Psi_h.
\]
The full quadratic coefficient is therefore
$(\mathcal J\otimes C_0^{-1})(\widehat D_U+M)$,
which is antisymmetric. Its Berezin integral is a fixed sign
times its Pfaffian. Squaring gives the determinant identity;
the normalization at $M=0$ fixes the square of the sign.
For $\widehat X$, the two off-diagonal blocks of
$\mathcal J\widehat X$ are $-X^{\mathsf T}$ and $X$.
Reversing the odd variables in the first term gives
$\phi[\widehat X]=\chi^{\mathsf T}C_0^{-1}X\psi
=\bar\psi X\psi$.

The identity $S^{\mathsf T}\mathcal JS=\mathcal J$ gives
$S^{-\mathsf T}\mathcal J=\mathcal JS$. Hence substitution
of $S^{-1}\widehat\Psi$ fixes the quadratic action, because
$S$ commutes with $P_\pm$, and carries the source coefficient
$\mathcal JM$ into $\mathcal JSMS^{-1}$. Its Berezin
Jacobian is one. This proves invariance on a torus.
On the infinite lattice the doubled Gaussian functional has
covariance
\[
 \sum_h(-P_h\mathcal J)\otimes D_{U,h}^{-1}C_0.
\]
The same symplectic identity preserves this covariance and
therefore every Wick contraction. Averaging in $U$ proves
the assertion for the state.

If $\hat n$ is the unit normal to $P$, direct substitution in
the reflection of the generators gives
\[
 \Theta_P\widehat\Psi(v)
  =(\mathcal J\tau_3\otimes C_0^{-1}\sigma(\hat n))
                                      \widehat\Psi(\theta_Pv).
\]
For an antilinear map reversing the order of odd factors, the
quadratic coefficient $A$ is changed to $B^{\mathsf T}A^\dagger B$,
where $B$ is the displayed matrix. Inserting
$A=\mathcal JM\otimes C_0^{-1}$ gives
$\mathcal J(\tau_3M^\dagger\tau_3)\otimes C_0^{-1}$.
The same computation on the generators gives
$\Theta_P\alpha_S\Theta_P=\alpha_{(S^\dagger)^{-1}}$,
which equals $\alpha_S$ for $S\in G$.

Under a colour transformation, $\chi$ transforms by the same
matrix as $\psi$, because $\epsilon u^{-\mathsf T}\epsilon^{-1}=u$.
The form $C_0^{-1}$ is invariant under this transformation,
so every $\phi[M]$ is gauge invariant. Each is quadratic in
the generators and hence even.
\end{proof}

\begin{lemma}[A basis obtained from mesons]\label{C-401-19}
The real space $\mathscr S_{\mathbb R}$ has a basis
\[
 \mathfrak B^{\mathrm{ext}}
   =\{S_b\widehat\Gamma_bS_b^{-1}:b=1,\ldots,
                                      2N_f(4N_f-1)\},
 \qquad S_b\in G,\quad\Gamma_b\in\mathfrak B.
\]
The matrices in this basis have operator norm at most one.
There is a constant $c(N_f)$ such that the sum of the absolute
values of the coefficients of any $M\in\mathscr S$ in this
basis is at most $c(N_f)\|M\|$.
\end{lemma}

\begin{proof}
Let $V$ be the real span of the displayed symplectic images of
meson basis elements. Unitarity and $[S,\tau_3]=0$ imply
$(SMS^{-1})^*=SM^*S^{-1}$, so $V\subset\mathscr S_{\mathbb R}$.
We show directly that $V+iV=\mathscr S$.
Write $r=(h,f)$ for a species--flavour index. Under
$M\mapsto\mathcal JM$, the meson labels give all blocks
with one $\psi$ index and one $\chi$ index. Their matrix
units are complex linear combinations of the meson basis.
The remaining blocks are antisymmetric forms in two $\psi$
indices or two $\chi$ indices, hence are spanned by pairs
$r\ne s$. For such a pair let $S\in G$ be $J$ on the plane
spanned by $\psi_r,\chi_r$ and the identity elsewhere.
Since $S^{-1}=-J$ on that plane,
$\alpha_S\psi_r=\chi_r$ and $\alpha_S\chi_r=-\psi_r$.
Thus
\[
 \alpha_S(\bar\psi_r\psi_s)
       =-\psi_r^{\mathsf T}C_0^{-1}\psi_s.
\]
Applying partial charge conjugation to the other index
instead gives the corresponding antidiquark. These images
span every remaining block of $\mathcal JM$, proving the claim.

For the real-form conclusion, write any
$M\in\mathscr S_{\mathbb R}$ as $v+iw$, with $v,w\in V$.
Since $v^*=v$, $w^*=w$ and $M^*=M$, we have $w=0$.
Thus $V=\mathscr S_{\mathbb R}$. Select a basis from the spanning
set. Unitary conjugation preserves operator norm, and the
coefficient estimate is equivalence of norms on this fixed
finite-dimensional space.
\end{proof}

\subsection{Uniform estimates for the enlarged family}

Put $\mathcal I^{\mathrm{ext}}=\{0\}\cup\mathfrak B^{\mathrm{ext}}$,
$\|f\|_0=\|f\|_\sharp$ and $\|f\|_M=\|f\|_\infty$ for
$M\ne0$. An extended tuple is a list
$\mathcal X=((\lambda_i,f_i))_{i=1}^n$ with these labels;
its truncated function is the cumulant of the indicated fields.
We extend it complex multilinearly in the labels and functions.
Write
\[
 p=432C_{\mathscr C}/\log L,\qquad \kappa=5+p/2.
\]
Here $C_{\mathscr C}$ is the source-energy counting constant
in \cite[Thm. F]{QCDM1}.

\begin{proposition}[All scalar bilinears at separated supports]
\label{C-401-48}\label{C-400-33}
For every $n\ge2$, $R,d>0$ and extended tuple whose real
Lipschitz functions have supports in balls of radius $R$, at
pairwise distances at least $d$,
\begin{equation}\label{eq:qcd2-extended-bound}
 |S_n^T(\mathcal X)|
 \le [B_{\mathrm{ext}}(R)(1+d^{-1})^{\kappa-1}]^n
                     (n!)^\kappa\prod_i\|f_i\|_{\lambda_i}.
\end{equation}
This holds on every torus of the construction and in every
fixed-cutoff infinite-volume state obtained from these tori.
The constant $B_{\mathrm{ext}}(R)$ is independent of cutoff
and volume. For arbitrary labels in $\mathscr S$, the right
side has the additional factor
$\prod_{i:\lambda_i\ne0}c(N_f)\|\lambda_i\|$.
\end{proposition}

\begin{proof}
We first explain why the one-cell estimate of
\cite[Prop. 10.17]{QCDM1} applies
to a symplectic image of a meson and to cells of side less
than one. Let a torus be divided into cubes of physical side
$b=L^{-q}$, with an even number of cubes in each direction
and walls through links. For a source $s$ supported in one
cube at distance $\delta>0$ from its walls, the same fold
used in that proposition has $\mathsf N=\ell^3/b^3$ copies.
The upper and lower partition estimates cost
$c_V\ell^3/\mathsf N=c_Vb^3$ after taking the
$\mathsf N$th root. The separated-energy estimate gives
the same counting term
$\mathscr C_K(\operatorname{supp}s;\delta)$; its proof is
local at the support and uses the distance to the walls.
Thus, writing $\lambda^{\mathfrak b}$ for the average of
the reflected single-source energies,
\begin{equation}\label{eq:qcd2-small-cell}
 \|e^{O(w)+\Phi_\Gamma(c)}\|_{\mathrm{ch}}
 \le\exp\left(\operatorname{Re}\lambda^{\mathfrak b}(w,c\Gamma)
       +c_Vb^3+\mathscr C_K(\operatorname{supp}(w,c);\delta)\right).
\end{equation}
The required source domains are unchanged, since reflection
preserves the source norms and the source vanishes near the
walls. To specify the seminorm, let $\operatorname{Fold}G$
be the product of the reflected copies of an even element
$G$ from one cell into all $\mathsf N$ cells. For the normalized
torus functional $\omega_\ell$, put
\[
 \|G\|_{\mathrm{ch}}
       =\omega_\ell(\operatorname{Fold}G)^{1/\mathsf N}.
\]
Reflection positivity makes the quantity under this root
nonnegative. The multiple-cell chessboard inequality is
\[
 \left|\omega_\ell\left(\prod_iG_i\right)\right|
       \le\prod_i\|G_i\|_{\mathrm{ch}},
\]
when the $G_i$ belong to distinct cells; the unused cells
contain $1$. This is the same inequality as in
\cite[Prop. 10.17 and its proof]{QCDM1}, with the finer board.

For $M=S\widehat\Gamma S^{-1}$ in the chosen basis,
\[
 e^{O(w)+\Phi[M](c)}
       =\alpha_S(e^{O(w)+\Phi_\Gamma(c)}).
\]
The automorphism $\alpha_S$ commutes with every reflection
in the fold and preserves the torus functional. Therefore
the two exponentials have equal chessboard seminorms, and
\eqref{eq:qcd2-small-cell} bounds the extended source too.
To apply the multiple-cell inequality, first bound each
individual seminorm in this way and then multiply the bounds.
Thus each cell may use its own matrix $S$: each individual
fold is compared with its meson fold by one global symmetry.

Choose $q\ge0$ least such that $\sqrt3L^{-q}<d$, and set
$\eta=L^{-q}/(16n)$. Then
\[
 \eta^{-1}\le16\sqrt3Ln(1+d^{-1}).
\]
First suppose $L^{K-q}\ge16n$. Split each $f_i$ by the
product hat partition at mesh $\eta$. There are at most
$(2R/\eta+3)^3$ pieces for each entry. Each is supported in
a cube of side $2\eta$, its supremum norm does not increase,
and its sharp norm increases by at most $1+80M/\eta$.
Fix one piece from each entry. In each coordinate a piece
forbids at most $4\eta L^K+1$ positions for a wall of a
board whose cells have $\Delta=L^{K-q}$ sites on a side.
The $n$ pieces forbid at most $\Delta/4+n<\Delta$
positions. Choose walls at the remaining half-integral
positions. They are at distance at least $\eta$ from all
pieces. The diameter of a cell is less than $d$, so no
cell contains pieces from two different entries.

For the chosen pieces, let $s_i$ be the corresponding meson
source before its symplectic transformation, and let
$\Lambda(z)=\sum_i\lambda^{\mathfrak b}(z_is_i)$.
The extended generating function is entire, and the
chessboard estimate gives, on each coordinate face indexed
by $A\subset\{1,\ldots,n\}$,
\[
 \left|e^{-\Lambda(z)}\frac{Z^{\mathrm{ext}}(z)}{Z}\right|
 \le\exp\left(|A|c_V+
       \sum_{i\in A}\mathscr C_K(\operatorname{supp}s_i;\eta)\right),
  \qquad |z_i|<\frac{2\varrho_{\lambda_i}}{
                                      \|f_i\chi_i\|_{\lambda_i}}.
\]
Here $\varrho_0=\varrho$, $\varrho_M=\varrho_J$ for $M\ne0$,
and $b^3\le1$. Since every term of $\Lambda$ depends on
one variable, its mixed derivative in all $n\ge2$ variables
vanishes. The logarithmic derivative formula, as in the proof
of \cite[Thm. F]{QCDM1}, is
\[
 \partial_{[n]}\log Z^{\mathrm{ext}}(0)
 =\sum_\pi(-1)^{|\pi|-1}(|\pi|-1)!
       \prod_{A\in\pi}\partial_A
           \left(e^{-\Lambda}Z^{\mathrm{ext}}/Z\right)(0).
\]
The sum of the absolute combinatorial coefficients is at
most $2^nn!$. Writing $G=e^{-\Lambda}Z^{\mathrm{ext}}/Z$,
Cauchy's estimate on smaller coordinate circles, followed by
the limit as their radii increase to the displayed radii, gives
\[
 |\partial_AG(0)|\le
 \exp\left(|A|c_V+
       \sum_{i\in A}\mathscr C_K(\operatorname{supp}s_i;\eta)\right)
 \prod_{i\in A}\frac{\|f_i\chi_i\|_{\lambda_i}}
                            {2\varrho_{\lambda_i}}.
\]
In the product over a partition, each index occurs exactly
once. Hence both the source counts and the inverse Cauchy
radii have the same total cost for every partition.
The refined source count of that theorem gives
\[
 \mathscr C_K(\operatorname{supp}s_i;\eta)
    \le216C_{\mathscr C}(1+\log_L(1/\eta))+C,
\]
because the support has diameter bounded by a fixed
multiple of $\eta$. Hence its exponential contributes
at most $C^n\eta^{-np/2}$.

Summing over the pieces and restoring the original test
functions therefore gives
\[
 |S_n^T(\mathcal X)|\le n!
 \left[\frac{Ce^{c_V}}{\varrho_{\min}}
       (2R/\eta+3)^3(1+80M/\eta)\eta^{-p/2}\right]^n
       \prod_i\|f_i\|_{\lambda_i}.
\]
Here $\varrho_{\min}=\min(\varrho,\varrho_J)$.
The three powers of $\eta^{-1}$ count the pieces, one
comes from their sharp norms, and $p/2$ comes from the
source-energy count. Since $\eta^{-1}\le Cn(1+d^{-1})$,
the right-hand side is at most
\[
 \left[C'(1+R)^3(1+d^{-1})^{4+p/2}\right]^n
       n!\,n^{(4+p/2)n}\prod_i\|f_i\|_{\lambda_i}.
\]
Applying $n^n\le e^nn!$ gives the factorial exponent
$1+(4+p/2)=5+p/2$ in \eqref{eq:qcd2-extended-bound}.

If $L^{K-q}<16n$, then
$a^{-1}<16\sqrt3Ln(1+d^{-1})$. The Pfaffian identity
above and $\|\widehat D_U^{-1}\|\le m_0^{-1}$ give
\[
 |Z_U[M]|\le\det D_U\,
                \exp\left(\frac{\|M\|_1}{2m_0}\right).
\]
For the source of the extended tuple,
$M(v)=-a\sum_i z_if_i(x_v)\lambda_i$; the doubled fibre
has $16N_f$ components and there are $48$ chambers per
site. Thus the fermion contribution is bounded by
$384N_fm_-^{-1}\sum_x\|\sum_i z_if_i(x)\lambda_i\|$.
The action contribution is bounded by
$2a^{-1}\sum_p|w(p)|$.
A radius-$R$ support has at most $27a^{-3}(1+R)^3$ sites
and $81a^{-3}(1+R)^3$ plaquettes. Since $a\le1$,
both costs are bounded by
\[
 B_a(R)=C(1+N_f/m_-)(1+R)^3a^{-4}
\]
times the appropriate test-function norm. Set
$r_i=(B_a(R)\|f_i\|_{\lambda_i})^{-1}$, omitting zero
test functions. On the face $A$, the normalized generating
function has modulus at most $e^{|A|}$ for $|z_i|\le r_i$.
Cauchy's estimate and the cumulant partition formula give
\[
 |S_n^T(\mathcal X)|\le
          2^nn!e^nB_a(R)^n\prod_i\|f_i\|_{\lambda_i}.
\]
Use $a^{-1}<Cn(1+d^{-1})$ and $n^{4n}\le e^{4n}(n!)^4$
to obtain
\[
 |S_n^T(\mathcal X)|
 \le[C(1+N_f/m_-)(1+R)^3(1+d^{-1})^4]^n
                         (n!)^5\prod_i\|f_i\|_{\lambda_i}.
\]
This is bounded by the claimed estimate. Multilinearity
and \cref{C-401-19} give arbitrary labels.

Finally the expectation of any fixed local Grassmann
polynomial converges along the tori defining a state,
as proved in \cite[Prop. 2.5]{QCDM1}. A cumulant is a polynomial in finitely
many such expectations. Passing to the limit proves the
same estimate in the infinite-volume state.
\end{proof}

\begin{proposition}[The enlarged limits]\label{C-401-49}
Let $S$ be a volume-first limit of \cite{QCDM1}, with a
fixed defining sequence. There is a subsequence along
which all extended truncated functions of orders $n\ge2$ converge on
pairwise disjoint $C_c^1$ supports. Their limits
$S^{\mathrm{ext}}$ restrict to $S$ on meson tuples and
satisfy \eqref{eq:qcd2-extended-bound}.

They are symmetric, translation invariant, invariant
under simultaneous conjugation of the labels by
$G_{\mathbb C}$.
At fixed labels, these truncated functions determine
distributions of locally finite order away from the diagonals.
The centered functional $\sigma^{\mathrm{ext}}$, obtained
by the moment--cumulant formula with one-point functions
zero, is reflection positive at every coordinate plane,
with the involution $M\mapsto M^*$. It defines a separable reflection Hilbert space
$\mathcal H^{\mathrm{ext}}$. Time translation gives a
strongly continuous self-adjoint contraction semigroup
$e^{-tH}$, $H\ge0$; spatial translations give a strongly
continuous unitary representation with generators
$P_1,P_2$ commuting strongly with $H$.
\end{proposition}

\begin{proof}
Fix finitely many pairwise disjoint closed supports,
each in a bounded ball. The uniform bound makes the
truncated functions equicontinuous in the product of
the indicated norms. Enumerate rational unions of
closed cubes, a countable dense family of $C_c^1$
functions on these supports, and the finite label
basis. A diagonal subsequence makes every corresponding
truncated function converge. Equicontinuity extends
the convergence to every tuple, and multilinearity
extends the labels to $\mathscr S$. The original
meson functions converge along the given defining
sequence, so their limits remain $S$.

Symmetry and $G_{\mathbb C}$-invariance are exact at
every cutoff. They pass to the limit. Lattice
translations approximate every real translation;
the estimate and translation continuity in $C^1$
give translation invariance. For a positive-side
tuple, replace each field by its lattice field minus
its expectation. The reflection sends this centered
field to the centered field with label $M^*$ and
reflected conjugate test function. Every reflection
form is therefore the finite sum of products of
truncated functions specified by the centered
moment--cumulant formula. Approximate a coordinate
plane by bond planes. The supports stay a positive
distance from the planes, and the reflected test
functions converge in their insertion norms. Lattice
reflection positivity and the bound pass the form
inequality to the limit.

On a product of open sets with disjoint closures,
the bound makes the limiting truncated functional continuous in
the $C^1$ norms. The Schwartz kernel theorem gives
a distribution there. Uniqueness makes these local
distributions agree on intersections, so they define
one distribution off the diagonals.

On finite linear combinations of positive-time tuples,
define the positive semidefinite form by
\[
 \langle\mathcal X,\mathcal Y\rangle
    =\sigma^{\mathrm{ext}}(\theta_3\mathcal X^*\cdot\mathcal Y).
\]
We first check that time translation descends to its quotient.
Write $T_t\mathcal X=\tau_{te_3}\mathcal X$, $t\ge0$.
Translation invariance and $\theta_3\tau_{te_3}
=\tau_{-te_3}\theta_3$ give
\[
 \langle T_sv,T_tv\rangle=\langle v,T_{s+t}v\rangle,
 \qquad \|T_tv\|^2=\langle v,T_{2t}v\rangle.
\]
For a fixed $v$, these correlations are bounded uniformly
for $s,t\ge0$. Indeed, each is a finite sum of products of
truncated correlations; the support radii stay fixed, the
separations within each tuple stay fixed, and the reflected
tuple stays a positive distance from the translated tuple.

Fix $t>0$ and put $a_k=\|T_{kt}v\|^2$. Cauchy--Schwarz
for the semidefinite form gives
\[
 a_k^2=|\langle T_{(k-1)t}v,T_{(k+1)t}v\rangle|^2
                 \le a_{k-1}a_{k+1}\qquad(k\ge1).
\]
If $a_0=0$, this implies $a_1=0$. If $a_0>0$ and
$a_1>a_0$, all subsequent terms are positive and the
ratios $a_{k+1}/a_k$ are nondecreasing. They are then
bounded below by $a_1/a_0>1$, contrary to boundedness.
Thus $\|T_tv\|\le\|v\|$. In particular null vectors
remain null, so $T_t$ defines a contraction on the quotient
and its completion $\mathcal H^{\mathrm{ext}}$.
The empty tuple gives the unit vector $\Omega$.

Translation is symmetric for the form, so each bounded
operator $T_t$ is self-adjoint. The semigroup law is the
translation law on tuples. In particular
\[
 \langle v,T_tv\rangle=\|T_{t/2}v\|^2\ge0.
\]
This is also the reflection form at the intervening plane.
For a vector represented by finitely many tuples,
\[
 \|T_tv-v\|^2
 =\langle v,T_{2t}v\rangle
       -2\operatorname{Re}\langle v,T_tv\rangle+\|v\|^2
       \longrightarrow0
       \quad(t\downarrow0),
\]
by continuity in the translated test functions. Contraction
extends strong continuity to the completion. The spectral
theorem for a strongly continuous self-adjoint contraction
semigroup now gives $T_t=e^{-tH}$ with $H\ge0$.

Spatial translations preserve the form and have
inverses. They extend to strongly continuous unitary
operators, and their commutation with the semigroup
on tuples gives strong commutation of the generators.
For separability, use the positive-time members of the
countable family used for extraction, and take finite
linear combinations with rational real and imaginary
coefficients. Any positive-time tuple can be approximated
by such tuples with supports in fixed closed sets that
remain separated and stay away from the reflection plane.
Its squared reflection-norm error is a finite combination
of moments, each a polynomial in continuous truncated
functions. The error therefore tends to zero, proving
density of this countable set.
\end{proof}

\begin{proposition}[The internal group on the enlarged space]
\label{C-401-46}
The map $[X]\mapsto[X]$ embeds the mesonic space
$\mathcal H(S)$ isometrically into $\mathcal H^{\mathrm{ext}}$
and intertwines time and spatial translations. Moreover
\[
 V(S)[\mathcal X]=[S\mathcal X S^{-1}],\qquad S\in G,
\]
is a strongly continuous unitary representation. It fixes
$\Omega$, commutes with $H,P_1,P_2$, and is trivial on
$(-1,-1)$. Let $\Sigma_3$ be $1$ on the $\psi$ components
and $-1$ on the $\chi$ components. The generator $Q$ defined by
\[
 V(e^{i\varphi\Sigma_3},e^{i\varphi\Sigma_3})=e^{-i\varphi Q}
\]
has spectrum in $2\mathbb Z$. The mesonic space lies in
$\ker Q$, and both eigenspaces $Q=2$ and $Q=-2$ are
nonzero. Charge conjugation $V(J,J)$ exchanges them.
\end{proposition}

\begin{proof}
The two centered functionals agree on meson tuples,
so the displayed embedding preserves inner products.
Both translations act on the functions and are
therefore intertwined. For $S\in G$,
$(SMS^{-1})^*=SM^*S^{-1}$. Invariance of
$\sigma^{\mathrm{ext}}$ proves that $V(S)$ preserves
the form and its null space. Its inverse is $V(S^{-1})$.
For a fixed tuple the orbit lies in a finite-dimensional
span of label tuples and depends continuously on $S$;
this proves strong continuity. The stated commutations
hold on tuples and hence on the completion.

The central element $(-1,-1)$ is the identity on all
quadratic labels. Setting $\varphi=\pi$ therefore
gives $e^{-i\pi Q}=1$, so the spectral theorem places
the spectrum of $Q$ in $2\mathbb Z$. Every meson
is neutral, and the mesonic space lies in its kernel.

To see that the charge-two space is nonzero, choose
the meson $N_f^{-1/2}\sum_f\bar\psi_{+,f}\psi_{-,f}$.
Its matrix is $\tau_+\otimes t_1$,
$\tau_+=(\tau_1+i\tau_2)/2$, a linear combination of
$\Gamma_{1,1}$ and $\Gamma_{2,1}$.
The leading two-point formula of
\cite[Prop. 11.1]{QCDM1}
has a positive definite leading coefficient on this
two-dimensional span. At a pair of reflected small
positive test functions its relative error tends to
zero as their scale tends to zero. Thus the associated
one-field vector has positive norm. The element of
$G$ equal to $J\otimes1$ on the plus half
and to the identity on the minus half sends this
meson to a nonzero multiple of
$\sum_f\psi_{+,f}^{\mathsf T}C_0^{-1}\psi_{-,f}$.
Its image under $V(S)$ has the same positive norm
and charge two. Finally $J\Sigma_3J^{-1}=-\Sigma_3$,
so $V(J,J)QV(J,J)^{-1}=-Q$. This gives the nonzero
charge-minus-two space as well.
\end{proof}

\section{Energy in the charged sectors}
\label{sec:qcd2-charged-energy}

The estimate needed in the continuum has constant one.
A direct propagator estimate initially gives a constant
depending on the cutoff. Reflection positivity removes
that constant before the cutoff is taken to zero.

\subsection{A sharp decay rate on the chamber lattice}

In this subsection lattice spacing is one. Let
\[
 L_\rho=\frac12\sum_{s=1}^3t_s(1-\widehat R_s)
                      \quad\hbox{on }\ell^2(W_3),
 \qquad E=\{g:ge_1=\pm e_3\},
\]
where $\widehat R_s$ is right multiplication by the
within-site chamber reflection $w_s$. The kernel of
$L_\rho$ consists of the constants; let $\lambda_\rho>0$
be its first positive eigenvalue. These are the
within-site operator and its gap used in
\cite[Section 5]{QCDM1}.
There are $16$ elements in $E$ and $48$ in $W_3$.
Define
\begin{equation}\label{C-400-3}
 \begin{split}
 h(k)&=-\lambda_{\min}(L_\rho-k1_E),\qquad
 k(\alpha)=\tfrac32(\cosh\alpha-1)+3\sinh\alpha,\\
 e(\alpha)&=h(k(\alpha)),\quad
 0\le\alpha\le\operatorname{artanh}(1/2),\qquad
 \mu(m_0)=\sup\{\alpha:e(\alpha)<m_0\}.
 \end{split}
\end{equation}

\begin{lemma}[Directional decay]\label{C-400-11}\label{C-400-12}
If $\varphi(v)=\phi((x_v)_3)$ is real and bounded, with increments of modulus
at most $\alpha\le\operatorname{artanh}(1/2)$, then
\[
 \operatorname{Re}(e^\varphi D_Ue^{-\varphi})
                                      \ge m_0-e(\alpha).
\]
If two sets of sites $A,B$ are at separation at least
$d$ in their third coordinates, then, for
$0\le\alpha<\mu(m_0)$,
\[
 \|1_AD_U^{-1}1_B\|
       \le(m_0-e(\alpha))^{-1}e^{-\alpha d}.
\]
The estimates are uniform in $U$. Moreover
$\mu(x)/x\to1$ as $x\downarrow0$, and consequently
$a^{-1}\mu(a\mathsf m)\to\mathsf m$ uniformly for
$\mathsf m$ in compact subsets of $(0,\infty)$.
\end{lemma}

\begin{proof}
Write $R_sf(v)=U(v,v\cdot s)f(v\cdot s)$. This is a
self-adjoint unitary involution. Conjugation changes
the coefficient of the $s$th hop to
$t_s(\tfrac12-\gamma(u_s(v)))e^{-w_s(v)}$, where
$w_s(v)=\varphi(v\cdot s)-\varphi(v)$ and
$w_s(v\cdot s)=-w_s(v)$. Its Hermitian part is
$B_sR_s$, with
\[
 B_s(v)=t_s\bigl(\tfrac12\cosh w_s(v)
                         +\gamma(u_s(v))\sinh w_s(v)\bigr).
\]
The matrix $B_s$ commutes with $R_s$. For $s=1,2,3$,
$w_s=0$, since the hop stays at the same site.
For $s=0$, it vanishes outside $E$. Its two eigenvalues
are $3(\frac12\cosh w\pm\sinh w)$, nonnegative in
the stated interval. Hence
\[
 \operatorname{Re}(e^\varphi D_Ue^{-\varphi})
 =m_0+L_\rho+(\tfrac32-B_0)+B_0(1-R_0)
 \ge m_0+L_\rho-k(\alpha)1_E.
\]
This proves the first estimate.

Choose $\phi(s)=\alpha\min\{\operatorname{dist}(s,B_3),d\}$,
where $B_3$ is the set of third coordinates of $B$.
The weight is zero on $B$ and $\alpha d$ on $A$.
The first estimate makes the conjugated inverse
bounded by $(m_0-e(\alpha))^{-1}$, proving decay.

Let $u$ be the constant unit vector in $\ell^2(W_3)$.
There are $16$ moving chambers among the $48$ chambers,
so $\langle u,1_Eu\rangle=16/48=1/3$. Moreover
\[
 \|(1-|u\rangle\langle u|)1_Eu\|^2
       =\langle u,1_Eu\rangle-|\langle u,1_Eu\rangle|^2
       =\tfrac13-\tfrac19=\tfrac29.
\]
Writing a vector as $cu+v$, $v\perp u$, and using
$L_\rho\ge\lambda_\rho$ on $u^\perp$, gives, for
$0\le k<3\lambda_\rho/2$,
\[
 \frac{k}{3}\le h(k)\le
        \frac{k}{3}+\frac{2k^2}{9\lambda_\rho-6k}.
\]
Indeed the lower bound uses the trial vector $u$.
For the upper bound the quadratic form is bounded
below by $-(k/3+\delta)\|(c,v)\|^2$ when
\[
 \begin{pmatrix}
 \delta&-\sqrt2k/3\\
 -\sqrt2k/3&\lambda_\rho-2k/3+\delta
 \end{pmatrix}\ge0;
 \qquad\delta=\frac{2k^2}{9\lambda_\rho-6k}
\]
satisfies this condition. Since $k(\alpha)=3\alpha+O(\alpha^2)$,
we have $e(\alpha)=\alpha+O(\alpha^2)$.
Inversion at zero gives $\mu(x)=x+O(x^2)$ and the
claimed uniform limit.
\end{proof}

\begin{lemma}[Reflection positivity fixes the prefactor]
\label{C-400-17}
Let $T$ be symmetric on a vector space with a positive
semidefinite Hermitian form. If
$\langle v,T^{2k}v\rangle\le C\lambda^{2k}$ for all
$k\ge0$, where $0<\lambda\le1$, then
\[
 \langle v,T^{2k}v\rangle\le\lambda^{2k}\langle v,v\rangle,
 \qquad |\langle v,T^kv\rangle|
                   \le\lambda^k\langle v,v\rangle.
\]
\end{lemma}
\begin{proof}
Put $a_k=\langle T^kv,T^kv\rangle$. Cauchy--Schwarz
gives $a_k^2\le a_{k-1}a_{k+1}$. If one $a_k$ with
$k\ge1$ vanishes, this inequality forces all such
$a_k$ to vanish. Otherwise the ratios $a_{k+1}/a_k$
are nondecreasing. A ratio larger than $\lambda^2$
would contradict $a_k\le C\lambda^{2k}$ at large
$k$. Thus $a_k\le\lambda^{2k}a_0$. A final
Cauchy--Schwarz inequality gives the second assertion.
\end{proof}

\begin{proposition}[The lattice charged estimate]\label{C-400-6}
Fix $m_0>0$ and an infinite-volume torus-limit state
$\omega_\nu$. Let $\Theta$ be reflection in $x_3=1/2$
and let $\tau$ be translation by $e_3$. For an even
local element $F$ supported in $x_3\ge1$, form its
reflection class $[F]$. Translation induces a positive
self-adjoint contraction $T$ on the resulting Hilbert
space.

Let $\mathcal F=\{+,-\}\times\{1,\ldots,N_f\}$.
The weight of a monomial is
$n_f=\#\psi_f-\#\bar\psi_f$ and
$|n|_1=\sum_f|n_f|$. If $F$ has weight $n$, then
\begin{equation}\label{eq:qcd2-lattice-charged}
 0\le\omega_\nu(\tau^kF\,\Theta F)
       \le e^{-\mu(m_0)|n|_1k}\,
                           \omega_\nu(F\,\Theta F),\qquad k\ge0.
\end{equation}
If the conditional Gaussian expectation
$h_F(U)=\langle F\rangle_U$ vanishes at every $U$,
the same estimate holds with $2$ in place of $|n|_1$.
\end{proposition}

\begin{proof}
The fixed-spacing states of \cite{QCDM1} are translation
invariant and reflection positive at coordinate planes
through sites and through links. Thus the form
$\langle[G],[F]\rangle=\omega_\nu(F\Theta G)$ is
positive, and $\Theta\tau=\tau^{-1}\Theta$ makes
translation symmetric. Wick's determinant formula
and $\|D_U^{-1}\|\le m_0^{-1}$ bound
$|\omega_\nu(\tau^kF\Theta F)|$ uniformly in $k$,
with a constant depending on $F,m_0$. Apply
\cref{C-400-17} with $\lambda=1$ to obtain a
contraction on the quotient. Positivity follows
from reflection through $x_3=1$, since
$\Theta F=\Theta_1(\tau F)$ and
$\omega_\nu(\tau F\Theta_1\tau F)\ge0$.

For clarity, consider a pair of monomials of total
degree $2p$. Wick's formula is a sum of at most $p!$
products of $p$ propagators. A propagator pairs only
a quark with an antiquark of the same species.
If the monomial on the positive side has weight $n$,
let $c_f^+$ and $c_f^-$ count its quarks and antiquarks
of species $f$ that are paired across the separation.
Each pairing confined to this side removes one quark
and one antiquark of the same species. Therefore
\[
 c_f^+-c_f^-=n_f,\qquad c_f^++c_f^-\ge|n_f|.
\]
Summing over species shows that at least $|n|_1$
propagators cross. Between the supports of
$\Theta F$ and $\tau^{2k}F$, their separation in
the third coordinate is at least $2k+1$.
By \cref{C-400-11}, for every $\alpha<\mu(m_0)$,
\[
 \omega_\nu(\tau^{2k}F\Theta F)
       \le C(F,\alpha)e^{-\alpha|n|_1(2k+1)}.
\]
The finite sum over the monomials only changes
$C(F,\alpha)$. Applying \cref{C-400-17} removes
this constant. Let $\alpha\uparrow\mu(m_0)$ to
obtain \eqref{eq:qcd2-lattice-charged}. This removal is
performed at each fixed cutoff. The preliminary constant
may depend on the cutoff; the final estimate depends
on $F$ only through its squared reflection norm, which
can subsequently pass to the continuum limit.

For the last assertion subtract the product of the
two conditional Gaussian expectations in Wick's
formula. The permutations preserving both sets
of factors cancel. Every remaining permutation
crosses the separation at least twice: if each
monomial is neutral, one crossing in each direction
is required; otherwise its even nonzero charge
already requires at least two crossings. Since
$h_F=0$, the subtracted product vanishes. This gives
$C(F,\alpha)e^{-2\alpha(2k+1)}$ and the same
argument proves the asserted bound.
\end{proof}

\subsection{Passing the estimate to the enlarged limit}

Let $\mathcal H_n^{\mathrm{ext}}$ be the closed span
of positive-time tuples of total weight $n$. A
quadratic label has quark number $0$ or $\pm2$,
and every label is a sum of weight vectors.
For an irreducible representation $\pi$ of $G$,
write $|\pi|$ for the largest $\ell^1$ norm of
one of its maximal-torus weights. This is the
$\ell^1$ norm of its highest weight: the weights
lie in the convex hull of its signed-permutation
Weyl orbit.

The three estimates below distinguish three notions of
charge. The weight $n=(n_f)_{f\in\mathcal F}$ records,
for each species, the number of $\psi_f$ factors minus
the number of $\bar\psi_f$ factors. On its weight space,
the maximal torus acts by $V(z)v=z^{-n}v$. Taking every
phase equal gives the quark number
\[
 q=\sum_{f\in\mathcal F}n_f,\qquad
 \{Q=q\}=\mathop{\bigoplus}_{\sum_fn_f=q}
                              \mathcal H_n^{\mathrm{ext}}.
\]
Thus quark number retains only the sum of the species
charges. A meson $\bar\psi_{+,f}\psi_{-,f}$ has quark
number zero but weight of $\ell^1$ norm two. The weight
estimate below therefore gives a positive lower bound
on such weight spaces within the quark-number-zero sector.

The $\pi$-isotypic space is the closed span of all
copies of the irreducible representation $\pi$.
It contains the whole multiplet, and may include
vectors of weight zero. Since $H$ commutes with $G$,
a spectral subspace of $H$ in this space contains
whole copies of $\pi$. Applying the weight estimate
to a highest-weight vector therefore bounds the
energy throughout the multiplet. This is stronger
than applying the weight estimate separately to its
zero-weight vectors.


\begin{theorem}[Energy and internal charges]\label{C-400-34}
Let $S^{\mathrm{ext}}$ be an extended volume-first
limit and let $\mathsf m$ be a limit point of the
rescaled mass parameters
$\mathsf m_j=a_j^{-1}m_{0,j}\in[m_-,m_+]$
along its defining sequence.
Then
\[
 H|_{\mathcal H_n^{\mathrm{ext}}}\ge|n|_1\mathsf m,
 \qquad
 H|_{\mathcal H_\pi^{\mathrm{ext}}}\ge|\pi|\mathsf m,
 \qquad
 H|_{\{Q=q\}}\ge|q|\mathsf m.
\]
Only even $|n|_1$ and even quark numbers occur.
Every nontrivial representation that occurs has
$|\pi|\ge2$. In particular,
\[
 H\ge2\mathsf m\quad\hbox{on }\bigl((\mathcal H^{\mathrm{ext}})^G\bigr)^\perp,
 \qquad \ker H\subset(\mathcal H^{\mathrm{ext}})^G.
\]
The spectral projection onto $\ker H$ vanishes on
every nonzero-weight space. The charge-two and
charge-minus-two spaces are nonzero.
\end{theorem}

The gap $2\mathsf m$ here is on the complement of the
$G$-invariant vectors. In the invariant sector these
estimates give only $H\ge0$. Since the vacuum vector
is invariant and has energy zero, the conclusion is
\[
 \mathbb C\Omega\subset\ker H
                     \subset(\mathcal H^{\mathrm{ext}})^G.
\]
Neither uniqueness of the vacuum nor a positive gap
above it within the invariant sector follows from
this theorem.


\begin{proof}
Pass to a subsequence with $\mathsf m_j\to\mathsf m$.
For a finite combination of positive-time tuples
$\mathcal X=\sum_i\zeta_i\mathcal X_i$, with
$\mathcal X_i=((\lambda_{ir},f_{ir}))_r$, put
\[
 F_j=\sum_i\zeta_i\prod_r
     \left(\Phi_j[\lambda_{ir}](f_{ir})
                 -\omega_{\nu_j}(\Phi_j[\lambda_{ir}](f_{ir}))\right).
\]
Here $\Phi_j[0](f)$ means the action-density insertion.
The moment--cumulant formula expresses the correlations
of these centered polynomials using only blocks of at
least two insertions. Thus only the truncated functions
already known to converge occur; convergence of the
individual one-point functions is unnecessary.
If $\mathcal X$ has weight $n$, then $F_j$ has
weight $n$: the mean of a nonzero-weight field
vanishes by the torus symmetry, and subtraction
does not change a zero weight.

For large $j$ the support lies above the bond plane
$x_3=a_j/2$. Fix $t$ in $\bigcup_kL^{-k}\mathbb N$.
Then $t/a_j$ is an integer for all large $j$, and
\cref{C-400-6} gives
\[
 \omega_{\nu_j}(\tau^{t/a_j}F_j\Theta_{a_j/2}F_j)
 \le e^{-|n|_1t\mu(a_j\mathsf m_j)/a_j}
                 \omega_{\nu_j}(F_j\Theta_{a_j/2}F_j).
\]
Reflection in $a_j/2$ differs from reflection in zero
by a translation of size $a_j$. For $C_c^1$ functions
that translation tends to the identity in both
$\|\cdot\|_\infty$ and $\|\cdot\|_\sharp$;
in the latter case use uniform continuity of the
gradient. The supports remain separated. Hence
\eqref{eq:qcd2-extended-bound}, multilinearity and
the moment--cumulant formula let both sides pass
to the limit. The limiting rate is $\mathsf m$
by \cref{C-400-11}, and
\[
 \langle[\mathcal X],e^{-tH}[\mathcal X]\rangle
       \le e^{-|n|_1\mathsf mt}\|[\mathcal X]\|^2.
\]
Strong continuity extends this from the dense set
of times to all $t\ge0$.

The maximal-torus projections
$\int z^nV(z)\,dz$ have ranges
$\mathcal H_n^{\mathrm{ext}}$; they are orthogonal,
sum to the identity and commute with $H$.
The Laplace transform inequality forces the
spectral measure of $H$ on each such space to
be supported in $[|n|_1\mathsf m,\infty)$.
Indeed positive mass on $[0,c]$ with
$c<|n|_1\mathsf m$ would contradict the inequality
as $t\to\infty$.

Every spectral projection of $H$ commutes with $G$.
If the isotypic subspace of $\pi$ had a nonzero
spectral part below $|\pi|\mathsf m$, that part
would contain a copy of $\pi$, hence a vector of
its highest weight. This contradicts the weight
estimate. An even monomial has
$|n|_1\equiv\sum_fn_f\equiv0\pmod2$.
The highest weight of any occurring representation
is an occurring weight, so its norm is even and
is at least two unless the representation is
trivial. Compact-group decomposition proves the
estimate on the complement of the invariant
vectors. The charge estimate follows by summing
weight spaces and using $|n|_1\ge|\sum_fn_f|$.
The conclusions about the kernel follow from
$\mathsf m\ge m_->0$, and the last assertion
is \cref{C-401-46}.
\end{proof}

\section{Decay of parity-odd meson correlations}
\label{sec:qcd2-parity}

The parity argument needs control at coinciding
supports only for a conditional expectation over
the fermions. The loop-variance estimate of
\cref{C-415-17} supplies that control. The part
with a quark propagator joining the two supports
is estimated at separated supports.

Write $h_F(U)=\langle F\rangle_U$ and
$z(f)=a\sum_xf(x)\operatorname{tr}D_{U,+}^{-1}(x,x)$,
where this trace is over spin, colour and the chambers
over $x$, for one flavour. We use the consequence
\begin{equation}\label{eq:qcd2-loop-variance-input}
 \operatorname{Var}_\nu z(f)
 \le\frac{\|f\|_\infty^2}{4\varrho_J^2}
       \min\left\{C(R)^{1/N_f},
                         \frac{C(R)-1}{N_f^2}\right\}
\end{equation}
of \cref{C-415-17}, for real bounded $f$ supported
in a radius-$R$ ball. Its hypotheses and the
corresponding restriction on the coupling are
assumed throughout this section.

\subsection{Decay of the part joining the supports}

Let $c_h=3(1+r)=9/2$ and $G_U^+=D_{U,+}^{-1}$.
For nonnegative functions define
\[
 \Pi_U(f,h)=a^2\sum_{x,y}f(x)h(y)
                                    \|G_U^+(x,y)\|_{\mathrm{HS}}^2.
\]

\begin{lemma}[A uniform estimate for one quark loop]
\label{C-401-42}
Fix $R,d>0$. There is a number $C'(R,d)$, independent
of the cutoff and the state, such that, when
$a\le d/4$ and $a\mathsf m\le2c_h$, the following
holds. Suppose $f,h\in C_c^1$ have supports in
$B(x_0,R),B(y_0,R)$ and $|x_0-y_0|\ge2R+3d$.
For any two meson labels $\Gamma,\Gamma'$,
the mean of their conditional connected function
is bounded in modulus by
\begin{equation}\label{eq:qcd2-quark-loop-decay}
 N_fC'(R,d)\|\Gamma\|\|\Gamma'\|
 \|f\|_\infty\|h\|_\infty
       e^{-m_-(|x_0-y_0|-2R-3d)/(2c_h)}.
\end{equation}
\end{lemma}

\begin{proof}
First, for the inverse between arbitrary site sets
$A,B$, conjugating by a bounded weight with nearest-site
increments at most $\alpha$ changes the hopping operator
by norm at most $c_h(e^\alpha-1)$. Choosing
$\alpha=\log(1+m_0/(2c_h))$ gives
\[
 \|1_AG_U^+1_B\|\le\frac2{m_0}
                   e^{-\alpha\operatorname{dist}_{\mathrm{graph}}(A,B)}.
\]
For $m_0\le2c_h$, $\alpha\ge m_0/(4c_h)$.

Let $d(x)=\operatorname{dist}(x,B(x_0,R))$ and
$\eta(x)=\min\{1,\max\{0,2-d(x)/d\}\}$.
The commutator $[D_{U,+},\eta]$ is supported in
\[
 A_d=\{x:d-a\le d(x)\le2d+a\},
 \qquad\|[D_{U,+},\eta]\|\le c_ha/d.
\]
For $f$ supported where $d(x)\le d$ and a set
$B\subset\{d(x)\ge2d\}$ at Euclidean distance
$D_0$ from $A_d$, the resolvent commutator identity
gives
\[
 f^{1/2}G_U^+1_B
    =f^{1/2}G_U^+[D_{U,+},\eta]G_U^+1_B.
\]
Take a Hilbert--Schmidt norm in the first factor and
an operator norm in the last factor, and square.
The preceding decay estimate yields
\begin{equation}\label{eq:qcd2-annulus}
 \Pi_U(f,1_B)
 \le\left(\frac{2c_h}{d\mathsf m}\right)^2
       e^{-\mathsf mD_0/(2c_h)}\Pi_U(f,1_{A_d}).
\end{equation}

Choose fixed smooth nonnegative functions $\varphi,g$
such that $\varphi=1$ on $B(x_0,R)$ and vanishes
where $d(x)\ge d/4$, while $g\ge1$ where
$3d/4\le d(x)\le9d/4$ and vanishes outside
$d/2<d(x)<5d/2$. Their supports have distance at
least $d/4$. For
$\Gamma_*=i\tau_1\otimes N_f^{-1/2}1$, Wick's
formula and $G_U^-= (G_U^+)^\dagger$ give
\[
 S_2^T[\nu]((\Gamma_*,\varphi),(\Gamma_*,g))
                         =2\int\Pi_U(\varphi,g)\,d\nu.
\]
The conditional one-point functions vanish for this
label. The separated two-point bound of \cite{QCDM1}
therefore bounds the nonnegative integral on the
right, uniformly in the cutoff and the state.

For two general labels the conditional connected
function is
\[
 -a^2\sum_{x,y}f(x)h(y)
                \operatorname{tr}[\Gamma G_U(x,y)\Gamma'G_U(y,x)].
\]
Decomposing into the two $\tau_3$ blocks and using
Cauchy--Schwarz for their Hilbert--Schmidt norms
bounds this by
$2N_f\|\Gamma\|\|\Gamma'\|\Pi_U(|f|,|h|)$.
The transposed charge-conjugation identity gives
\[
 \|G_U^+(x,y)\|_{\mathrm{HS}}
 =\|G_U^+(y,x)\|_{\mathrm{HS}}.
\]
Apply \eqref{eq:qcd2-annulus}; the separation
from $\operatorname{supp}h$ to $A_d$ is at least
$|x_0-y_0|-2R-3d$. The displayed separated meson
bound then proves \eqref{eq:qcd2-quark-loop-decay},
with, for example,
\[
 C'(R,d)=\left(\frac{2c_h}{dm_-}\right)^2
                   C_2(d/4;R+3d)\|g\|_\infty,
\]
where $C_2$ is the companion's separated two-point
constant.
\end{proof}

\subsection{The parity tail}

Fix a coordinate direction $e_\mu$. Let $c$ and $w>0$
belong respectively to
$\bigcup_k L^{-k}\frac12\mathbb Z$ and
$\bigcup_k L^{-k}\mathbb Z$. Let $\theta_c$ be
reflection in $x_\mu=c$, and $\theta_0$ reflection
in the middle plane $x_\mu=c+w/2$. For
\[
 \mathcal F=\sum_{k=1}^p c_k(\Gamma_k,f_k),\qquad
 \|\mathcal F\|=\sum_k|c_k|\|\Gamma_k\|\|f_k\|_\infty,
\]
assume $\Gamma_k\in\mathfrak B$ and that the real functions
$f_k\in C_c^1$ are supported
in the open slab $c<x_\mu<c+w$ and in a ball
$B(x_*,R)$ centered on its middle plane. The parity
signs $\pi_\Gamma$ are those of \cite{QCDM1}. Put
\[
 \mathcal F^-=\frac12\sum_kc_k
       \bigl((\Gamma_k,f_k)-\pi_{\Gamma_k}
                                    (\Gamma_k,f_k\circ\theta_0)\bigr),
 \qquad s_{\mathcal F}(t)
   =\sigma_S(\theta_c(\mathcal F^-)^*\cdot
                               \tau_{te_\mu}\mathcal F^-).
\]

\begin{proposition}[Decay of the parity-odd part]\label{C-401-45}
Let $S$ be a volume-first limit satisfying the
hypotheses of \cref{C-415-17}. Then $s_{\mathcal F}$
is continuous, nonnegative and nonincreasing on
$(0,\infty)$. For each cutoff of its defining
sequence let $F_j^-$ be the polynomial associated
to $\mathcal F^-$ and set
\[
 C_{\mathcal F}=\limsup_j\int|h_{F_j^-}|^2\,d\nu_j.
\]
This number is finite, with
\[
 C_{\mathcal F}\le\varrho_J^{-2}\|\mathcal F\|^2
                  \min\{N_f^2C(R)^{1/N_f},C(R)-1\}.
\]
For every $d,t_0>0$ with $t_0+w\ge2R+3d$,
\begin{equation}\label{eq:qcd2-parity-tail}
 \int_{t_0}^\infty s_{\mathcal F}(t)\,dt
 \le(t_0+w)C_{\mathcal F}
   +N_fC'(R,d)\|\mathcal F\|^2\frac{2c_h}{m_-}
                e^{-m_-(t_0+w-2R-3d)/(2c_h)}.
\end{equation}
Consequently $s_{\mathcal F}(t)\to0$ as $t\to\infty$.
\end{proposition}

\begin{proof}
At all sufficiently fine cutoffs the two planes
are lattice planes and $m=w/a$ is an integer.
Let $V_{\theta_0}$ be the linear parity transformation
on the fermions, using the matrix
$\sigma(e_\mu)\otimes\tau_2$. Then
$V_{\theta_0}\Phi_\Gamma(f)
=\pi_\Gamma\Phi_\Gamma(f\circ\theta_0)$.
Thus $V_{\theta_0}F_j^-=-F_j^-$ and its expectation
vanishes. Put $A_j=h_{F_j^-}$ and
$\Xi_j=[F_j^-]$ in the reflection space at $x_\mu=c$.

For a meson label $\Gamma=1_{\mathrm s}\otimes
X\otimes1_{\mathrm c}$, Wick's rule gives
\[
 h_{\Phi_\Gamma(f)}=-\alpha_\Gamma z(f)
                              -\beta_\Gamma\overline{z(f)},
 \quad\alpha_\Gamma=\operatorname{tr}(e_+X),
 \quad\beta_\Gamma=\operatorname{tr}(e_-X),
 \quad |\alpha_\Gamma|,|\beta_\Gamma|\le N_f\|\Gamma\|.
\]
Since $A_j$ has mean zero, subtracting the means
of these loops does not change their combination.
All functions involved are supported in the
same ball, which is preserved by $\theta_0$.
Minkowski's inequality and
\eqref{eq:qcd2-loop-variance-input} give
\[
 \|A_j\|_2\le
  \varrho_J^{-1}\|\mathcal F\|
          \min\{N_fC(R)^{1/(2N_f)},(C(R)-1)^{1/2}\}.
\]
This proves the stated uniform estimate.

We next retain the sign coming from parity.
Define the antilinear conjugation on generators by
$\psi^\sharp=(\bar\psi\tau_2)^{\mathsf T}$,
$\bar\psi^\sharp=(\tau_2\psi)^{\mathsf T}$,
and reverse products. The identities on the doubled
field give
\[
 \Theta_0=\alpha_{(1,-1)}V_{\theta_0}\sharp,
 \qquad h_{F^\sharp}=\overline{h_F},
 \qquad h_{\alpha_SF}=h_F.
\]
As $\theta_c=\tau^{-m}\theta_0$ in lattice units,
putting $B_j=\tau^{-m}\alpha_{(1,-1)}((F_j^-)^\sharp)$
gives $\Theta_cF_j^-=-B_j$ and
$h_{B_j}=\tau^{-m}\overline{A_j}$. Hence
\begin{equation}\label{C-401-32}
 s_j(n):=\langle\Xi_j,T^n\Xi_j\rangle
   =-\langle A_j,\tau^{n+m}A_j\rangle_{L^2(\nu_j)}-R_j(n),
\end{equation}
where $R_j(n)$ is the averaged conditional connected
function of $\tau^nF_j^-$ and $B_j$.
The two supports are in radius-$R$ balls whose
centers have distance $na+w$. Applying
\cref{C-401-42} to their finitely many meson terms
gives, for $na+w\ge2R+3d$,
\[
 |R_j(n)|\le N_fC'(R,d)\|\mathcal F\|^2
                   e^{-m_-(na+w-2R-3d)/(2c_h)}.
\]
Each of the four reflected/unreflected pairings
has coefficient $1/4$, so no extra factor is needed.

The sequence $l\mapsto\langle A_j,\tau^lA_j\rangle$
is of positive type, since translation is unitary
on $L^2(\nu_j)$. The relevant finite inequality is
\[
 0\le\frac1N\left\|\sum_{l=0}^{N-1}\tau^lA_j\right\|_2^2
 =\|A_j\|_2^2+
       2\sum_{l=1}^{N-1}(1-l/N)
                     \operatorname{Re}\langle A_j,\tau^lA_j\rangle.
\]
Let $n_0$ be the least integer with $n_0a\ge t_0$ and
take $N>m+n_0$. For $l<m+n_0$ bound the real part by
$\|A_j\|_2^2$. For $l\ge m+n_0$, the minus sign in
\eqref{C-401-32} gives
\[
 \operatorname{Re}\langle A_j,\tau^lA_j\rangle
       \le-s_j(l-m)+|R_j(l-m)|.
\]
Consequently
\begin{align*}
 2\sum_{n=n_0}^{N-m-1}\left(1-\frac{n+m}{N}\right)s_j(n)
 &\le \left[1+2\sum_{l=1}^{m+n_0-1}
                   \left(1-\frac lN\right)\right]\|A_j\|_2^2\\
 &\quad+2\sum_{n=n_0}^{N-m-1}
              \left(1-\frac{n+m}{N}\right)|R_j(n)|\\
 &\le[2(m+n_0)-1]\|A_j\|_2^2
                         +2\sum_{n\ge n_0}|R_j(n)|.
\end{align*}
Because $0\le T\le1$, every $s_j(n)$ is nonnegative.
The weights on the left increase to one as $N\to\infty$,
whereas the remainder on the right is summable by its
exponential bound. Monotone convergence therefore gives
\[
 \sum_{n\ge n_0}s_j(n)
 \le(m+n_0-\tfrac12)\|A_j\|_2^2
                                      +\sum_{n\ge n_0}|R_j(n)|.
\]
Multiply by $a$. The first coefficient is at most
$w+t_0+a/2$, and the geometric-series factor in
the second term is
$a/(1-e^{-m_-a/(2c_h)})\to2c_h/m_-$.

At every time $t$ in $\bigcup_kL^{-k}\mathbb N$,
$s_j(t/a_j)\to s_{\mathcal F}(t)$ by the defining
correlation limits. The reflected supports stay
separated, so continuity of test-function translation
and the companion's two-point bound make
$s_{\mathcal F}$ continuous. Nonnegativity and
monotonicity pass from $s_j$ on the dense set of
times to all times. Fix a finite partition
$t_0<t_1<\cdots<t_p$, with $t_1,\ldots,t_p$ in that
dense set. For all sufficiently fine cutoffs, monotonicity
gives the explicit lower sum
\[
 a_j\sum_{n\ge n_0}s_j(n)
    \ge\sum_{i=1}^p(t_i-t_{i-1}-a_j)s_j(t_i/a_j).
\]
Only finitely many continuum correlation limits are needed
on the right. Pass along the defining subsequence, use
the upper bound on the left, and then take the supremum
over partitions of each fixed interval. Finally let the
upper endpoint tend to infinity. Thus the continuum
integral is obtained after the finite-cutoff tail bound;
no uniform integrable majorant for the correlations has
been assumed. The preceding estimates give
\eqref{eq:qcd2-parity-tail}. Finally
$(t-t_0)s_{\mathcal F}(t)\le\int_{t_0}^t
s_{\mathcal F}(s)\,ds$, proving decay.
\end{proof}

For a real $u$ symmetric about a coordinate plane,
the field $\Phi_{\tau_3\otimes1}(u)$ is one of
these parity-odd combinations. Choose a slightly
wider lattice-compatible slab about that plane.
Its ordinary two-point correlation with its
translate is $s_{\mathcal F}$ with a fixed shift
in the translation parameter. It therefore tends
to zero and has the corresponding integrable tail.

\section{The short-distance law for two currents}
\label{sec:qcdm2-currents}

The conserved lattice current is supported at the two ends of a
link.  The source estimates in \cite{QCDM1} concern bilinears at
one site.  We first compare these two insertions at a fixed gauge
configuration.  The error contains a covariant difference of a
fermion propagator away from its source.  An energy inequality
controls this difference by a positive pseudoscalar two-point
function.  The source estimates then compare the site insertion
with its free value, whose leading term is computed by a Pauli
matrix trace.

The construction and all data remain those of the chamber model
with $r=1/2$.  We retain $r$ in a few local identities to display
the coefficients of the hopping operator.  For a neighbouring
pair of places $v,v'=v\cdot s$, write
\[
 h(v,v')=t_s(r-\gamma(u_s(v))),\qquad
 t^T(v,v')=\bar\psi(v)h(v,v')T U(v,v')\psi(v'),
\]
where $T=e_+\otimes t_++e_-\otimes t_-$ acts on species and
flavour. Thus $T$ commutes with the gamma matrices and with colour.
Its reflection involution is
$T^*=\tau_3T^\dagger\tau_3=T^\dagger$.
The chamber graph, wall directions and transporters are those in
\cref{sec:qcdm2-setting}. In this section we label sites by
$x\in a\mathbb Z^3$: the place $(x,q)$ corresponds to
$(a^{-1}x,q)$ in \eqref{eq:qcdm2-Dirac}. Set
\begin{equation}\label{C-402-3}
 \begin{split}
 \mathsf j^T(v,v')&=t^T(v',v)-t^T(v,v'),\\
 j^T_\mu(x)&=\sum_{g:\,ge_1=e_\mu}
       \mathsf j^T((x,g),(x,g)\cdot0),\\
 J^T(f)&=ia\sum_{x,\mu}
       f_\mu(x+\tfrac12ae_\mu)j^T_\mu(x).
 \end{split}
\end{equation}
The site sum has eight terms.  The factor $i$ is part of the
definition of $J^T$; it determines the sign of its two-point
function and its reflection law.

\begin{proposition}[Lattice conservation]\label{prop:qcdm2-current-Ward}
Let $\delta_v^T$ be the derivation which fixes the links and
satisfies
\[
 \delta_v^T\psi(u)=\delta_{uv}T\psi(u),\qquad
 \delta_v^T\bar\psi(u)=-\delta_{uv}\bar\psi(u)T.
\]
At every gauge configuration, with
$(\operatorname{div}\mathsf j^T)(v)=\sum_{v'\sim v}\mathsf j^T(v,v')$,
the conditional Gaussian functional satisfies
\begin{equation}\label{eq:qcdm2-current-Ward}
 \langle(\operatorname{div}\mathsf j^T)(v)F\rangle_U
 =-\langle\delta_v^TF\rangle_U
\end{equation}
for every polynomial $F$ in finitely many fermion generators.
The identity also holds after gauge averaging. At a site,
\[
 \sum_{v:x_v=x}(\operatorname{div}\mathsf j^T)(v)
 =\sum_\mu\bigl[j^T_\mu(x)-j^T_\mu(x-ae_\mu)\bigr].
\]
Thus the divergence vanishes in correlations away from the
fermion insertions; at an insertion it gives the contact
term on the right of \eqref{eq:qcdm2-current-Ward}.
\end{proposition}

\begin{proof}
The diagonal mass term is invariant under $\delta_v^T$.
Since $T$ commutes with the hopping matrices and the links,
\[
 \delta_v^T(\bar\psi D_U\psi)
 =-(\operatorname{div}\mathsf j^T)(v).
\]
The changes $\psi(v)\mapsto e^{tT}\psi(v)$ and
$\bar\psi(v)\mapsto\bar\psi(v)e^{-tT}$ have cancelling
Berezin Jacobians. Differentiating their change-of-variables
identity proves \eqref{eq:qcdm2-current-Ward} on a torus.
On the infinite lattice the same identity follows from
$D_UD_U^{-1}=1$ and Wick's rule. Summing over the chambers
at one site cancels the internal edges and leaves the
displayed site divergence. Finally, $\delta_v^TF=0$ when
$F$ has no generator at $v$.
\end{proof}

Fix $0<\vartheta<1$.  In this section a test pair at scale $d$
means two nonzero, nonnegative Lipschitz functions $f,h$ such that
each support lies in a ball of radius $d$,
\begin{equation}\label{eq:qcdm2-current-test-pair}
 \begin{gathered}
 d\le\operatorname{dist}(\operatorname{supp}f,
                         \operatorname{supp}h)\le6d,\\
 \int f\ge\vartheta d^3\|f\|_\infty,\qquad
 \int h\ge\vartheta d^3\|h\|_\infty,\qquad
 \operatorname{Lip}f\le\frac{\|f\|_\infty}{\vartheta d},
 \quad
 \operatorname{Lip}h\le\frac{\|h\|_\infty}{\vartheta d}.
 \end{gathered}
\end{equation}
On a torus we also require $20d\le\ell$ and read the supports in
one Euclidean ball of radius $5d$.  In particular
$d\le|x-y|\le10d$ on the product of the supports.  Put
\[
 k_{\mu\nu}(x)=\frac{\delta_{\mu\nu}|x|^2-2x_\mu x_\nu}{|x|^6},
 \qquad
 \mathcal K_4(f,h)=\iint f(x)h(y)|x-y|^{-4}\,dx\,dy.
\]
The test-pair assumptions give
\begin{equation}\label{eq:qcdm2-current-Klower}
 10^{-4}\vartheta^2d^2\|f\|_\infty\|h\|_\infty
       \le\mathcal K_4(f,h)
       \le(4\pi/3)^2d^2\|f\|_\infty\|h\|_\infty.
\end{equation}

\begin{theorem}[Two currents]\label{C-417-14}
There are a constant $C_V'(g)$ and a nondecreasing function
$\epsilon_V(g,\vartheta;d)$ tending to zero as $d\downarrow0$
with the following properties.  Let
$0<d\le d_*:=\min\{1/200,1/(20{m_+})\}$ and $16a\le d$,
and let $(f,h)$ be a test pair as in
\eqref{eq:qcdm2-current-test-pair}.  The state may be a torus
functional or a functional $\omega_\nu$ with
$\nu\in\mathcal G^{\rm tor}_{K,\beta,m_0}$.
For unit vectors $e,e'$, $a_0,b_0\in\{0,3\}$ and Hermitian
flavour matrices $t,t'$, put $T=\tau_{a_0}\otimes t$,
$T'=\tau_{b_0}\otimes t'$.  Then
\begin{align}
 \kappa_2(J^T(fe),J^{T'}(he'))
  =\frac{\operatorname{tr}(tt')}{2\pi^2}
  \left[\delta_{a_0b_0}
       \iint f(x)h(y)\langle e,k(x-y)e'\rangle\,dx\,dy
                     +\varepsilon\mathcal K_4(f,h)\right],
 \label{eq:qcdm2-current-law}
\end{align}
where
\begin{equation}\label{eq:qcdm2-current-law-error}
 |\varepsilon|\le\epsilon_V(g,\vartheta;d)
       +C_V'(g)\left[
          \left(\frac{a}{\vartheta^4d}\right)^{1/2}
                         +\frac{a}{\vartheta^3d}\right].
\end{equation}
If $\operatorname{tr}(tt')=0$, the cumulant is exactly zero.

At fixed test pair, these cumulants are bounded along every
sequence with $K\to\infty$.  Every limit point satisfies
\eqref{eq:qcdm2-current-law} with
$|\varepsilon|\le\epsilon_V(g,\vartheta;d)$.
This holds on a fixed torus and for infinite-lattice states;
it also holds along tori on which the fixed pair is read with
$20d\le\ell$.
\end{theorem}

The assertion about limit points concerns the indicated two-point
functions.  It does not by itself construct all joint current
correlations.  For an orthonormal Hermitian flavour basis,
$\operatorname{tr}(t_jt_k)=\delta_{jk}$, the leading coefficient
is $\delta_{a_0b_0}\delta_{jk}/(2\pi^2)$.

\subsection{Conditional loops and their signs}

We use the resolvent in lattice units in this section:
\[
 G_U=(D_+[U]+m_0)^{-1}=a^{-1}\mathsf G_+.
\]
Let $h_+(v,v')=t_s(r-\sigma(u_s(v)))$ be the upper-species
hop.  Define $Y(v,v')$ by its two nonzero blocks
\[
 Y(v,v')[v',v]=h_+(v',v)U(v',v),\qquad
 Y(v,v')[v,v']=-h_+(v,v')U(v,v'),
\]
and sum these over the eight chamber edges above a link to obtain
$Y_\mu(x)$.  For a site vector function $F$ set
\[
 Y(F)=\sum_{x,\mu}F_\mu(x)Y_\mu(x),\qquad
 N(F)(x,g)=\sigma(F(x)).
\]
For a scalar function $f$ and a unit vector $e$ use
$F^{f,e}_\mu(x)=f(x+ae_\mu/2)e_\mu\cdot e$, and abbreviate
$Y(fe)=Y(F^{f,e})$, $N(fe)=N(F^{f,e})$.
The two loop expectations to be compared are
\begin{equation}\label{C-417-2}
 \begin{split}
 E_\nu(f,e;h,e')&=a^2\mathsf E_\nu
                \operatorname{Tr}(G_UY(fe)G_UY(he')),\\
 E_\nu^{\rm s}(f,e;h,e')&=a^2\mathsf E_\nu
                \operatorname{Tr}(G_UN(fe)G_UN(he')).
 \end{split}
\end{equation}
The trace includes places, spin and colour, and no flavour
variable.  Subscript $1$ denotes the value at $U=1$.
We also use the positive form
\begin{equation}\label{eq:qcdm2-positive-propagator}
 P_\nu(F,H)=2a^2\sum_{u,u'}F(x_u)H(x_{u'})
                 \mathsf E_\nu\|G_U(u,u')\|_{\rm HS}^2
\end{equation}
for scalar $F,H\ge0$.
For a site matrix $A$ acting trivially on colour, write
\[
 \Phi_A=a\sum_v\bar\psi(v)A(x_v)\psi(v).
\]
This is the bilinear in \eqref{eq:qcdm2-meson},
now allowing $A$ to act on spin as well.

\begin{lemma}\label{C-417-3}\label{C-402-24}
The conditional mean of every current is zero.  For real $f,h$,
\begin{equation}\label{eq:qcdm2-current-loop}
 \kappa_2(J^T(fe),J^{T'}(he'))
 =\operatorname{tr}(t_+t'_+)E_\nu(f,e;h,e')
       +\operatorname{tr}(t_-t'_-)\overline{E_\nu(f,e;h,e')}.
\end{equation}
For real site vector functions $F,H$, choose a Hermitian $t$
with $\operatorname{tr}t^2=1$ and let
$N_F=\sigma(F)\otimes e_+\otimes t$, acting as the identity
on colour.  Then
\[
 \omega_\nu(\Phi_{N_F})=0,\qquad
 \kappa_2(\Phi_{N_F},\Phi_{N_H})
             =-a^2\mathsf E_\nu\operatorname{Tr}(G_UN(F)G_UN(H)).
\]
For $\Gamma=i\tau_1\otimes t$ the corresponding scalar
bilinears satisfy
$\kappa_2(\Phi_\Gamma(F),\Phi_\Gamma(H))=P_\nu(F,H)$.
\end{lemma}

\begin{proof}
Let $C=\sigma_2\otimes\epsilon$ on spin and colour, where
$\epsilon U^{\mathsf T}\epsilon^{-1}=U^{-1}$ for
$U\in\mathrm{SU}(2)$.  Charge conjugation gives
\[
 G_U(u,u')^{\mathsf T}=C^{-1}G_U(u',u)C,
 \qquad \sigma_2\sigma_\mu^{\mathsf T}\sigma_2=-\sigma_\mu.
\]
Consequently the two orders of $G_U(u,u')$ have equal
Hilbert--Schmidt norms and
$\operatorname{tr}(\sigma_\mu G_U(v,v))=0$.
Let $\mathcal C$ be $C$ extended by the identity on species and
flavour. The change of Grassmann variables
$\psi\mapsto(\bar\psi\mathcal C)^{\mathsf T}$,
$\bar\psi\mapsto-(\mathcal C^{-1}\psi)^{\mathsf T}$ preserves
the action and its Berezin measure and sends
$\mathsf j^T$ to $-\mathsf j^{T^{\mathsf T}}$.
On the infinite lattice the same invariance follows directly
from the covariance identity and Wick's rule.

At fixed $U$ the current mean is a flavour-invariant linear
form in $T$, hence a linear combination of
$\operatorname{tr}t_+$ and $\operatorname{tr}t_-$.
It is unchanged by transposition and changed in sign by
charge conjugation, so it is zero.  Thus the covariance
of conditional means makes no contribution.  Wick's rule
gives $-\operatorname{Tr}(GYGY')$ for two bilinears, while
the two factors $ia$ in the currents give $-a^2$.
In the lower species $G_-=G_+^*$ and $Y_-=-Y_+^*$;
cyclicity of the trace gives the complex conjugate loop.
This proves \eqref{eq:qcdm2-current-loop}.

The vanishing Pauli trace proves the assertion about the
site mean; its second cumulant is the stated Wick trace.
Finally $i\tau_1$ exchanges the two species.  Its two
orientations in the loop give twice the squared
Hilbert--Schmidt norm, since $G_-=G_+^*$ and the two
orders have the same norm.  This proves the last identity.
\end{proof}

\subsection{Replacing the two endpoints by a site}

For a chamber function $\varphi$, let
$\nabla_s\varphi(v)=\varphi(v)-U(v,v\cdot s)\varphi(v\cdot s)$.
For a set of sites $S$ write
\[
 \|\varphi\|_S^2=\sum_{x_v\in S}|\varphi(v)|^2,
 \qquad
 \mathfrak D_S(\varphi)=\sum_{x_v\in S}\sum_{s=0}^3
                                      |\nabla_s\varphi(v)|^2.
\]

\begin{lemma}\label{C-417-4}
For a site vector function $F$ supported in $S$,
$Y(F)=N(F)+R(F)$, where
\begin{equation}\label{eq:qcdm2-current-vertex-error}
 |\langle\chi,R(F)\varphi\rangle|
 \le c_Y\|F\|_\infty\bigl(
       \mathfrak D_S(\chi)^{1/2}\|\varphi\|_S
       +\|\chi\|_S\mathfrak D_S(\varphi)^{1/2}\bigr),
 \qquad c_Y=\sqrt3(51+3r).
\end{equation}
\end{lemma}

\begin{proof}
For $v=(x,g)$ with $ge_1=e_\mu$, the wall-zero hops are
$3(r-\sigma_\mu)$ and $3(r+\sigma_\mu)$.  Expanding the two
endpoint values about $v$ gives
\begin{align*}
 \langle\chi,Y(v,v\cdot0)\varphi\rangle
 &=6\chi(v)^*\sigma_\mu\varphi(v)
  -3(\nabla_0\chi(v))^*(r+\sigma_\mu)\varphi(v)\\
 &\hspace{2em}+3\chi(v)^*(r-\sigma_\mu)\nabla_0\varphi(v).
\end{align*}
The last two terms sum to at most $3(1+r)$ times the
right-hand expression in \eqref{eq:qcdm2-current-vertex-error}
at $x$.  To replace the first term by the sum over all
48 chambers, choose a map from these chambers onto the
eight with $ge_1=e_\mu$, with six preimages each.
Put $b(g)=\chi(x,g)^*\sigma_\mu\varphi(x,g)$ and denote
the map by $\pi_\mu$. Then
\[
 6\sum_{g:\,ge_1=e_\mu}b(g)
 =\sum_gb(g)+\sum_g\bigl(b(\pi_\mu(g))-b(g)\bigr).
\]
The first sum is the site insertion. Join each chamber to
its image by a shortest gallery using walls $1,2,3$.
These walls connect the chambers of one site, since their
positive quadratic form has precisely the chamber-independent
fields as its kernel. Their transporters are the identity.
For one step $w\mapsto w\cdot s$ of a gallery,
\[
 b(w\cdot s)-b(w)
 =-\chi(w)^*\sigma_\mu\nabla_s\varphi(w)
   -(\nabla_s\chi(w))^*\sigma_\mu\varphi(w\cdot s).
\]
A shortest gallery visits a chamber at most once. There are
48 galleries, so each chamber and each chamber-wall pair
occurs at most 48 times. Cauchy--Schwarz therefore gives
\[
 \left|\sum_g\bigl(b(\pi_\mu(g))-b(g)\bigr)\right|
 \le48\bigl(\|\chi\|_x\mathfrak D_x(\varphi)^{1/2}
        +\mathfrak D_x(\chi)^{1/2}\|\varphi\|_x\bigr).
\]
Thus the coefficient is $51+3r$ for one direction.
Sum over directions using $\sum_\mu|F_\mu(x)|\le\sqrt3|F(x)|$
and then apply Cauchy--Schwarz over sites.
\end{proof}

\begin{lemma}[Energy away from the source]\label{C-417-5}
Let $0\le\eta\le1$ be a site cutoff with finite support, let
$\delta$ be the largest difference of $\eta$ at neighbouring
sites, and let $A_\eta$ contain the places at both ends of
every hop where $\eta$ changes.  Suppose $\eta=1$ on $S$
and its neighbours.  If either
$(D_+[U]+m_0)\varphi=0$ or
$(D_+[U]^*+m_0)\varphi=0$ wherever $\eta>0$, then
\begin{equation}\label{eq:qcdm2-Caccioppoli}
 \mathfrak D_S(\varphi)
       \le C_E\delta\|\varphi\|_{A_\eta}^2,
 \qquad C_E=2(1+r)/r.
\end{equation}
\end{lemma}

\begin{proof}
Write $D$ for either of these two operators without its
mass.  Its Hermitian part is
$\Delta=\sum_srt_s(1-R_s)$, where $R_s$ is a self-adjoint
unitary and $t_s\ge3$. For every place $v$,
\[
 ((D+m_0)(\eta\varphi))(v)
 =\eta(x_v)((D+m_0)\varphi)(v)
 +\sum_sD(v,v\cdot s)
       (\eta(x_{v\cdot s})-\eta(x_v))\varphi(v\cdot s).
\]
The first term is zero where $\eta>0$ by the equation,
and elsewhere by its factor $\eta$. Thus
$(D+m_0)(\eta\varphi)=[D,\eta]\varphi$.
Only wall zero changes the site.  Its hop has norm
$3(1+r)$, and $v\mapsto v\cdot0$ is a bijection, so
$\|[D,\eta]\varphi\|\le3(1+r)\delta
\|\varphi\|_{A_\eta}$. Let $T_\eta$ be the set of places
at which the commutator can be nonzero. It is contained
in $A_\eta$. Since $m_0>0$ and $0\le\eta\le1$,
\begin{align*}
 \langle\eta\varphi,\Delta\eta\varphi\rangle
 &\le\operatorname{Re}
       \langle\eta\varphi,(D+m_0)\eta\varphi\rangle\\
 &=\operatorname{Re}\langle\eta\varphi,[D,\eta]\varphi\rangle
 \le\|\varphi\|_{T_\eta}\|[D,\eta]\varphi\|
 \le3(1+r)\delta\|\varphi\|_{A_\eta}^2.
\end{align*}
On the other hand this energy is
$\frac12\sum_{v,s}rt_s|\nabla_s(\eta\varphi)(v)|^2$,
which is at least $(3r/2)\mathfrak D_S(\varphi)$.
\end{proof}

\begin{lemma}[Point splitting]\label{C-417-6}
Let $16a\le d$ and
$\operatorname{dist}(\operatorname{supp}f,
\operatorname{supp}h)\ge d$.  Let $B_f,B'_f$ be the sets
of sites at distance at most $a$ and $d/4+5a$, respectively,
from $\operatorname{supp}f$, and similarly for $h$.
Then
\begin{align}
 |E_\nu(f,e;h,e')-E_\nu^{\rm s}(f,e;h,e')|
 &\le C_{\rm split}(a/d)^{1/2}\|f\|_\infty\|h\|_\infty
       \bigl(P_\nu(1_{B'_f},1_{B_h})
                      +P_\nu(1_{B_f},1_{B'_h})\bigr),
 \label{eq:qcdm2-point-splitting}
\end{align}
where $C_{\rm split}=12(1+r)c_YC_E^{1/2}$.
The same estimate holds at $U=1$.
\end{lemma}

\begin{proof}
Write $Y(fe)=N(fe)+R_f$ and similarly for $h$.  The
difference of traces is
\[
 \operatorname{Tr}(G_UR_fG_UY(he'))
       +\operatorname{Tr}(G_UN(fe)G_UR_h).
\]
Let $S_f$ be the site support of $F^{f,e}$; it lies within
$a/2$ of $\operatorname{supp}f$.  Choose a cutoff equal
to one within $2a$ of $S_f$, decreasing linearly to zero
over width $d/4$.  Its neighbouring-site difference is
at most $4a/d$, its transition region lies in $B'_f$,
and it vanishes on $B_h$.

Expand $Y(he')$ in its nonzero blocks $Y'_{pq}$.
Their norms are at most $3(1+r)\|h\|_\infty$, and each
place is a row and column of at most one such block.
The matrix elements of $(G_UR_fG_U)(q,p)$ are
$\langle G_U^*(\cdot,q)e_\alpha,
R_fG_U(\cdot,p)e_\beta\rangle$.  Both columns satisfy
the homogeneous equation on the support of the cutoff.
For
$X_{q,\alpha}=G_U^*(\cdot,q)e_\alpha$ and
$\Phi_{p,\beta}=G_U(\cdot,p)e_\beta$,
\cref{C-417-5,C-417-4} give
\[
 |\langle X_{q,\alpha},R_f\Phi_{p,\beta}\rangle|
 \le2c_Y\|f\|_\infty(C_E\delta)^{1/2}
       \|X_{q,\alpha}\|_{B'_f}\|\Phi_{p,\beta}\|_{B'_f}.
\]
Set
\[
 \mathcal X_q=\sum_{u\in B'_f}\|G_U(q,u)\|_{\rm HS}^2,
 \qquad
 \mathcal Y_p=\sum_{u\in B'_f}\|G_U(u,p)\|_{\rm HS}^2,
\]
where a sum over a site set includes all its places.
Summing the squares over the four spin-colour basis
vectors gives
\[
 \|(G_UR_fG_U)(q,p)\|_{\rm HS}
 \le2c_Y\|f\|_\infty(C_E\delta)^{1/2}
                      (\mathcal X_q\mathcal Y_p)^{1/2}.
\]
On this four-dimensional fibre the trace norm is at most
twice the Hilbert--Schmidt norm. With $c_h=3(1+r)$,
\begin{align*}
 |\operatorname{tr}((G_UR_fG_U)(q,p)Y'_{pq})|
 &\le4c_hc_Y\|f\|_\infty\|h\|_\infty
          (C_E\delta)^{1/2}(\mathcal X_q\mathcal Y_p)^{1/2}\\
 &\le2c_hc_Y\|f\|_\infty\|h\|_\infty
          (C_E\delta)^{1/2}(\mathcal X_q+\mathcal Y_p).
\end{align*}
Each $p$ and $q$ occurs at most once, and
$(C_E\delta)^{1/2}\le2C_E^{1/2}(a/d)^{1/2}$.
The first trace is therefore at most
\[
 C_{\rm split}(a/d)^{1/2}\|f\|_\infty\|h\|_\infty
 \sum_{u\in B'_f,\,u'\in B_h}
       \bigl(\|G_U(u,u')\|_{\rm HS}^2
                         +\|G_U(u',u)\|_{\rm HS}^2\bigr).
\]
For the second trace interchange $f,h$; the blocks of
$N(fe)$ have norm at most $\|f\|_\infty$ and one block
per place.  Its constant is no larger.  Multiply by
$a^2$, integrate, and use the equality of the two
Hilbert--Schmidt norms from \cref{C-417-3}.
\end{proof}

\subsection{Source estimates with a spin matrix}

The site current insertion contains $\sigma_\mu$.  We therefore
need the mass source to range in
\[
 \mathcal M=\operatorname{End}
       (\mathbb C^2_{\rm spin}\otimes\mathbb C^2_\tau
                          \otimes\mathbb C^{N_f})\otimes1_{\rm colour}.
\]
A source $J(x)\in\mathcal M$ acts identically at the 48
chambers of a site. Write $\hat J$ for this multiplication
operator, $\|J\|=\sup_x\|J(x)\|$, and $\mathsf n=-aJ$.
The fine operator is $D[U]+m_0+\hat{\mathsf n}$.
There are three source domains in this argument:
\begin{align*}
 \mathscr O_k^{\mathcal M}
   &=\{J:\lambda_k\|J\|<r_M/L\},\\
 \mathcal P^{\mathcal M}
   &=\{J:\|J\|\le2\varrho_J\},\\
 \mathcal O_j^{{\rm lev},\mathcal M}
   &=\{J:\lambda_j\|J\|<R_J\},
 \qquad R_J=\min\{r_M,\varrho_S,\varrho'_S/L\}.
\end{align*}
The radii are those of \cite[Sections 10 and 11]{QCDM1}.
Here $r_M$ is chosen from the inverse and walk estimates,
and $\varrho_S,\varrho'_S$ are the free-operator radii.
The physical radius satisfies
$2\varrho_J\le c_J{m_-}$, where
$c_J\le\min\{1/2,\varrho_*,r_M/L\}$; hence
$\mathcal P^{\mathcal M}$ lies strictly inside every
$\mathscr O_k^{\mathcal M}$.
The source chain through level $k$ is formed on
$\mathscr O_k^{\mathcal M}$. To estimate the energy made
at level $j$ we use the larger level domain
$\mathcal O_j^{{\rm lev},\mathcal M}\subset
\mathscr O_{j-1}^{\mathcal M}$.
Indeed $\lambda_j=L\lambda_{j-1}$ and $R_J\le r_M$.
Only the supremum norm of $J$ enters these domains; no
derivative of $J$ is required.

For this proof let $\mathfrak G(d)$ consist of bounded site
matrix pairs $(F,H)$ whose supports lie in balls of radius
$2d$ with centres at distance at most $8d$, and whose
supports have separation at least $d/4$.  The same notation
will be used for scalar or vector pairs after placing them
in fixed matrices of norm at most one.

\begin{proposition}\label{C-417-7}\label{C-417-8}
There is a nondecreasing function $A_g(d)$, finite on
$(0,1/200]$ and tending to zero as $d\downarrow0$, such that
for $(F,H)\in\mathfrak G(d)$, $16a\le d$, and $m_0\le1/2$,
\begin{equation}\label{eq:qcdm2-matrix-source-comparison}
 \left|\kappa_2(\Phi_F,\Phi_H)
       +a^2\operatorname{Tr}(\hat FG_0\hat HG_0)\right|
       \le A_g(d)d^2\|F\|\|H\|.
\end{equation}
Here $G_0=(D[1]+m_0)^{-1}$ includes the full fibre.
The estimate holds on admitted tori and in their
infinite-lattice limit states.  There is also a constant
$C_{\rm fr}(g)$ such that, for nonnegative scalar pairs
in $\mathfrak G(d)$,
\begin{equation}\label{eq:qcdm2-flat-propagator-bound}
 P_1(F,H)\le C_{\rm fr}(g)d^2\|F\|\|H\|,\qquad
 P_\nu(F,H)\le(C_{\rm fr}(g)+A_g(d))d^2\|F\|\|H\|.
\end{equation}
\end{proposition}

\begin{proof}
We first check the operator estimates and local extraction
on $\mathcal M$. We then carry them through the source
chain and the sum over histories, which supplies the cell
bound used in the geometric reduction.

\emph{The inverse and background maps.}
On a history, let $\Omega_1$ be the region of the fine
fermion variables being integrated. The layered operator
on this region has the form
\[
 \mathcal K_i(\mathsf n)
  =(D[U]+m_0)_{\Omega_1\Omega_1}
       +\mathsf W_i+\hat{\mathsf n}
  =\mathcal K_i+\hat{\mathsf n}.
\]
Here $\mathsf W_i$ is the sum of block-average penalties
and $\mathsf s_i$ is the smoothed local scale of
\cite[Section 6]{QCDM1}. Let $\iota$ embed a site spinor
isometrically as a field constant over its $48$ chambers;
then $\Pi_0=\iota\iota^*$ is the chamber-constant projection.
Multiplication by $\mathsf n$
commutes with $\mathsf s_i$, with scalar site functions,
and with the chamber projection $\Pi_0$. Moreover,
$\iota^*\hat{\mathsf n}\iota=\mathsf n$ on the
chamber-independent subspace. These identities hold
because the same matrix acts at each chamber of a site.

The Neumann estimate can be written without a scalar-mass
assumption. Put
\[
 B_i=\mathsf s_i^{-1/2}\mathcal K_i^{-1}\mathsf s_i^{-1/2},
 \qquad
 N_i=\mathsf s_i^{1/2}\hat{\mathsf n}\mathsf s_i^{1/2}.
\]
In the weighted row and column kernel norm,
$\|B_i\|\le C_{\mathcal K}$ and
$\|N_i\|\le\sup_x\mathsf s_i(x)\|\mathsf n(x)\|$.
For $i\le j$, this supremum is at most $2r_M$.
The choice $r_M\le(4C_{\mathcal K})^{-1}$ gives
$\|B_iN_i\|\le2C_{\mathcal K}r_M\le1/2$.
These are the source perturbation bounds of
\cite[Lem. 6.5 and Section 10]{QCDM1}. Thus
\begin{equation}\label{eq:qcdm2-current-Neumann}
 \mathsf s_i^{-1/2}\mathcal K_i(\mathsf n)^{-1}
       \mathsf s_i^{-1/2}
 =\sum_{q\ge0}(-B_iN_i)^qB_i,
 \qquad
 \|\mathsf s_i^{-1/2}\mathcal K_i(\mathsf n)^{-1}
       \mathsf s_i^{-1/2}\|\le2\|B_i\|.
\end{equation}
The same calculation after an exponential conjugation
gives the weighted inverse bound. The local inverses in
the walk construction use the pulled-back site matrix,
so their support and decay bounds are unchanged.
The heavy block has real part at least
$\lambda_\rho+m_0-\|\mathsf n\|$; its bound uses the
norm of the perturbation and the commutation with
$\Pi_0$. Positivity of the complex mass matrix itself
is not needed.

For a background column with prescribed coarse data,
write $\psi_P(\mathsf n)=\mathcal K_i(\mathsf n)^{-1}g_P$.
The resolvent identity gives
\begin{equation}\label{eq:qcdm2-current-background-change}
 \psi_P(\mathsf n)-\psi_P(0)
 =-\mathcal K_i(\mathsf n)^{-1}
                    \hat{\mathsf n}\psi_P(0).
\end{equation}
The row twin obeys the corresponding identity, with the
order of the factors reversed. For the fluctuation
covariance, let $A_i$ be the integrated places at level $i$
and let $q_i$ be the one-step fermion average. Put
$T_A=1_{A_i}q_i\cdots q_1$. The block-kernel parameters are
\[
 \theta=b_F>0,\qquad
 \theta_i=\frac{b_F(1-L^{-1})}{1-L^{-i}}\quad(i\ge1),
\]
as in \cite[Lem. 5.4]{QCDM1}.
The Schur formula at the start of \cite[Section 6]{QCDM1}
then reads, for $i\ge1$,
\[
 \Gamma_i(\mathsf n)^{-1}
 =c_i^{-1}+c_i^{-1}\theta_iT_A
       \mathcal K_{i+1}(\mathsf n)^{-1}T_A^+
                         \theta_ic_i^{-1},
 \qquad c_i=\theta_i+\theta q_{i+1}^+q_{i+1}.
\]
At $i=0$ the identity is
$\Gamma_0(\mathsf n)^{-1}=\mathcal K_1(\mathsf n)^{-1}$.
Here $+$ denotes the holomorphic row twin, agreeing
with the adjoint at real links. The averages do
not depend on $\mathsf n$. Substitution of
\eqref{eq:qcdm2-current-Neumann} therefore gives the
same covariance, background, and Gram estimates as
\cite[Prop. 6.7]{QCDM1}, with its fixed enlarged
constants. All these objects are holomorphic in
$\mathsf n$ on the stated balls.

\emph{The source coordinates and extraction.}
At zero gauge source the independent coordinates are
\[
 \mathfrak c=(0,\mathsf n,(\sigma_i)_{i\le k}),\qquad
 \mathscr D_k^{\mathcal M}
 =\{L^{k+1}\|\mathsf n\|<r_M,
        \ \|\sigma_i\|<B_\sigma\phi_i\ (i\le k)\}.
\]
The doubled domain doubles each radius. Since
$L^{k+1}\|\mathsf n\|<2r_M$ there, the preceding
operator estimates apply through $\mathcal K_{k+1}$
and $\Gamma_k$. The local matrix coordinates
$\sigma_i$, the increments, and the extracted zeroth
moments already take values in $\mathcal M$ in the
construction of \cite[Prop. 10.1]{QCDM1}.

More explicitly, let $k_X(\mathfrak c)$ be the quadratic
kernel at the trivial configuration before extraction,
and set
\[
 \operatorname{ord}_0[k](y)
   =\tfrac12\sum_{y'}[k(y,y')+k(y',y)].
\]
On a torus these sums use the lifted paths in the
definition of the local moment. The increment is
\[
 \nu_{k+1}(y;\mathfrak c)=\mu_{k+1}
 -\mathfrak s_{k+1}^{-1}\sum_X
  \bigl(\operatorname{ord}_0[k_X(\mathfrak c)](y)
             -\operatorname{ord}_0[k_X(0)](y)\bigr),
\]
with the same restriction to polymers of diameter less
than one quarter of the torus side. Let $A_{k+1}(y,z)$
be the moment map of a local mass form, evaluated at
the trivial gauge configuration and $\mathsf n=0$.
It is a linear map on $\mathcal M$. On the infinite
lattice, constant background fields give
\begin{equation}\label{eq:qcdm2-current-matrix-extraction}
 \sum_zA_{k+1}(y,z)N=\mathfrak s_{k+1}N
                  \qquad(N\in\mathcal M).
\end{equation}
On a torus, both this row sum and the corresponding
column sum differ from $\mathfrak s_{k+1}N$ by at most
$C_Ae^{-4c_t\ell/\lambda_{k+1}}\|N\|$. Indeed, the local
pieces which do not wrap around agree with their
infinite-lattice counterparts; the remaining pieces
are bounded by the exponential walk and polymer tails.
For nonzero $\mathsf n$, write $A_{k+1}^{\mathsf n}$
for the moment map. Its change from $A_{k+1}$ has the
weighted column-sum bound of
\cite[Prop. 10.1, proof]{QCDM1}, with factor
$C_{\rm bg}^{\rm w}L^{k+1}\|\mathsf n\|/r_M$.

The extraction acts on the entire matrix coefficient.
It need not leave a kernel with zero zeroth moment:
a varying coefficient $\nu_{k+1}$ leaves differences
of the form
\[
 \sum_z A_{k+1}(y,z)
           [\nu_{k+1}(z)-\nu_{k+1}(y)].
\]
The cell estimates retain these differences, the
source variation, and the torus error. The extraction
forms have coefficients independent of
the gauge configuration and commuting with colour,
which preserves their joint gauge invariance.
The weighted row sums and cell-unit estimates use
operator norms of these coefficients, so apply to
the larger source space with the same constants.

The source chain is
\[
 \sigma_0=0,\qquad
 \sigma_{k+1}=L\sigma_k+
        \nu_{k+1}(\mathfrak c_k(J))-\mu_{k+1},
 \qquad
 \mathfrak c_k(J)=(0,-aJ,(\sigma_i(J))_{i\le k}).
\]
The coordinate Cauchy estimate in the doubled domain
has the same distance as before, with the gauge-source
entry omitted:
\[
 \|\mathfrak c-\mathfrak c'\|_{k,y}
 =\sup_xe^{-\kappa_\sigma|x-y|/L^{k+1}}
 \max\left\{
   \frac{L^{k+1}|\mathsf n-\mathsf n'|(x)}{r_M},
   \max_{i\le k}\frac{|\sigma_i-\sigma_i'|(x)}
                            {B_\sigma\phi_i}\right\}.
\]
Here $x,y$ are in fine lattice coordinates. In particular
\[
 |\nu_{k+1}(\mathfrak c)(y)-\nu_{k+1}(\mathfrak c')(y)|
 \le B_\nu\phi_{k+1}\|\mathfrak c-\mathfrak c'\|_{k,y}.
\]
The induction of \cite[Lem. 10.4]{QCDM1} therefore
keeps the chain in $\mathscr D_k^{\mathcal M}$ and gives
\[
 \|\sigma_k(J)\|
 \le B_\nu S_\uparrow\phi_k
                   \lambda_k\|J\|/r_M.
\]
If $J$ and $J'$ agree outside $S$, the same induction
gives the tail
$C_t\phi_k e^{-\kappa_\sigma\operatorname{dist}(x,S)/\lambda_k}$,
now in physical coordinates. Neither estimate uses a
commutation of $J$ with the spin matrices.

\emph{Symmetries and the sum over histories.}
The four source operations are
\[
 J^\star=\gamma_4J^\dagger\gamma_4,\quad
 J^{\mathsf c}=C_sJ^{\mathsf T}C_s^{-1},\quad
 J\mapsto VJV^{-1},\quad
 J\mapsto S_{\mathsf t}(J\circ\mathsf t^{-1})S_{\mathsf t}^{-1},
 \qquad C_s=\sigma_2\otimes1_\tau\otimes1_{N_f}.
\]
Here $V=e_+\otimes V_1+e_-\otimes V_2$, with
$V_1,V_2\in\mathrm U(N_f)$, acts trivially on spin.
The map $\mathsf t$ is a lattice translation followed by a
signed coordinate permutation preserving the block partition;
$S_{\mathsf t}$ is its spinor lift from
\cite[proof of Prop. 2.5]{QCDM1}. The sign ambiguity
in the lift disappears in conjugation.
They preserve $\mathcal M$ and its norm. At real links,
for example, write $D_J=D[U]+m_0-a\hat J$. Then
\[
 D_J^*=\gamma_4D_{J^\star}\gamma_4,
 \qquad D_J^{\mathsf T}=\mathcal C^{-1}D_{J^{\mathsf c}}\mathcal C.
\]
The flavour and isometry identities are obtained by
conjugating $D_J$ with their respective lifts.
The finite-dimensional integral gives the exact identity
$Z[J^\star]=\overline{Z[J]}$. The diameter restriction in
local extraction requires the torus qualifications of
\cite[Prop. 10.1]{QCDM1}; it does not give an exact
covariance identity for every local increment.

The exact conjugation identity for the sum of the
source energies follows separately, as in
\cite[Prop. 10.6, proof]{QCDM1}. Sum the new activities
and localized normalizers before evaluating their scalar
parts. On the background pairs used in the density
identity, their sum is the logarithm of the whole step
integral. The trivial pair belongs to this domain, and
conjugation carries its scalar part to the complex
conjugate at the conjugated source. Subtracting the real
zero-source value proves the identity for the level
sum. The assignment to individual polymers need not
be preserved.

The history comparison uses the same fine matrix
$\mathsf n$ in the chain and on the history; only the
coordinates $\sigma_i$ differ. Its distance estimate
therefore applies with the same constants. In the
history integrals, the fermion covariance and stored
background columns are those of
\eqref{eq:qcdm2-current-Neumann} and
\eqref{eq:qcdm2-current-background-change}. Schur's
test, the row bounds, and the Berezin Gram bound depend
only on their norms. Products such as
$(1-\mathfrak s_j^{-1}\sum_yA_j^{\mathsf n}(y,z))a_{X,p}(z)$
are estimated in that order, so no commutation of
their matrix factors is assumed.

Write $\mathcal H$ and $\bar{\mathcal H}$ for the composed
column and row background maps from the variables before
step $j$ to the original fermion variables. They are the
maps obtained by successive square completion in
\cite[Section 5.3]{QCDM1}. In the following formula,
the subscript $AA$ restricts both coarse arguments to the
variables integrated at that step. The determinant
variation used in the local energy is
\[
 \log\det\Gamma_{j-1}(\mathsf n)
       -\log\det\Gamma_{j-1}(0)
 =\int_0^1\operatorname{Tr}\left[
  \Gamma_{j-1}(t\mathsf n)^{-1}
  (\bar{\mathcal H}\hat{\mathsf n}\mathcal H)_{AA}
                         (t\mathsf n)\right]dt.
\]
Expand the inverse and the two background maps into
walks and assign each trace to their joint support.
The fibre trace is bounded by $8N_f$ times the product
of the operator norms. This proves the same local
energy estimate as \cite[Prop. 10.6]{QCDM1}.
A change of source outside a polymer can be seen only
by a walk reaching that change; reserving part of its
decay gives the source-distance tail.

At zero gauge source, every variation of the gauge
weight and every associated Gaussian or linear term
vanishes. The remaining closed fermion loops are
gauge invariant. Their first derivative in curvature
at the trivial configuration is zero, since a linear
form invariant under the adjoint action of
$\mathrm{SU}(2)$ vanishes. Consequently the history
majorants and their sums in
\cite[Prop. 10.14]{QCDM1} apply with the same
constants. The bare density is entire in $J$, and
its coefficients are continuous on the compact link
space, as required when that sum is identified with
the original partition function. This proves the
first-order source and energy bounds needed here.
No second-order improvement in $\phi$ for the enlarged
source space is being used.

\emph{Reflection and the cell bound.}
The cell bound requires the correct reflection
of a spin matrix.  Directly on generators,
\[
 \Theta_P(\bar\psi(v)J(v)\psi(v))
  =\bar\psi(\theta v)\gamma(\hat n)J(v)^\dagger
                                  \gamma(\hat n)\psi(\theta v).
\]
For the reflection lift $S_P=\gamma_4\gamma(\hat n)$ this
is the linear reflected source followed by $\star$.
The norm is preserved. The operation $\star$ commutes
with conjugation by the reflection lift, since
$\gamma_4S_P=-S_P\gamma_4$ and the sign drops out of
the conjugation. For a word $\mathsf r$ in the board
reflections, define
\[
 \mathsf r_{\rm lin}J
 =S_{\mathsf r}(J\circ\mathsf r^{-1})S_{\mathsf r}^{-1},
 \qquad
 \mathsf r_\theta J=
 \begin{cases}
 \mathsf r_{\rm lin}J,&|\mathsf r|\text{ even},\\
 (\mathsf r_{\rm lin}J)^\star,&|\mathsf r|\text{ odd}.
 \end{cases}
\]
If $J$ is supported strictly inside one cell, the reflected
exponentials are even and have disjoint supports. Hence
their product is
\[
 \mathrm{Fold}\,e^{\Phi_J}=e^{\Phi_{J^{\rm ch}}},
 \qquad J^{\rm ch}=\sum_{\mathsf r}\mathsf r_\theta J,
 \qquad \|J^{\rm ch}\|=\|J\|.
\]
The chessboard estimate applies to this product. The
partition-function bound just obtained and the separated
energy estimate in \cite[Prop. 10.16]{QCDM1}
then give the cell bound of \cite[Prop. 10.17]{QCDM1}.
Odd reflection words cause no change in its real energy
bound, because $E_K[J^\star]=\overline{E_K[J]}$.
For sources in the central ball of radius $1/4$ this
bound takes the form
\begin{equation}\label{eq:qcdm2-current-cell-bound}
 \Psi(J):=e^{-\Lambda(J)}\frac{Z[J]}{Z[0]},
 \quad \Psi(0)=1,\quad |\Psi(J)|\le B'_Z,
 \quad B'_Z=e^{c_V+\mathscr C_*},
\end{equation}
Here $\mathscr C_*$ bounds the source count uniformly for
supports in the central ball of radius $1/4$, as in
\cite[Prop. 10.17]{QCDM1}, and
$\Lambda(J)=\mathsf N^{-1}\sum_{\mathsf r}
 (E_K[\mathsf r_{\rm lin}J]-E_K[0])$
and $\mathsf N=\ell^3$.

\emph{Packing and averaging.}
A pair in $\mathfrak G(d)$ lies in a ball of radius
$6d$. Choose a grid of mesh $14d$ inside the ball of
radius $1/4-7d\ge0.215$, and move each grid point to
the nearest permitted lattice translate of the centre
of this pair. The displacement is at most
$\sqrt3a/2<d/8$. The translated balls stay inside the
cell ball, and their separation is at least
$2d-\sqrt3a\ge d$. An inscribed cube supplies at least
\[
 N\ge\left\lfloor\frac{0.43}{14\sqrt3d}\right\rfloor^3
     \ge(112d)^{-3}
\]
points: $0.43/(14\sqrt3)>0.0177$, and
$(0.0177-1/112)d^{-1}\ge1$ for $d\le1/200$.

For separated matrix sources put
$\|F\|_\circ=\|F\|/\varrho_J$ and similarly for $H$.
The function $\Psi(z_1F+z_2H)$ is bounded by $B'_Z$
on $|z_1|<2/\|F\|_\circ$, $|z_2|<2/\|H\|_\circ$.
On the bidisc smaller by $1/(4B'_Z)$, Schwarz's lemma
on complex lines gives $|\Psi-1|\le1/2$. Its logarithm
exists there, and Cauchy's estimate gives
\[
 |\partial_1\partial_2\log\Psi(0)|
       \le4(B'_Z)^2\|F\|_\circ\|H\|_\circ.
\]
Apply this to
$\sum_t\varepsilon_tF^t$ and $\sum_t\varepsilon_tH^t$,
with independent signs $\varepsilon_t\in\{-1,1\}$.
The supports are disjoint, so these sums have the
norms of one pair. Since
$\mathsf E_\varepsilon\varepsilon_t\varepsilon_{t'}
=\delta_{tt'}$, averaging the bilinear derivative
retains only matching translates. Translation invariance
then yields
\[
 \kappa_2(\Phi_F,\Phi_H)
 =\frac1{N\mathsf N}\sum_{t,\mathsf r}
       \partial_1\partial_2E_K(F^{t,\mathsf r},H^{t,\mathsf r})
       +R_{\rm cell},
 \quad
 |R_{\rm cell}|\le C_{\rm cell}(g)d^3\varrho_J^{-2}\|F\|\|H\|,
\]
where $\mathsf N=\ell^3$ is the number of unit cells.
One may take $C_{\rm cell}(g)=4\cdot112^3
\exp(2c_V+2\mathscr C_*)$ in the notation of that lemma.

\emph{The energy decomposition.}
At zero gauge source the energy decomposition
\cite[Lem. 11.4]{QCDM1} has just its fermion and interaction
parts. The relative gauge weight is zero, as are the
linear gauge source and its Gaussian variations. In
particular the term of fourth order in the gauge source
is zero. Completing the fermion square at each step is
an algebraic identity for every $J\in\mathcal M$. The
coarse determinant from one step is the input determinant
of the next. Write $\mathbb D_K(1;J)$ for the final coarse
Dirac matrix after these $K$ square completions at the
trivial gauge configuration with source $J$. Their product gives
\[
 \sum_{j,X}E_j^{\rm F}(X;J)
 =\log\frac{\det(D[1]+m_0-a\hat J)}{\det(D[1]+m_0)}
       -\log\frac{\det\mathbb D_K(1;J)}{\det\mathbb D_K(1;0)}.
\]
This is an identity of logarithmic germs at zero.
On the larger level-source domains its continuation
is defined by the product of the invertible fluctuation
determinants. Separate global logarithms of the two
displayed determinants are not required.

For completeness, the small factor in the interaction
energy comes from the part of the step potential left
after the explicit normalizers and Gaussian terms.
On the level-$j$ domain the induction in the proof of that
lemma applies to the independent matrix coordinates
$(0,-aJ,(\sigma_k)_{k<j})$, with
$\|\sigma_k\|<2B_\sigma\phi(x_k)$, and gives
\[
 \|\mathcal V_i^{\rm rest}\|
       \le C_{\rm rest}\phi(x_{i-1}),\qquad i\le j,
 \qquad \|\sigma_i\|<2B_\sigma\phi_i.
\]
At zero gauge source its gauge vertices are independent
of $J$, its fermion vertices have the source-step
bound proved above, and its closed-loop variations
have the weighted trace bound. The localized variation
of an older activity has the bound from the preceding
induction level. Thus the scalar output at the trivial
configuration, after removing the normalizers, is at
most $C_w^{\rm s}\phi(x_{j-1})$. Subtracting its value
at zero source gives
\begin{equation}\label{eq:qcdm2-current-interaction-energy}
 |E_j^{\rm I}(X;J)|
 \le2C_w^{\rm s}\phi(x_{j-1})e^{-\kappa d_j(X)}.
\end{equation}
This induction also retains holomorphy and locality in the
independent coordinates. The matrix-source vertices and
covariances are holomorphic and read $J$ on their assigned
walks. Scalar localization preserves the parameter dependence
of the old arrays; a member assigned to $X$ reads only their
restrictions to $X$. The cluster expansion, its pointed
families and the Cauchy integrals defining the remainder
therefore have the same properties. These are exactly the
coordinate clauses used in the interaction-tail proof of
\cite[Lem. 11.4]{QCDM1}; the matrix extension above verifies
their inputs.

Let $S=\operatorname{supp}(J-J')$ and
$D=\operatorname{dist}(X,S)>0$. Multiply the remainder
potential by $z$. Its scalar output at the trivial pair,
before the normalizers are removed, is a holomorphic
function $\mathcal E(z;\mathfrak c)$ on the full independent
coordinate domain, with
\[
 |\mathcal E(z;\mathfrak c)|
 \le K_M(c_{\rm cl}+1)e^{-\kappa d_j(X)},\qquad
 |z|\le R=\frac{c_{\rm cl}}{4C_{\rm rest}\phi(x_{j-1})}\ge2.
\]
The two source chains have the same $J$ on $X$, and
their induced coordinates differ there by at most
$C_{\rm t}\phi(x_k)e^{-\kappa_\sigma D/\lambda_j}$.
Keep $J'$ fixed and replace all $\sigma_k(J')$ simultaneously
by $\sigma_k(J')+\zeta(\sigma_k(J)-\sigma_k(J'))1_X$.
Since the chain lies in the inner window
$\|\sigma_k(J')\|\le B_\sigma\phi(x_k)$, this curve stays
in the coordinate domain for $|\zeta|<1/\theta$, where
$\theta=(C_{\rm t}/B_\sigma)e^{-\kappa_\sigma D/\lambda_j}$.
Locality and Schwarz's lemma bound the difference of the two
outputs by
\[
 C_{\rm tail}K_M(c_{\rm cl}+1)e^{-\kappa d_j(X)}
                e^{-\kappa_\sigma D/\lambda_j},\qquad
 C_{\rm tail}=2\max\{1,C_{\rm t}/B_\sigma\}.
\]
For $\theta\ge1$ the two separate bounds suffice; for
$\theta=0$ the outputs agree. There is no sum over $k$
in this interpolation.

The output difference vanishes at $z=0$, where its scalar
potential and quadratic origin read only the zero gauge
source. At $z=1$ it equals the difference of $E_j^{\rm I}$:
the gauge normalizer and Gaussian part agree, and the free
fermion normalizer cancels with $E_j^{\rm F}$. Schwarz's
lemma in $z$ now supplies $R^{-1}$. Thus the constant is
at most $4C_{\rm rest}C_{\rm tail}K_M(1+c_{\rm cl}^{-1})$,
already included in the companion's $C_I$, and
\begin{equation}\label{eq:qcdm2-current-interaction-tail}
 |E_j^{\rm I}(X;J)-E_j^{\rm I}(X;J')|
 \le C_I\phi(x_j)e^{-\kappa d_j(X)}
       e^{-\kappa_\sigma\operatorname{dist}(X,S)/\lambda_j}.
\end{equation}
If $X$ meets $S$, use the two bounds
\eqref{eq:qcdm2-current-interaction-energy}; the fixed
comparison of $\phi(x_{j-1})$ and $\phi(x_j)$ is included
in $C_I$. This explains why the interaction tail, as
well as its undifferentiated size, retains a factor
$\phi$.

\emph{Summing the mixed derivatives.}
For the averaged interaction sum, the level-$j$ source
radii are $R_J/(\lambda_j\|F\|)$ and
$R_J/(\lambda_j\|H\|)$.  At levels with $\lambda_j<d$
the separation $d/4$ gives the tail
$e^{-\kappa_\sigma d/(16\lambda_j)}$ and at most
$C(d/\lambda_j)^3$ cubes near one support.
At levels $\lambda_j\ge d$, use the independent signs
for the packed sources before applying Cauchy's
estimate. The packed pairs and their board images
meet at most $CN\mathsf N(d/\lambda_j)^3$ cubes:
their union occupies a fixed-volume region in each
unit cell and $N\ge c d^{-3}$. The signs have the
same source norm as a single member. Dividing by
$N\mathsf N$ therefore leaves the factor
$C(d/\lambda_j)^3$, even at these coarser levels.

After summing the polymers through each cube,
Cauchy's estimate and
\eqref{eq:qcdm2-current-interaction-tail} bound the
mean by $C R_J^{-2}\|F\|\|H\|$ times
\begin{equation}\label{eq:qcdm2-current-level-sum}
 \sum_j\phi(x_j)\lambda_j^2(d/\lambda_j)^3
 \begin{cases}
 e^{-\kappa_\sigma d/(16\lambda_j)},&\lambda_j<d,\\
 1,&\lambda_j\ge d.
 \end{cases}
\end{equation}
The two powers of $\lambda_j$ are the inverse Cauchy
radii; the other factor is the number of cubes after
averaging. To see its summability, put $t=\mathsf b/d$.
The scale comparison $x_j\asymp\mathsf b/\lambda_j$
and $\phi(u)=u^{-1/2}(1+\log u)^{p_0}$ give, with
fixed comparison constants,
\[
 \frac{\phi(\mathsf b/\lambda)}{\phi(t)}
 \le
 \begin{cases}
 (\lambda/d)^{1/2}(1+\log(d/\lambda))^{p_0},&\lambda\le d,\\
 (\lambda/d)^{1/2},&\lambda\ge d.
 \end{cases}
\]
All arguments are in the admitted decreasing range of
$\phi$. Consequently, after division by $d^2\phi(t)$,
the terms with $\lambda_j\ge d$ form a geometric sum
with ratio $L^{-1/2}$. For $\lambda_j<d$, the sum is
bounded by a fixed multiple of
\[
 \sum_{u\ge0}L^{u/2}(1+(u+1)\log L)^{p_0}
                e^{-(\kappa_\sigma/16)L^{u-1}}<\infty.
\]
This is the summation in \cite[Lem. 11.6]{QCDM1},
with the support separation $d/4$. It gives
\[
 |E^{\rm I}_{\rm averaged}|
       \le C_M^\natural R_J^{-2}\phi(\mathsf b/d)
                                      d^2\|F\|\|H\|.
\]
\emph{The last determinant.}
Put $A(z_1,z_2)=\mathbb D_K(1;z_1F+z_2H)$ and
$A=A(0,0)$. The real-part bound gives
$\|A^{-1}\|\le(c_b\mathsf m)^{-1}$.
Differentiating the inverse formula inserts $-a\hat F$
between two background maps. Their squared row sums
are $O(a^2)$, so the first derivative satisfies
\[
 \|A_1\|_1\le Ca^3\sum_x\|F(x)\|
       \le Cd^3\|F\|,
 \qquad \|A_2\|_1\le Cd^3\|H\|.
\]
Here $\|\cdot\|_1$ is the trace norm. The supports lie
in balls of radius $2d$, so their site counts are
bounded by a constant times $(d/a)^3$.
In the mixed derivative the inserted matrices are
separated by one free propagator, with bound
$Ca^2(16d^{-2}+\mathsf m^{-1})$ at separation $d/4$.
The same row estimate therefore gives
\[
 \|A_{12}\|_1
 \le C(d^4+d^6\mathsf m^{-1})\|F\|\|H\|.
\]
The logarithmic derivative is
\[
 \partial_1\partial_2\log\det A(z)\big|_{z=0}
 =\operatorname{Tr}(A^{-1}A_{12}
                      -A^{-1}A_2A^{-1}A_1).
\]
Using the trace norm for the localized derivatives,
and the operator norm for $A^{-1}$, bounds it by
\[
 C_{\rm last}^\natural
          (d^6\mathsf m^{-2}+d^4\mathsf m^{-1})\|F\|\|H\|,
\]
as in the last paragraph of \cite[Lem. 11.5]{QCDM1}.
No factor proportional to the total volume enters
these two traces.
The free loop is unchanged by the translates and reflected
matrix sources.  Thus \eqref{eq:qcdm2-matrix-source-comparison}
holds with
\[
 \alpha_g(d)=C_{\rm cell}(g)\varrho_J^{-2}d
       +C_M^\natural R_J^{-2}
                 \sup_{\mathsf b\in I_g}\phi(\mathsf b/d)
       +C_{\rm last}^\natural
                    (d^4{m_-}^{-2}+d^2{m_-}^{-1}),
 \qquad A_g(d)=\sup_{0<s\le d}\alpha_g(s).
\]
This function is finite and tends to zero.  The mass
labels lie in a compact interval, and the same is true
of $\mathsf b\in I_g$ at fixed $g$.
On the infinite lattice, pass from the tori by convergence
of local Gaussian functionals and by the vanishing of
the massive free image terms.

For the positive free form, the propagator estimate
in the proof of \cite[Lem. 11.5]{QCDM1} gives
$\|G_1(u,u')\|\le C'_G(g)|n|^{-2}$, with
$n=(x_u-x_{u'})/a$. Its leading term has this bound
because $(1+t)e^{-t}\le1$ and $te^{-t}\le e^{-1}$.
The remainder is $O(|n|^{-3})$, and the torus images
are bounded by $C'_{\rm im}(\ell/a)^{-2}$.
These constants are included in $C'_G(g)$.
The support balls contain at most
\[
 \frac{4\pi}{3}(2d/a+\sqrt3/2)^3
       \le36\pi(d/a)^3
\]
sites each, with 48 places per site, and
$|n|\ge d/(4a)$. The four-dimensional spin-colour
Hilbert--Schmidt norm is at most twice the operator
norm. Thus
\begin{align*}
 P_1(F,H)
 &\le a^2\,8C'_G(g)^2\,48^2(36\pi)^2
        (d/a)^6(4a/d)^4\|F\|\|H\|\\
 &=2048(36\pi)^2 48^2C'_G(g)^2d^2\|F\|\|H\|.
\end{align*}
This gives the stated choice of $C_{\rm fr}(g)$.
Apply \eqref{eq:qcdm2-matrix-source-comparison} with
$\Gamma=i\tau_1\otimes t$ and use \cref{C-417-3} to obtain
the second bound in \eqref{eq:qcdm2-flat-propagator-bound}.
\end{proof}

\begin{theorem}[Comparison with the free current loop]\label{C-417-9}
Suppose $16a\le d\le1/200$ and $m_0\le1/2$.
Let $f,h$ be real bounded
functions, each supported in a ball of radius $d$, with
support separation between $d$ and $6d$.  Then
\begin{equation}\label{eq:qcdm2-current-comparison}
 |E_\nu(f,e;h,e')-E_1(f,e;h,e')|
 \le\left(A_g(d)+C_{\rm pt}(g)(a/d)^{1/2}\right)
                                   d^2\|f\|_\infty\|h\|_\infty,
\end{equation}
where
$C_{\rm pt}(g)=2C_{\rm split}(2C_{\rm fr}(g)+A_g(1/200))$.
\end{theorem}

\begin{proof}
The pairs $(1_{B'_f},1_{B_h})$ and
$(1_{B_f},1_{B'_h})$ in \cref{C-417-6} belong to
$\mathfrak G(d)$.  Indeed the support enlargement is
at most $d/4+5a\le9d/16$, so each enlarged support
fits in a ball of radius $2d$, and the remaining
separation is at least $d-9d/16-a\ge d/4$.
The shifted site pair $(F^{f,e},F^{h,e'})$ also belongs
to this class.  Apply \eqref{eq:qcdm2-flat-propagator-bound}
to the two point-splitting errors, one under $\nu$ and
one at $U=1$, and apply
\eqref{eq:qcdm2-matrix-source-comparison} to $N_F,N_H$.
Adding the three differences gives
\eqref{eq:qcdm2-current-comparison}.
\end{proof}

\subsection{The free tensor and the current classes}

\begin{lemma}\label{C-417-13}
For a test pair, $16a\le d$ and $10\mathsf m d\le1$,
\begin{align}
 \left|E_1^{\rm s}(f,e;h,e')
       -\frac1{4\pi^2}\iint
          f(x)h(y)\langle e,k(x-y)e'\rangle\,dx\,dy\right|
 \le C_{\rm cr}\left(\frac a{\vartheta d}
                +\mathsf m d+(d/\ell)^2\right)
                                      d^2\|f\|_\infty\|h\|_\infty.
 \label{eq:qcdm2-free-current}
\end{align}
The last term is absent on the infinite lattice.
The constant depends only on the fixed model and
${m_-}\ell_G$, where $\ell_G$ is the least admitted side.
\end{lemma}

\begin{proof}
The free propagator expansion in the proof of
\cite[Lem. 11.5]{QCDM1} gives, in its light upper block,
\[
 G_1(n)=v(n)\sigma(n)+s(n)+O(|n|^{-3}),\qquad
 v(n)=\frac{(1+m_0|n|)e^{-m_0|n|}}{4\pi|n|^3},\quad
 s(n)=\frac{m_0e^{-m_0|n|}}{4\pi|n|}.
\]
The light factor between individual chambers is $48^{-1}$;
the sum over the two chamber indices cancels its square.
The mixed and heavy parts have bounds $C|n|^{-3}$ and
$C|n|^{-4}$, hence cost $O(a/d)$ in this loop.
Use
\[
 \operatorname{tr}(\sigma(n)\sigma(F)\sigma(n)\sigma(H))
       =2\bigl(2(n\cdot F)(n\cdot H)-|n|^2F\cdot H\bigr).
\]
The sign of the second propagator and the colour factor
two give for the $v^2$ term
\[
 \frac{(1+t)^2e^{-2t}}{4\pi^2}\langle F,k(n)H\rangle,
 \qquad t=m_0|n|.
\]
The two $vs$ terms involve three Pauli matrices.
Their ratio to the size $4v^2|n|^2|F||H|$ of the
leading term is at most $2t$, since
\[
 \frac{s(n)}{v(n)|n|}=\frac{t}{1+t}\le t.
\]
The $s^2$ term has ratio at most $t^2$, and
$(1+t)^2e^{-2t}=1+O(t^2)$. On the shifted supports,
\[
 15d/16\le d-a\le|x-y|\le10d+a\le11d,
 \qquad t\le11\mathsf m d\le2.
\]
The two mixed terms therefore give an
$O(\mathsf m d)$ error; they cannot be replaced by
an $O((\mathsf m d)^2)$ error in this estimate.

Homogeneity gives $a^2k(n)=a^6k(x-y)$.
The discrete error just estimated is bounded by
\[
 C(\mathsf m d+a/d)a^6
   \sum_{x,y}|F(x)||H(y)||x-y|^{-4}
 \le C(\mathsf m d+a/d)d^2\|f\|_\infty\|h\|_\infty.
\]
Replacing midpoint values by site values uses
\[
 |F^{f,e}(x)-f(x)e|
 \le\tfrac{\sqrt3}{2}a\operatorname{Lip}f
 \le\frac{a}{\vartheta d}\|f\|_\infty,
\]
and the same bound for $H$. For the remaining Riemann
sum, $|k(x-y)|\le Cd^{-4}$ and
$|\nabla k(x-y)|\le Cd^{-5}$ on the separated supports.
The change across a lattice cube is therefore bounded
by a constant times
\[
 a\bigl(d^{-4}\operatorname{Lip}f\|h\|_\infty
       +d^{-4}\|f\|_\infty\operatorname{Lip}h
       +d^{-5}\|f\|_\infty\|h\|_\infty\bigr).
\]
Its support in the product space has volume $O(d^6)$.
Summing the cube errors gives
$Ca/(\vartheta d)$ times
$d^2\|f\|_\infty\|h\|_\infty$.
On the torus the antiperiodic image sum is bounded
entrywise by $C({m_-}\ell_G)(\ell/a)^{-2}$;
its contribution to the loop is $C(d/\ell)^2$ at
the same normalization.  These bounds prove the claim.
\end{proof}

\begin{proof}[Proof of \cref{C-417-14}]
For $T=\tau_{a_0}\otimes t$ use $t_+=t$,
$t_-=\epsilon_{a_0}t$, with $\epsilon_0=1$ and
$\epsilon_3=-1$.  Equation \eqref{eq:qcdm2-current-loop}
is therefore $\operatorname{tr}(tt')$ times
$E_\nu+\epsilon_{a_0}\epsilon_{b_0}\overline{E_\nu}$.
The leading tensor in \cref{C-417-13} is real; adding
its conjugate multiplies it by
$1+\epsilon_{a_0}\epsilon_{b_0}=2\delta_{a_0b_0}$.
The remaining error is at most twice the sum of the
interacting-to-free difference, the free point-splitting
error, and \eqref{eq:qcdm2-free-current}.
Divide by the lower bound in
\eqref{eq:qcdm2-current-Klower}.  For example, take
\[
 \epsilon_V(g,\vartheta;d)
 =\frac{2\cdot10^4}{c_B\vartheta^2}
       \left[A_g(d)+C_{\rm cr}
          ({m_+}d+(d/\ell_G)^2)\right],
 \qquad c_B=\frac1{2\pi^2},
\]
and
\[
 C_V'(g)=\frac{2\cdot10^4}{c_B}
       (C_{\rm pt}(g)+2C_{\rm split}C_{\rm fr}(g)+C_{\rm cr}).
\]
This gives \eqref{eq:qcdm2-current-law-error}; when the
flavour trace vanishes the exact loop formula gives zero.
At fixed $d,f,h$, both mesh terms tend to zero with $K$.
The displayed formulas bound the cumulants and every
limit point as claimed.  If torus sides vary, the
assumption $20d\le\ell$ makes the Euclidean pair and
its two kernel integrals the same at every cutoff.
\end{proof}

\begin{proposition}[A nonzero current class]\label{prop:qcdm2-nonzero-current}
For every sufficiently small fixed $d>0$ and all sufficiently
large $K$, every nonzero Hermitian flavour matrix $t$ and
$T=\tau_{a_0}\otimes t$, $a_0\in\{0,3\}$, admit a current
supported on the positive side of
$P_{1/2}=\{x_3=a/2\}$ whose class in the lattice
reflection-positive Hilbert space is nonzero.
More precisely, let $\vartheta_c=1/500$,
$z_d=(5d/8+a/2)e_3$, and
$f(x)=(d/8-|x-z_d|)_+$.  For $\mu=1,2,3$,
\begin{equation}\label{eq:qcdm2-current-norm}
 \omega_\nu\bigl(J^T(fe_\mu)\Theta_{P_{1/2}}J^T(fe_\mu)\bigr)
       \ge\tfrac14c_B\operatorname{tr}(t^2)
                  \mathcal K_4(f,f\circ\theta)>0.
\end{equation}
The bound is uniform in admitted labels and torus-limit
states.  It also holds on admitted tori with $20d\le\ell$.
\end{proposition}

\begin{proof}
The two tents form a test pair at scale $d$ and
$\vartheta_c$.  Their centres are $5d/4$ apart and
their support radii are $d/8$.  Also
$\operatorname{Lip}f=1$, $\|f\|_\infty=d/8$, and
$\int f=(\pi/1536)d^3\|f\|_\infty
\ge\vartheta_cd^3\|f\|_\infty$.
For large $K$, every endpoint of a bond with midpoint
in its support is strictly above $P_{1/2}$.

On generators, reflection reverses the ordered edge
and sends $\mathsf j^T(v,v')$ to
$-\mathsf j^{T^*}(\theta v,\theta v')$.
Antilinearity also sends $i$ to $-i$.  After summing
the oriented bonds this gives
\begin{equation}\label{C-402-26}
 \Theta_PJ^T(f)=J^{T^*}(\varrho\bar f\circ\theta),
\end{equation}
where $\varrho$ is the linear reflection.  The positive
form of \cite[Thm. B and Thm. A]{QCDM1} therefore
evaluates the left side of \eqref{eq:qcdm2-current-norm}
by the two-current law at $e=e_\mu$ and
$e'=\varrho e_\mu$; the current means are zero.

On these supports $n=x-y$ has $n_3\ge d$ and
$|n_\perp|\le d/4$, so
\[
 \langle e_\mu,k(n)\varrho e_\mu\rangle
       \ge\frac{15}{17}|n|^{-4}.
\]
Choose $d$ so that $\epsilon_V(g,\vartheta_c;d)<1/4$,
and then $K$ so that the two mesh errors in
\eqref{eq:qcdm2-current-law-error} total at most $1/4$.
The leading term less the error is at least
$(15/17-1/2)c_B\operatorname{tr}(t^2)\mathcal K_4$,
which proves \eqref{eq:qcdm2-current-norm}.
The same supports lie in the positive slab of the
torus when $20d\le\ell$.
\end{proof}

\section{Analytic dependence at fixed volume}\label{sec:analytic-labels}

A constant source changes the bare coupling or the bare mass.
This observation connects the source estimates to analytic
dependence on the physical parameters. The partition function
itself need not be bounded as the cutoff is removed. The useful
quantity is the partition function after the same vacuum energy
has been removed from its upper and lower bounds.

Throughout this section the physical torus $T_\ell$ is fixed.
Fix $0<g\le\gamma$ and an allowed side $\ell\in\mathfrak L(g)$,
and write
\[
 I_g=[g^{-2}+\bar B,\,L(g^{-2}+B_\sharp)-\bar B],
 \qquad W=I_g\times[m_-,m_+].
\]
The constants and the cofinal set $\mathfrak L(g)$ are those of
\cite[Prop. 10.20]{QCDM1}. In particular the estimates in this section
hold for every $K\ge1$ and every $(\mathsf b,\mathsf m)\in W$.
We use the chamber operator with $r=1/2$.

For a source $s=(w,J)$ put
\[
 Z_K[s]=\int e^{-(\mathsf b/a)A(U)+O_K(w)}
       \det\bigl(D[U]+a\mathsf m-a\widehat J\bigr)\,dU,
 \qquad a=L^{-K}.
\]
The determinant includes both internal blocks and all flavours.
Thus $Z_K[0]>0$. The source $w$ is a function on plaquette
centres and $J$ takes values in the complex span of the physical
meson basis $\mathfrak B$. The definitions of $O_K$ and the meson insertions are
those in \cref{sec:qcdm2-setting}. The action-source norm is
$\|w\|_\sharp=\max\{\|w\|_\infty,40M\operatorname{Lip}w\}$.
Let
\[
 \mathcal P=\{(w,J):\|w\|_\sharp\le2\varrho,
                              \ \|J\|_\infty\le2\varrho_J\}.
\]
The radii are independent of $K$ and $\ell$, and
$2\varrho_J\le c_Jm_-$, $c_J\le1/2$.

We first recall precisely the source input. It is the combination
of the local-energy, history and counting estimates of
\cite[Props. 10.6 and 10.14, and the counting estimate
preceding Prop. 10.16]{QCDM1}.
There is a holomorphic energy
\begin{equation}\label{eq:qcdm2-energy-decomposition}
 E_K(s)=E_K(0)+\sum_{j=1}^K\sum_{Y\in\mathcal D_j}e_j(Y;s),
 \qquad Z_K[s]=e^{E_K(s)}W_K[s],
\end{equation}
defined on a neighbourhood of $\mathcal P$. Here $Y$ is a
connected union of blocks of physical side $M\lambda_j$,
$\lambda_j=L^{j-K}$, and $d_j(Y)$ is the block-tree length
used in the construction. For fixed positive constants,
\begin{align}
 |e_j(Y;s)|&\le C_Ee^{-\kappa d_j(Y)},\label{eq:qcdm2-local-energy}\\
 |e_j(Y;s)-e_j(Y;s')|&\le C_Ee^{-\kappa d_j(Y)}
 e^{-\kappa_\sigma\operatorname{dist}(Y,\operatorname{supp}(s-s'))/\lambda_j}.
 \label{eq:qcdm2-energy-tail}
\end{align}
Distances in this section are physical. The geometry gives
$\operatorname{diam}Y\le c_DM\lambda_j(1+d_j(Y))$.
The construction chooses $\kappa$ so that summing
$e^{-\kappa d_j(Y)/2}$ over polymers containing one specified
block costs at most $C_{\rm cnt}$. Finally,
\begin{equation}\label{eq:qcdm2-partition-window}
 |W_K[s]|\le e^{(C+2)\ell^3},\qquad
 W_K[0]\ge e^{-C_P(1+\log X_g)\ell^3},\qquad
 X_g=L(g^{-2}+B_\sharp).
\end{equation}
The conjugation of the source makes $E_K$ real for constant real
coupling and mass shifts. No positivity at complex sources is
used.

\subsection{A zero-free neighbourhood}

Set
\begin{equation}\label{eq:qcdm2-analytic-radii}
 V_\ell=\exp\{[C+2+C_P(1+\log X_g)]\ell^3\},\qquad
 \rho_b=\frac{\varrho}{1+V_\ell},\qquad
 \rho_m=\frac{\varrho_J}{1+V_\ell}.
\end{equation}
Let $W_{\mathbb C}$ be the union, over $(b,m)\in W$, of the
open polydiscs $|b'-b|<\rho_b$, $|m'-m|<\rho_m$.

\begin{proposition}[Complex parameters]\label{C-403-41}
For every $K$, the partition function
$Z_{K,\ell}(b',m')=Z_K[0]$ formed with bare parameters
$b'/a$ and $am'$ has no zero on $W_{\mathbb C}$.
For a real base point $(b,m)\in W$, write
$c=(- (b'-b)1,-(m'-m)1)$ and
$R_K(b',m')=W_K[c]/W_K[0]$, with the energies and partition
functions based at $(b,m)$. Then
\begin{equation}\label{eq:qcdm2-zero-free}
 |R_K-1|\le(1+V_\ell)
 \max\left\{\frac{|b'-b|}{2\varrho},
             \frac{|m'-m|}{2\varrho_J}\right\}<\frac12.
\end{equation}
On this polydisc the pressure has the decomposition
\[
 \ell^{-3}\log Z_{K,\ell}(b',m')
=\ell^{-3}E_K(c)+\ell^{-3}\log W_K[c].
\]
The logarithm is chosen real at the real base point, and
\[
 |\log R_K|\le\log2,\qquad
 \ell^{-3}|\log W_K[c]|
 \le C+2+C_P(1+\log X_g)+\ell^{-3}\log2.
\]
Consequently every sequence of cutoffs has a subsequence on
which $\ell^{-3}\log R_K$ converges locally uniformly on this
polydisc, with all parameter derivatives.
\end{proposition}

\begin{proof}
The constant source identity
$Z_{K,\ell}(b',m')=Z_K[c]$ holds before any integration:
the action coefficient becomes $b'/a$, and the matrix in the
determinant becomes $D+am'$. The ratio $R_K$ is holomorphic on
the larger polydisc with radii $2\varrho,2\varrho_J$,
$R_K(b,m)=1$, and $|R_K|\le V_\ell$ by
\eqref{eq:qcdm2-partition-window}.
For a nonzero displacement let $r$ be the maximum in
\eqref{eq:qcdm2-zero-free}. Restrict $R_K-1$ to the complex
line through that displacement, rescaled so that its unit
disc lies in the larger polydisc. Dividing by $1+V_\ell$
gives a function of modulus at most one, vanishing at zero.
The given displacement corresponds to radius $r$ on this
disc. Schwarz's lemma proves
\eqref{eq:qcdm2-zero-free}. Thus $R_K$ has a logarithm with
$|\log R_K|\le\sum_{n\ge1}2^{-n}/n=\log2$.
The remaining bounds follow from the positive real lower
bound and the upper bound in \eqref{eq:qcdm2-partition-window}.
Local boundedness gives a normally convergent subsequence;
Cauchy's formula gives convergence of every derivative.
\end{proof}

For example, if $p+q\ge1$, differentiation at a real base point
gives
\begin{equation}\label{eq:qcdm2-integrated-cumulants}
 \left|S^T_{p+q,K}((0,1)^p,(1,1)^q)
       -(-\partial_{b'})^p(-\partial_{m'})^q
          E_K(c(b',m'))\big|_{(b',m')=(b,m)}\right|
 \le p!q!\log2\,\rho_b^{-p}\rho_m^{-q}.
\end{equation}
The RG base point $(b,m)$ is held fixed in this derivative.
Here the second label $1$ in $(1,1)$ is the constant test
function and the first is the identity meson matrix, extended
linearly from the chosen basis. The sign comes from the
constant-source identity. The vacuum term may diverge with
$K$; this estimate concerns its difference from the cumulant.
The radii in \eqref{eq:qcdm2-analytic-radii} depend on the
physical volume. No infinite-volume analyticity follows from
this argument.

\subsection{Supports with empty total intersection}

For closed nonempty sets $S_1,\ldots,S_n\subset T_\ell$ define
\[
 d_\cap(S_1,\ldots,S_n)
 =\min_{x\in T_\ell}\max_i\operatorname{dist}(x,S_i).
\]
Compactness shows that $d_\cap>0$ precisely when
$\bigcap_iS_i=\emptyset$. Pairwise intersections are allowed.
A tuple $X=((\alpha_i,f_i))_{i=1}^n$ is admissible in this
section when $n\ge2$, $\alpha_i\in\{0\}\cup\mathfrak B$,
$f_i$ is smooth, and its supports have
empty total intersection. We write $d_\cap(X)$ for this
distance. As usual $\|f\|_0=\|f\|_\sharp$ and
$\|f\|_\Gamma=\|f\|_\infty$ for $\Gamma\in\mathfrak B$;
these basis matrices have operator norm at most one.

\begin{lemma}[Mixed derivatives of the vacuum energy]\label{C-412-53}
Let $s_*+\sum_i z_is_i\in\mathcal P$ for $|z_i|\le r_i$,
and suppose the supports of $s_i$ have $d_\cap=d>0$.
There is a constant $C_{\rm ov}$, independent of $K$, $\ell$
and the sources, such that
\begin{equation}\label{eq:qcdm2-overlap-energy}
 \left|\partial_1\cdots\partial_n
       E_K(s_*+\textstyle\sum_i z_is_i)\big|_{z=0}\right|
 \le 2^{n-1}C_{\rm ov}\ell^3(1+d^{-3})\prod_i r_i^{-1}.
\end{equation}
\end{lemma}

\begin{proof}
Let $P_i$ set $z_i=0$. A mixed derivative kills every term
independent of one variable. It can therefore be estimated by
Cauchy's formula applied to $\prod_i(1-P_i)e_j(Y;s_*+s(z))$.
Explicitly, this finite difference is
\[
 \prod_i(1-P_i)e_j(Y;s_*+s(z))
 =\sum_{A\subset\{1,\ldots,n\}}(-1)^{n-|A|}
        e_j\left(Y;s_*+\sum_{i\in A}z_is_i\right).
\]
The operators $1-P_i$ commute. We may therefore choose one
index $i$ separately for each polymer, apply the source-tail
bound to $(1-P_i)e_j$, and bound the other $n-1$
differences by their $2^{n-1}$ terms. This is why empty
total intersection suffices: the source that is distant
from a small polymer need not be the same for every polymer.

Choose $y\in Y$. Some support $S_i$ has
$\operatorname{dist}(y,S_i)\ge d$. If
$\operatorname{diam}Y\ge d/2$, half the polymer weight pays
for this diameter:
\[
 e^{-\kappa d_j(Y)/2}
 \le e^{\kappa/2}e^{-\kappa d/(4c_DM\lambda_j)}.
\]
Otherwise $\operatorname{dist}(Y,S_i)\ge d/2$, and the source
tail supplies $e^{-\kappa_\sigma d/(2\lambda_j)}$.
Set
\[
 \eta=\min\{\kappa/(4c_DM),\kappa_\sigma/2\}.
\]
The derivative of one energy is at most
\[
 2^{n-1}C_Ee^{\kappa/2}\prod_i r_i^{-1}
 e^{-\eta d/\lambda_j}e^{-\kappa d_j(Y)/2}.
\]
Assigning each polymer to one of its blocks and using the
polymer count bounds the sum over $Y$ by this expression
times $C_{\rm cnt}(\ell/(M\lambda_j))^3$.
Finally
\[
 \sum_{j=1}^K\lambda_j^{-3}e^{-\eta d/\lambda_j}
 \le C_\eta(1+d^{-3}).
\]
Indeed, split at $\lambda_j=d$. The coarse part is a geometric
sum bounded by $Cd^{-3}$, and the fine part is bounded by
$Cd^{-3}\sum_{u\ge0}L^{3u}e^{-cL^u}$; when $d\ge1$ only
the latter estimate is needed. Summing the energy decomposition
proves the assertion.
\end{proof}

Let $k_n=\sum_{\pi}(|\pi|-1)!$, with the sum over partitions
of $\{1,\ldots,n\}$, and let
$\varrho_{\min}=\min(\varrho,\varrho_J)$.

\begin{proposition}[Uniform cumulants at complex parameters]\label{C-412-54}
For every admissible tuple $X$ the lattice cumulant
$S^T_{n,K}(b,m;X)$ is holomorphic on $W_{\mathbb C}$, and
\begin{equation}\label{eq:qcdm2-overlap-cumulant}
 |S^T_{n,K}(b,m;X)|
 \le\left(\frac n{\varrho_{\min}}\right)^n
 \left[k_n(2V_\ell)^n+
       2^{n-1}C_{\rm ov}\ell^3(1+d_\cap(X)^{-3})\right]
 \prod_i\|f_i\|_{\alpha_i}.
\end{equation}
For all $p,q\ge0$,
\begin{equation}\label{eq:qcdm2-parameter-insertion}
 (-\partial_b)^p(-\partial_m)^q S^T_{n,K}(X)
 =S^T_{n+p+q,K}(X,(0,1)^p,(1,1)^q).
\end{equation}
The tuple on the right has the same $d_\cap$.
\end{proposition}

\begin{proof}
Work about a real base point and represent the new parameters
by the constant source $c$ of \cref{C-403-41}. For a zero test
function the assertion follows from multilinearity. Otherwise
choose $r_i=\varrho_{\alpha_i}/(n\|f_i\|_{\alpha_i})$.
The sum of the action-source norms is at most $\varrho$,
and the sum of the meson-source norms is at most $\varrho_J$.
Since $V_\ell\ge1$, the constant displacement has norms
less than $\rho_b\le\varrho/2$ and
$\rho_m\le\varrho_J/2$. In particular
\[
 \left\|c_w+\sum_i z_iw_i\right\|_\sharp
        <\tfrac32\varrho<2\varrho,
 \qquad
 \left\|c_J+\sum_i z_iJ_i\right\|_\infty
        <\tfrac32\varrho_J<2\varrho_J.
\]
Thus the whole Cauchy polydisc stays inside $\mathcal P$.
The triangle inequality gives this source budget even
when the supports overlap.

Set $G(z)=W_K[c+s(z)]/W_K[c]$. The zero-free estimate gives
$G(0)=1$ and $|G(z)|\le2V_\ell$. As germs at zero,
\[
 \log Z_K[c+s(z)]-\log Z_K[c]
 =E_K(c+s(z))-E_K(c)+\log G(z).
\]
The mixed derivative of the last term is
\[
 \sum_\pi(-1)^{|\pi|-1}(|\pi|-1)!
                 \prod_{B\in\pi}\partial^BG(0).
\]
For each block $B$ of a partition, Cauchy's estimate gives
\[
 |\partial^BG(0)|\le2V_\ell\prod_{i\in B}r_i^{-1}.
\]
The product over the blocks contains each inverse radius
once and costs $(2V_\ell)^{|\pi|}\le(2V_\ell)^n$.
Summing the partition coefficients therefore bounds the
mixed logarithmic derivative by
$k_n(2V_\ell)^n\prod_i r_i^{-1}$. This calculation uses
the logarithmic germ at zero; the generating function
may have zeros elsewhere in its source polydisc.
Apply \cref{C-412-53} to the energy and use
$\prod_i r_i^{-1}\le(n/\varrho_{\min})^n
\prod_i\|f_i\|_{\alpha_i}$. This proves
\eqref{eq:qcdm2-overlap-cumulant}.

The finite-dimensional integral is entire in the parameters
and the sources. Its partition function is nonzero at zero
source throughout $W_{\mathbb C}$; local logarithms therefore
give holomorphic cumulants there. Differentiating the
constant-source identity proves
\eqref{eq:qcdm2-parameter-insertion}. Adding an entry supported
on the whole torus adds zero to the maximum defining
$d_\cap$, so it leaves that distance unchanged.
\end{proof}

\begin{theorem}[Holomorphic families of continuum correlations]\label{C-412-55}
Every sequence $K_j\to\infty$ has a subsequence for which
$S^T_{n,K_j}(b,m;X)$ converges for every admissible tuple $X$,
locally uniformly in $(b,m)\in W_{\mathbb C}$ and with all
parameter derivatives. Denote the limit by $S_n(b,m;X)$.
It obeys \eqref{eq:qcdm2-overlap-cumulant} and
\eqref{eq:qcdm2-parameter-insertion} with the cutoff removed.
At real parameters its restriction to separated tuples is a
continuum limit of the original lattice model on $T_\ell$.
For each fixed choice of indices it determines a unique
distribution on
\[
 T_\ell^n\setminus\{x_1=\cdots=x_n\}.
\]
This distribution is symmetric and translation invariant,
has the flavour and cubic covariance of the lattice
correlations, and depends holomorphically on the parameters
when tested on admissible tensors.
\end{theorem}

\begin{proof}
Choose a countable collection of products whose factors are
finite unions of closed dyadic boxes and have empty total
intersection. On each product choose
countable dense sets of smooth test functions, in the $C^1$
norm for an action-density entry and the uniform norm for a
meson entry. There are finitely many basis indices at each
position. The bound \eqref{eq:qcdm2-overlap-cumulant} and
Montel's theorem, followed by a diagonal selection, give
locally uniform convergence on every tuple in these sets.
Use one enumeration containing all orders $n$, all label
strings, all these support tuples, and their dense test
families. A countable compact exhaustion of $W_{\mathbb C}$
makes the local uniform convergence part of this same
selection. Thus one subsequence works for all these data.

For an arbitrary admissible tuple the supports have
$d=d_\cap(X)>0$. Enlarge them by less than $d/4$ and cover
the enlarged sets by sufficiently fine closed dyadic cubes.
The resulting sets still have $d_\cap\ge d/2$. Approximation
within these sets extends the convergence to $X$ uniformly
on every compact parameter set. To see the uniformity,
write $T_K(f_1,\ldots,f_n)=S^T_{n,K}(b,m;X)$ with fixed
labels and these enlarged supports. The bound gives a
common constant $B$ there. Multilinearity gives
\[
 |T_K(f_1,\ldots,f_n)-T_K(h_1,\ldots,h_n)|
 \le B\sum_{i=1}^n\|f_i-h_i\|_{\alpha_i}
       \prod_{r<i}\|h_r\|_{\alpha_r}
       \prod_{r>i}\|f_r\|_{\alpha_r}.
\]
The approximating functions may be chosen with bounded
norms, so the right-hand side tends to zero uniformly
in $K$ and the parameters. Applying the same estimate
to two cutoffs proves convergence for the arbitrary tuple.

For a compact parameter set, choose slightly larger compact
polydiscs still contained in $W_{\mathbb C}$. Cauchy's
formula on their boundaries gives convergence of every
parameter derivative on the smaller sets. The constant
insertions in \eqref{eq:qcdm2-parameter-insertion} also
belong to the common extraction: adding a whole-torus
support preserves $d_\cap>0$. Hence both sides of the
insertion identity converge to the stated identity.

A point outside the total diagonal has two distinct
coordinates. It therefore has a product neighbourhood with
two disjoint closed factors. On this neighbourhood the
estimate defines a continuous multilinear functional on
smooth functions, hence a distribution. These local
distributions agree on overlaps: products of test functions
are dense, and the same limiting functional gives their
values. A partition of unity glues them uniquely.

Permutation, flavour and cubic identities hold at each cutoff
and pass to the limits. Approximate any torus translation by
lattice translations. Translation is continuous in the
test-function norms, so the uniform bound gives the full
translation identity. At real fixed parameters the same
subsequence defines the stated continuum limit. If parameters
also converge inside $W$, local uniform convergence shows that
the limit is the value of this same holomorphic family at the
limiting parameters.
\end{proof}

\begin{corollary}[Mass differentiation]\label{C-412-36}
For every admissible tuple the function $m\mapsto S_n(b,m;X)$
is real analytic on the allowed real interval and
\[
 \partial_m S_n(b,m;X)=-S_{n+1}(b,m;X,(1,1)).
\]
The inserted cumulants are bounded uniformly in the cutoff
at fixed $g$, $\ell$, $n$ and supports. The same holds for
every higher mass derivative and for mixed coupling and mass
derivatives.
\end{corollary}
\begin{proof}
Apply the preceding theorem and the insertion identity. The
uniform bound for the inserted tuple uses its unchanged
$d_\cap>0$.
\end{proof}

\begin{remark}\label{rem:qcdm2-total-diagonal}
The theorem fixes correlations on partial diagonals. It does
not fix an extension to the total diagonal. If two
translation-invariant distributional extensions exist, their
difference is supported there. In relative coordinates it is
a finite sum of derivatives of the delta distribution with
constant coefficients. Thus the remaining ambiguity is local;
it cannot be prescribed independently on a proper partial
diagonal.
\end{remark}

\section{Continuation in the mass}\label{sec:qcdm2-mass-continuation}

At fixed real coupling and fixed volume, one subsequence gives
meromorphic continuations of all separated correlations in
the mass, with a common discrete exceptional set. We prove
that this set misses a neighbourhood of a longer positive
real mass interval. Its width is at least a constant times
$[\ell^3(1+\log\ell)]^{-1}$; see \cref{C-404-60,C-404-72}.
The proof first makes the vacuum normalization independent of
the real mass used to define it. Reflection into small cells
then gives a differential inequality on the real axis, and
a harmonic majorant extends its control to complex masses.
The final two propositions retain the short-distance kernels
and the volume-independent real-mass cumulant bounds.
The continuum tests in this section are $C^1$ at the action
index and continuous at the scalar-bilinear indices; the
finite-cutoff estimates also apply to Lipschitz tests.

\subsection{A common normalization as the mass varies}

We now fix the real coupling \(b\in I_g\) and continue in the
mass alone. Write \(I=[m_-,m_+]\) and
\[
 U_r=\{z\in\mathbb C:\operatorname{dist}(z,I)<r\},
 \qquad r_J=r_M/L.
\]
The source radii satisfy
\[
 2\varrho_J\le \tfrac12r_J,\qquad
 2\varrho_J\le\tfrac12m_-.
\]
The torus, \(g\), and \(b\) remain fixed in every complex
continuation below. All cutoff limits are subsequential.

For a real centre \(x\in I\), denote its source energy by
\(E_K^{(x)}(s)\). Put
\[
 q_z^{(x)}=(0,-(z-x)1),\qquad
 Z_K(z;t)=Z_K^{(x)}\bigl[q_z^{(x)}+\textstyle\sum_i t_is_i\bigr].
\]
Here \(s_i\) is the source of the entry \((\alpha_i,f_i)\).
The right side is independent of the centre: in its determinant
the mass matrix is
\[
 ax-a\bigl(J-(z-x)1\bigr)=az-aJ.
\]
Thus \(Z_K(z;t)\) is the partition function at bare mass \(az\),
with the indicated insertions. It is entire in \(z,t\).
Write
\[
 W_K^{(x)}[s]=e^{-E_K^{(x)}(s)}Z_K^{(x)}[s].
\]
This is the normalized partition function denoted by $W_K$
above. When a centre is fixed, we suppress its superscript.
The energies at two centres need a separate comparison.

Choose $n_\pi\ge1$ with
$(1+1/n_\pi)^{p_0}\le4/3$, and put
\[
 \pi_j=(1+(K-j)/n_\pi)^{p_0},\qquad
 \|w\|_{\sharp,j}
 =\max\{\|w\|_\infty,40M\lambda_j\pi_j\operatorname{Lip}w\}.
\]
Here $p_0$ is the exponent already fixed in \cref{sec:qcdm2-mixed-proof};
no new RG weight is chosen.
We use the source construction of \cite{QCDM1} on
\[
 \mathscr O_j=\{(w,J):\|w\|_{\sharp,j}<4\varrho,\
                         \lambda_j\|J\|_\infty<r_J\},
 \qquad \lambda_j=L^{j-K}.
\]
In particular its local energies are holomorphic on
\(\mathscr O_{j-1}\), with the size and tail bounds recalled
above. The history bound applies on
\[
 \mathcal P^\#_j
 =\{(w,J):\|w\|_\sharp<2\varrho,\
                       \lambda_j\|J\|_\infty\le r_J/2\}.
\]
We will also use the almost-additivity estimate at a base source
and its bilinear improvement from
\cite[Props. 10.15 and 10.18]{QCDM1}.
For bilinear sources \(s_0=(0,J_0)\), \(s_i=(0,J_i)\) with
\[
 \|J_0\|_\infty\le r_J/2,\qquad
 \|J_i\|_\infty\le\epsilon\le r_J/8,
\]
whose supports have mutual distance at least \(\delta\) and
lie in balls of radii \(R_i\), that improvement is
\begin{equation}\label{eq:qcdm2-mass-additivity}
 \left|E_K\left(s_0+\sum_i s_i\right)
       -\sum_iE_K(s_0+s_i)+(N-1)E_K(s_0)\right|
 \le C_\flat\sum_i\left(1+(R_i/\delta)^3\right).
\end{equation}
The constant \(C_\flat\) is independent of the cutoff and side.
At the fine scales this bound uses the separation exponential.
At the coarse scales Schwarz's lemma in the bilinear source
adds a factor \(\lambda_{j-1}/4\). This makes the sum independent
of the number of coarse scales. That improvement is needed
when the cells below become small.

\begin{proposition}[Compatibility of the vacuum energies]
\label{C-404-66}
Let \(x,x'\in I\), \(\zeta=x'-x\), and
\(s_\zeta=(w,J-\zeta1)\) for \(s=(w,J)\).
For \(1\le k\le K\), if \(s,s_\zeta\in\mathscr O_k\), then
\[
 E_k^{(x)}(s_\zeta)+\mathfrak f_k^{(x)}
       =E_k^{(x')}(s)+\mathfrak f_k^{(x')}.
\]
Here \(\mathfrak f_k^{(x)}\) is the real finite-size constant
in the empty-history scalar normalization of \cite[Lem. 9.15]{QCDM1};
\(|\mathfrak f_k^{(x)}|\le C_{\rm fs}\).
\end{proposition}

\begin{proof}
The two bare densities agree by the mass-source identity.
Their block integration also agrees before it is expressed
in terms of extracted coefficients. Indeed the fermion block
kernel reads only \(b_F\) and the block averages, while the gauge
chart integral and the density test read the calibrations
\(\mathsf t_j=L^{K-j}b\). None of these depends on the mass
centre or on a coefficient extracted by the flow.

More explicitly, if \(d_j\) is the density of the empty-history
term at level \(j\), the next density is
\[
 d_{j+1}(W)=\chi_{j+1}(W)\,
       \mathsf I_j\!\left[\operatorname{Ber}_j
                    (\operatorname{Ker}_j d_j)\right](W).
\]
Here $\chi_{j+1}$ is the density-test indicator,
$\mathsf I_j$ integrates the gauge chart, $\operatorname{Ber}_j$
integrates the old Grassmann variables, and
$\operatorname{Ker}_j$ is the fermion block kernel.
The density is zero when \(W\) fails the test. This is the exact
tested chart and block-integration identity of
\cite[Lem. 7.1 and Prop. 10.1]{QCDM1}. The two inputs have \(d_0=d'_0\);
linearity and the identical kernel, test and chart give
\(d_j=d'_j\) at every subsequent level.

We compare the scalar coefficients at the trivial block field.
In the extracted form of the empty-history density all
fermion quadratic terms and activities with legs have positive
Grassmann degree. Thus its degree-zero coefficient, at a
tested block field \(V\), is
\begin{equation}\label{eq:qcdm2-mass-scalar-density}
 g_0(V)=\exp\left[
 E_k^{(x)}(s_\zeta)-A_k(U;s_\zeta)
       +\sum_{j\le k}\sum_X
          E^{(j)}_0(X;\mathbf p_X^{(j)}(U);s_\zeta)
                         \right],
\end{equation}
where \(U\) is any background in the minimal orbit over \(V\).
Gauge invariance makes this expression independent of that
choice. The analogous formula defines \(g'_0\).

The tested block fields form an open set containing \(1\).
On a neighbourhood of \(1\), the canonical critical-background
theorem of \cite[Thm. E]{YMM1} provides a continuous representative
\(V\mapsto U^\circ(V)\) of the minimal orbit. To apply it,
write \(V=e^{iB_V}\) bond by bond using the principal logarithm.
Write $r_{\rm B}$ for the positive logarithmic row radius
in that theorem. For \(\|V-1\|<\min\{1,2r_{\rm B}/\pi\}\),
\(\|B_V\|\le(\pi/2)\|V-1\|<r_{\rm B}\).
The logarithm is continuous and the representative is
holomorphic in \(B_V\) on this ball. Their composition has
the required continuity. The geometric hypotheses are those
already verified in the tested-block-field lemma of \cite{QCDM1}.
Its restricted pairs depend continuously
on \(V\); the scalar activities are holomorphic on an open
pair domain containing them. The action and the currents
entering those pairs are continuous functions of the link
matrices. There are only finitely many polymers at fixed
cutoff and volume. Therefore \eqref{eq:qcdm2-mass-scalar-density}
is continuous near \(1\), and so is \(g'_0\).

The exact density identity gives equal integrals of \(g_0\)
and \(g'_0\) over every Borel subset of the block-field space.
The reference measure there is Haar measure and charges every
nonempty open set. Two continuous densities agreeing almost
everywhere must consequently agree throughout this neighbourhood.
At \(V=1\) we may take \(U=1\). The action vanishes, and the
scalar normalization in the source step gives
\[
 g_0(1)=\exp\bigl(E_k^{(x)}(s_\zeta)+\mathfrak f_k^{(x)}\bigr),
 \qquad
 g'_0(1)=\exp\bigl(E_k^{(x')}(s)+\mathfrak f_k^{(x')}\bigr).
\]
Their equality determines the difference of the exponents
modulo \(2\pi i\mathbb Z\).

The common domain
\(\{s:s,s_\zeta\in\mathscr O_k\}\) is convex.
The difference of the two exponents is continuous there and
takes values in \(2\pi i\mathbb Z\), hence is constant.
If the domain is nonempty, its mass inequalities imply
\(|\zeta|<2r_J/\lambda_k\); it therefore contains
\(s_*=(0,\zeta1/2)\). The two sources at this point are
\((0,\zeta1/2)\) and \((0,-\zeta1/2)\).
They are fixed by source conjugation, so both energies and
both finite-size constants are real. Their difference is
therefore zero. This fixes the branch and proves the identity.
\end{proof}

\begin{proposition}[A common normalizer]\label{C-404-67}
For every \(K\ge1\) there is a holomorphic function
\[
 \mathcal E_K(z)
   =E_K^{(x)}(q_z^{(x)})+\mathfrak f_K^{(x)},
 \qquad |z-x|<r_J,\quad x\in I,
\]
on \(U_{r_J}\). It satisfies
\(\mathcal E_K(\bar z)=\overline{\mathcal E_K(z)}\).
Set
\[
 M_0=\exp\!\left(\max\{C+2,C_P(1+\log X_g)\}\ell^3\right),
 \qquad h_K(z)=Z_K(z)e^{-\mathcal E_K(z)}.
\]
Then
\begin{equation}\label{eq:qcdm2-common-normalizer}
 |h_K(z)|\le M_0e^{C_{\rm fs}}
       \quad(\operatorname{dist}(z,I)\le2\varrho_J),\qquad
 h_K(x)\ge(M_0e^{C_{\rm fs}})^{-1}\quad(x\in I).
\end{equation}
For \(0<r<2\varrho_J\), put
\[
 \rho_{n,r}=\min\{\varrho/n,2\varrho_J-r\},\qquad
 M_{n,d}=M_0\exp\!\left[nC_{\rm count}(1+d^{-3})\ell^3\right].
\]
Here \(C_{\rm count}\) is the torus source-count constant.
If \(n\ge2\), the supports of \(X=((\alpha_i,f_i))_{i=1}^n\)
are separated by at least \(d\), and
\(\|f_i\|_{\alpha_i}\le\rho_{n,r}\), then there are holomorphic
functions \(\mathcal E_{K,i}(z,t_i)\) on
\(U_r\times\{|t_i|<1\}\), zero at \(t_i=0\), such that
\begin{equation}\label{eq:qcdm2-normalized-insertion-bound}
 \left|Z_K(z;t)
  e^{-\mathcal E_K(z)-\sum_i\mathcal E_{K,i}(z,t_i)}\right|
       \le M_{n,d}e^{C_{\rm fs}} .
\end{equation}
\end{proposition}

\begin{proof}
If \(|z-x|,|z-x'|<r_J\), apply \cref{C-404-66} to
\(s=q_z^{(x')}\). Its shifted source is \(q_z^{(x)}\).
Thus the displayed definitions of \(\mathcal E_K\) agree.
Each is holomorphic on its disc, and the discs cover \(U_{r_J}\).
Source conjugation gives the stated reality identity.
The upper bound follows from \eqref{eq:qcdm2-partition-window}
at \(q_z^{(x)}\), and the lower bound at \(x\in I\)
follows from the zero-source lower bound at that centre.
The factor \(e^{C_{\rm fs}}\) accounts for
\(\mathfrak f_K^{(x)}\).

Define locally
\[
 \mathcal E_{K,i}(z,t_i)
   =E_K^{(x)}(q_z^{(x)}+t_is_i)-E_K^{(x)}(q_z^{(x)}).
\]
For \(|z-x|<r\), \(|t_i|<1\), every partial sum of these
sources lies in the interior of \(\mathcal P\): the action
norm is less than \(n\rho_{n,r}\le\varrho\), and disjointness
of the supports gives the matrix-source norm
\[
 |z-x|+\max_i|t_i|\,\|f_i\|_\infty<2\varrho_J.
\]
Apply \cref{C-404-66} with and without \(t_is_i\), and
subtract. The finite-size constants cancel, proving that
\(\mathcal E_{K,i}\) is independent of the chosen centre.

Almost additivity at the base \(q_z^{(x)}\) gives
\[
 \left|E_K^{(x)}\left(q_z^{(x)}+\sum_i t_is_i\right)
       -E_K^{(x)}(q_z^{(x)})
       -\sum_i\mathcal E_{K,i}(z,t_i)\right|
 \le\sum_i\mathscr C_K(\operatorname{supp}s_i;d)
 \le nC_{\rm count}(1+d^{-3})\ell^3.
\]
Combining this with the partition upper bound proves
\eqref{eq:qcdm2-normalized-insertion-bound}.
\end{proof}

\subsection{A lower bound on a longer real mass interval}

The common normalizer controls possible zeros uniformly on
compact subsets of \(U_{2\varrho_J}\). To obtain a useful
zero-free neighbourhood we first bound the normalized partition
function on a longer real interval.

Fix \(0<\epsilon\le\varrho_J/2\), and set
\begin{align*}
 R_\epsilon&=2\varrho_J-2\epsilon,&
 c_g&=C_P(1+\log X_g),\\
 c_\epsilon&=512\epsilon^{-1}L^3e^{1+19C_\flat},&
 A_\epsilon&=(C+3+c_g)e^{c_\epsilon R_\epsilon},\\
 K_\epsilon&=1+\left\lceil\tfrac13\log_L A_\epsilon\right\rceil .
\end{align*}
These constants are independent of the side and cutoff.

\begin{proposition}[The normalized partition function at real masses]
\label{C-404-70}
Let \(x_0\in I\), \(K\ge K_\epsilon\), and
\[
 h_{K,x_0}(x)=Z_K(x)
          e^{-E_K^{(x_0)}(q_x^{(x_0)})}.
\]
For real \(|x-x_0|\le R_\epsilon\),
\begin{equation}\label{eq:qcdm2-real-mass-floor}
 h_{K,x_0}(x)\ge e^{-A_\epsilon\ell^3},\qquad
 |(\log h_{K,x_0})'(x)|\le c_\epsilon A_\epsilon\ell^3.
\end{equation}
\end{proposition}

\begin{proof}
Write \(E=E_K^{(x_0)}\), \(q_x=q_x^{(x_0)}\), \(h=h_{K,x_0}\),
and \(u(x)=-\ell^{-3}\log h(x)\).
All these real masses exceed \(m_-/2\). The positive
real-mass density gives $Z_K(x)>0$, and \(E(q_x)\) is real,
so \(h>0\) and \(u\) is
real analytic.

For \(1\le i\le K-1\), divide the torus into cells of
physical side \(\lambda=L^{-i}\). Their lattice side
\(\lambda_{\rm l}=L^{K-i}\) is odd and at least \(L\);
the number of cells in each direction is even.
Let \(N_{\rm c}=(\ell/\lambda)^3\).
In a cell centred at \(y^c\), in lattice coordinates, put
\[
 \eta_c(y)=\prod_{\nu=1}^3\tau(y_\nu-y^c_\nu),\qquad
 \tau(v)=\min\left\{1,\max\left\{0,
          \frac{3\lambda_{\rm l}/8-|v|}{\lambda_{\rm l}/4}
                                      \right\}\right\}.
\]
These functions are carried to one another by wall reflections.
Their supports have mutual distance at least \(\lambda/4\),
and each lies in a ball of radius \(3\sqrt3\lambda/8\).
Consequently \(1+(R/\delta)^3\le19\).

For \(|t|\le\epsilon\), apply the chessboard inequality to
\(e^{\Phi(q_x+t\eta_c1)}\) in one cell and to
\(e^{\Phi(q_x)}\) in the others, with each constant source
restricted to its own cell. Since \(q_x\) is real and scalar,
its fold is the same constant source. Thus
\[
 |Z[q_x+t\eta_c1]|
 \le Z[q_x+(t\eta)^{\rm ch}]^{1/N_{\rm c}}
                         Z[q_x]^{1-1/N_{\rm c}},
\]
where the coefficient in a reflected cell is \(t\) or \(\bar t\).
The folded source has norm at most
\(R_\epsilon+\epsilon<2\varrho_J\).
Use \eqref{eq:qcdm2-mass-additivity} and the partition upper
bound. With
\[
 \Lambda(t)=N_{\rm c}^{-1}\sum_{c'}
           \bigl[E(q_x+t\eta_{c'}1)-E(q_x)\bigr],
\]
source conjugation identifies the real parts of the odd and
even reflected terms. We obtain
\[
 \psi_c(t)=e^{-\Lambda(t)}\frac{Z[q_x+t\eta_c1]}{Z[q_x]},
 \qquad \psi_c(0)=1,
\]
\[
 |\psi_c(t)|
 \le \exp\left[19C_\flat+(C+2+u(x))\lambda^3\right]
       =:B_1(x,\lambda).
\]
These functions are holomorphic past \(|t|=\epsilon\),
so \(|\psi'_c(0)|\le B_1/\epsilon\).

For a site function \(v\), define the linear functional
\[
 \mathcal V(v)=
 \left.\frac{d}{dt}\log W_K[q_x+tv1]\right|_{t=0}.
\]
For one offset \(o\) of the board,
\(\sum_c\psi'_c(0)=\mathcal V(\eta^{(o)})\), where
\(\eta^{(o)}=\sum_c\eta_c\). Hence
\(|\mathcal V(\eta^{(o)})|\le N_{\rm c}B_1/\epsilon\).
Summing over the \(\lambda_{\rm l}^3\) offsets gives
\[
 \sum_o\eta^{(o)}=\sigma^3,\qquad
 \sigma=\sum_{v\in\mathbb Z}\tau(v)\ge\lambda_{\rm l}/8.
\]
Since \(q_x+(0,t1)=q_{x-t}\), this yields
\[
 |u'(x)|\le512\epsilon^{-1}\lambda^{-3}
       \exp\left[19C_\flat+(C+2+u(x))\lambda^3\right].
\]

Set \(Y=C+2+\max\{u,1\}\). At a point where
\(Y\le A_\epsilon\), choose
\(i=\lceil\frac13\log_L Y\rceil\). Then
\[
 1\le i\le K-1,\qquad
 \lambda^3Y\le1<L^3\lambda^3Y.
\]
The preceding bound gives \(|u'|\le c_\epsilon Y\) and
\(|Y'|\le c_\epsilon Y\) almost everywhere.
At \(x_0\), the zero-source floor gives \(u(x_0)\le c_g\),
so \(Y(x_0)\le C+3+c_g\).
Integrate the differential inequality on the maximal interval
where \(Y\le A_\epsilon\). If its right endpoint were less
than \(x_0+R_\epsilon\), then
\[
 Y(x)\le(C+3+c_g)e^{c_\epsilon(x-x_0)}<A_\epsilon
\]
at that endpoint, contradicting maximality and continuity.
The same argument to the left gives \(Y\le A_\epsilon\)
throughout \(|x-x_0|\le R_\epsilon\).
The two assertions now follow from \(u\le A_\epsilon\)
and \(|(\log h)'|=\ell^3|u'|\).
\end{proof}

\begin{lemma}[A complex neighbourhood from a real derivative bound]
\label{C-404-71}
Let \(h\) be holomorphic near \(\{|z-x|\le r\}\), where
\(x\) is real, and suppose
\[
 h(\bar z)=\overline{h(z)},\qquad h>0
       \text{ on }[x-r,x+r],
\]
\[
 \log|h(z)/h(x)|\le W,\qquad
 \log(h(x')/h(x))\le\Lambda|x'-x|
       \quad(x'\in[x-r,x+r]).
\]
For \(z=x+\xi+i\eta\), \(|z-x|\le r/2\), \(\eta\ne0\),
\begin{equation}\label{eq:qcdm2-half-disc-bound}
 \log|h(z)/h(x)|
 \le\frac{16W}{3\pi}\frac{|\eta|}{r}
       +\Lambda|\xi|
       +\frac2\pi\Lambda|\eta|\log\frac r{|\eta|}.
\end{equation}
Consequently, if \(0<\rho\le r/e\) and
\[
 \rho\left[\frac{16W}{3\pi r}+\Lambda
                  +\frac2\pi\Lambda\log(r/\rho)\right]\le1,
\]
then \(|h(z)|\ge h(x)/2\) for
\(|z-x|\le\rho/(2(e+1))\).
\end{lemma}

\begin{proof}
By conjugation it suffices to consider \(\eta>0\).
On the upper unit half-disc put
\[
 \zeta=(z-x)/r,\qquad
 u(\zeta)=\log|h(x+r\zeta)/h(x)|.
\]
This function is subharmonic and bounded above, including
when \(h\) has zeros. The harmonic measure of the semicircle is
\[
 \omega(\zeta)=\frac2\pi
       \arctan\frac{2\operatorname{Im}\zeta}{1-|\zeta|^2}.
\]
The fractional linear map \((1+\zeta)/(1-\zeta)\) sends the
half-disc to the first quadrant and its semicircle to the
positive imaginary axis, giving this formula. For
\(|\zeta|\le1/2\),
\(\omega(\zeta)\le16\operatorname{Im}\zeta/(3\pi)\).
The second harmonic function
\[
 q(\zeta)=\operatorname{Re}
       \left(\zeta+\frac{2i}{\pi}\zeta\log\zeta\right)
 =\operatorname{Re}\zeta
       \left(1-\frac2\pi\arg\zeta\right)
       +\frac2\pi\operatorname{Im}\zeta\log\frac1{|\zeta|}
\]
uses the principal logarithm. It extends continuously to the
diameter with value \(|s|\), and on the semicircle its value
is \(\cos t(1-2t/\pi)\ge0\).
The boundary assumptions therefore give
\(u\le W\omega+\Lambda r q\) on the half-disc by the
maximum principle. The two endpoints cause no difficulty
because \(u\) is bounded above. Bound
\(|1-2\arg\zeta/\pi|\le1\) and
\(\log(1/|\zeta|)\le\log(1/\operatorname{Im}\zeta)\).
This proves \eqref{eq:qcdm2-half-disc-bound}.

Since \(y\log(r/y)\) increases for \(0<y\le r/e\),
the hypothesis on \(\rho\) gives
\(|h|\le e\,h(x)\) on the disc of radius \(\rho\).
Schwarz's lemma applied to \(h-h(x)\), bounded there by
\((e+1)h(x)\), gives
\[
 |h(z)-h(x)|\le(e+1)h(x)|z-x|/\rho.
\]
The asserted lower bound follows.
\end{proof}

Set
\[
 I_\epsilon=[m_--2\varrho_J+3\epsilon,\
                         m_++2\varrho_J-3\epsilon],
 \quad W_\ell=(C+2+A_\epsilon)\ell^3,\quad
 \Lambda_\ell=c_\epsilon A_\epsilon\ell^3.
\]
Let \(\rho_\ell\) be the largest number in \((0,\epsilon/e]\)
such that
\begin{equation}\label{eq:qcdm2-mass-width}
 \rho_\ell\left[\frac{16W_\ell}{3\pi\epsilon}
       +\Lambda_\ell+
          \frac2\pi\Lambda_\ell\log(\epsilon/\rho_\ell)\right]
       \le1,
 \qquad w_\ell=\frac{\rho_\ell}{2(e+1)}.
\end{equation}
The left side is increasing on this interval and tends to zero
at zero. Thus \(\rho_\ell\) exists. Define
\[
 \Omega_\epsilon=\{z:\operatorname{dist}(z,I_\epsilon)<w_\ell\},
 \qquad \rho_n^\epsilon=\min\{\varrho/n,2\epsilon\}.
\]
All real points of \(\Omega_\epsilon\) exceed \(m_-/2\).

\begin{proposition}[A zero-free neighbourhood of the longer interval]
\label{C-404-72}
For \(K\ge K_\epsilon\) and \(z\in\Omega_\epsilon\),
\begin{equation}\label{eq:qcdm2-long-zero-free}
 |Z_K(z)|\ge\tfrac12e^{-A_\epsilon\ell^3-C_{\rm fs}}
                         |e^{\mathcal E_K(z)}|>0.
\end{equation}
If \(X\) has \(n\ge2\) entries with pairwise separation
at least \(d>0\), then
\begin{equation}\label{eq:qcdm2-long-cumulants}
 |S^T_{n,K}(z;X)|
 \le\left[\frac{2(2e^{A_\epsilon\ell^3}M_{n,d}+1)}
                         {\rho_n^\epsilon}\right]^n
                           \prod_i\|f_i\|_{\alpha_i}.
\end{equation}
There is \(c_w>0\), depending on the fixed constants, \(g\)
and \(\epsilon\), but not on the side, such that
\[
 w_\ell\ge\frac{c_w}{\ell^3(1+\log\ell)}.
\]
\end{proposition}

\begin{proof}
For \(z\in\Omega_\epsilon\), choose \(x\in I_\epsilon\)
with \(|z-x|<w_\ell\), and let \(x_0\in I\) be nearest to
\(x\). Then
\[
 |x-x_0|\le R_\epsilon-\epsilon,\qquad
 \{|\zeta-x|\le\epsilon\}\subset
                        \{|\zeta-x_0|\le R_\epsilon\}.
\]
On this disc \(h_{K,x_0}\) is holomorphic,
\(|h_{K,x_0}|\le e^{(C+2)\ell^3}\), and
\(h_{K,x_0}(x)\ge e^{-A_\epsilon\ell^3}\).
On its real diameter \cref{C-404-70} gives the logarithmic
derivative bound \(\Lambda_\ell\).
Apply \cref{C-404-71} with \(r=\epsilon\),
\(W=W_\ell\), and \(\Lambda=\Lambda_\ell\). It gives
\[
 |h_{K,x_0}(z)|\ge\tfrac12e^{-A_\epsilon\ell^3}.
\]
Since
\(\mathcal E_K(z)=E_K^{(x_0)}(q_z^{(x_0)})
                       +\mathfrak f_K^{(x_0)}\),
this is \eqref{eq:qcdm2-long-zero-free}.

Also \(|z-x_0|<R_\epsilon\). For
\(\|f_i\|_{\alpha_i}\le\rho_n^\epsilon\), the local insertion
bound used in \cref{C-404-67} applies at this centre with
radius \(R_\epsilon=2\varrho_J-2\epsilon\).
The holomorphic function
\[
 \Psi(t)=
 \frac{Z_K(z;t)
       e^{-E_K^{(x_0)}(q_z^{(x_0)})
              -\sum_i\mathcal E_{K,i}(z,t_i)}}
      {h_{K,x_0}(z)}
\]
has \(\Psi(0)=1\) and
\(|\Psi|\le2e^{A_\epsilon\ell^3}M_{n,d}\) on
\(\max_i|t_i|<1\).
On each line through zero, Schwarz's lemma bounds
\(|\Psi(t)-1|\) by
\((2e^{A_\epsilon\ell^3}M_{n,d}+1)\max_i|t_i|\).
On the polydisc of radius
\([2(2e^{A_\epsilon\ell^3}M_{n,d}+1)]^{-1}\),
its logarithm has modulus at most \(\log2\).
For \(n\ge2\), its full mixed derivative equals that of
\(\log Z_K(z;t)\): each subtracted insertion energy depends
on only one \(t_i\). Cauchy's inequality and multilinearity
give \eqref{eq:qcdm2-long-cumulants}.

For the width, put
\[
 \alpha=\frac{16(C+2+A_\epsilon)}{3\pi\epsilon}
                         +c_\epsilon A_\epsilon,\quad
 \beta=\frac{2c_\epsilon A_\epsilon}{\pi},\quad
 L_\ell=\max\{1,\log(2\epsilon\ell^3(\alpha+\beta))\}.
\]
The defining expression in \eqref{eq:qcdm2-mass-width} is
\(\rho\ell^3[\alpha+\beta\log(\epsilon/\rho)]\).
For
\(\rho_*=[2\ell^3(\alpha+2\beta L_\ell)]^{-1}\),
\[
 \log(\epsilon/\rho_*)\le
 L_\ell+\log(2L_\ell)\le2L_\ell.
\]
Thus its value at \(\rho_*\) is at most one when
\(\rho_*\le\epsilon/e\); if \(\rho_*>\epsilon/e\),
its value at \(\epsilon/e\) is less than \(1/2\).
It follows that
\[
 \rho_\ell\ge\min\{\epsilon/e,\rho_*\}.
\]
For \(\ell\ge1\),
\[
 L_\ell\le a_w+3\log\ell
       \le3a_w(1+\log\ell),\qquad
 a_w=1+|\log(2\epsilon(\alpha+\beta))|.
\]
One may therefore take
\[
 c_w=\frac1{2(e+1)}
       \min\{\epsilon/e,[2(\alpha+6\beta a_w)]^{-1}\}.
\]
\end{proof}

\subsection{A common meromorphic family}

\begin{theorem}[Continuation in the mass]
\label{C-404-60}\label{C-404-73}
Fix \(b\in I_g\) and \(0<\epsilon\le\varrho_J/2\).
Every sequence \(K_j\to\infty\) has a subsequence and a
holomorphic function \(h\) on \(U_{2\varrho_J}\), not
identically zero, with the following properties.

For every separated tuple \(X\) of order \(n\ge2\),
\(S^T_{n,K_j}(z;X)\) converges locally uniformly off the
discrete set \(\mathcal Z=\{h=0\}\) to a meromorphic
function \(\mathcal S_n(z;X)\) on \(U_{2\varrho_J}\), and
\(h^n\mathcal S_n(\cdot;X)\) is holomorphic there.
The same \(h\), subsequence and exceptional set work for all
orders, labels and tuples. Moreover
\[
 h_{K_j}\longrightarrow h
        \quad\text{locally uniformly on }U_{2\varrho_J}.
\]
The set \(\mathcal Z\) is symmetric under conjugation and
is the set of accumulation points, in this open domain,
of zeros of \(Z_{K_j}\).

The set \(\mathcal Z\) is disjoint from \(\Omega_\epsilon\).
On \(\Omega_\epsilon\) the limits are holomorphic, obey
\eqref{eq:qcdm2-long-cumulants}, and converge with all mass
derivatives. If \(z_j\to z\in\Omega_\epsilon\), then
\[
 S^T_{n,K_j}(z_j;X)\longrightarrow\mathcal S_n(z;X).
\]
At \(m\in I\) these values are continuum limits with labels
\((b,m)\). Every continuum limit at fixed \(b\), obtained
with masses \(m_j\to m\in I\), occurs this way after a
subsequence. On all real masses in \(\Omega_\epsilon\)
the centered limiting functional is reflection positive on
the separated polynomial words used in its construction.
It has the permutation, translation, cubic, flavour and
charge-conjugation symmetries of the finite lattice theory.

Two such meromorphic families which agree on a subset of
\(I\) having an accumulation point agree throughout
\(U_{2\varrho_J}\). This is a uniqueness assertion for
analytic continuation from those values; it does not assert
that distinct cutoff subsequences have the same limits.
\end{theorem}

\begin{proof}
The functions \(h_K\) are uniformly bounded on \(U_{2\varrho_J}\)
by \eqref{eq:qcdm2-common-normalizer}. We construct bounded
holomorphic numerators for the cumulants on the same domain.
Fix \(0<r<2\varrho_J\), a separated tuple \(X\), and first
rescale its entries so that
\(\|f_i\|_{\alpha_i}\le\rho_{n,r}\).
Put
\[
 \Psi_K(z,t)=Z_K(z;t)
       e^{-\mathcal E_K(z)-\sum_i\mathcal E_{K,i}(z,t_i)}.
\]
It is bounded by \(M_{n,d}e^{C_{\rm fs}}\) on
\(U_r\times\{\max_i|t_i|<1\}\), and
\(\Psi_K(z,0)=h_K(z)\).
Cauchy's formula on smaller circles, followed by their radius
limit, gives the same bound for
\(\partial_B\Psi_K(z,0)\), \(B\subset\{1,\ldots,n\}\).
The logarithmic derivative formula used in
\cref{C-412-54} gives, where \(h_K\ne0\),
\[
 S^T_{n,K}(z;X)=\frac{N_K(z;X)}{h_K(z)^n},
\]
\[
 N_K(z;X)=
 \sum_\pi(-1)^{|\pi|-1}(|\pi|-1)!\,
 h_K(z)^{n-|\pi|}
             \prod_{B\in\pi}\partial_B\Psi_K(z,0).
\]
The one-variable insertion energies have zero full mixed
derivative because \(n\ge2\).
Each numerator is holomorphic. The formulas obtained on
different \(U_r\)'s agree away from the discrete zeros of
\(Z_K\), hence everywhere by continuity. They therefore
define one holomorphic \(N_K(\cdot;X)\) on \(U_{2\varrho_J}\).
For the original entries its bound on \(U_r\) is
\[
 |N_K(z;X)|\le
 k_n\left(\frac{M_{n,d}e^{C_{\rm fs}}}{\rho_{n,r}}\right)^n
                     \prod_i\|f_i\|_{\alpha_i}.
\]

Apply Montel's theorem to \(h_K\) and these numerators on a
countable compact exhaustion of \(U_{2\varrho_J}\).
Use one enumeration of all orders, basis labels and dense
smooth test families with supports in fixed separated closed
sets. The same diagonal selection gives convergence for
every member of this enumeration. The displayed multilinear
bound extends convergence to arbitrary separated tuples:
enlarge their supports slightly while keeping positive
separation, approximate in the action \(C^1\) norm and meson
uniform norm, and telescope the difference of the multilinear
forms one entry at a time. The resulting error tends to zero
uniformly on each compact mass set and uniformly in \(K\).

Denote the limits by \(h,N\).
On \(I\), \eqref{eq:qcdm2-common-normalizer} gives
\(h\ge(M_0e^{C_{\rm fs}})^{-1}\), so \(h\not\equiv0\).
Its zeros are discrete. On a compact set avoiding them,
\(h_{K_j}\) is eventually bounded away from zero, and hence
\[
 S^T_{n,K_j}\longrightarrow \mathcal S_n=N/h^n
\]
locally uniformly there. Source conjugation makes \(h_K\)
real under conjugation, and the same is true of \(h\).
If \(h(z)\ne0\), uniform convergence excludes nearby zeros
of \(h_{K_j}\) for all sufficiently large \(j\).
Conversely, on a small circle around a zero of \(h\) and
containing no other zero, Rouch\'e's theorem gives zeros of
\(h_{K_j}\) inside for every sufficiently large \(j\).
Since \(h_K\) and \(Z_K\) have the same zeros, this proves
the assertion about \(\mathcal Z\).

By \eqref{eq:qcdm2-long-zero-free},
\[
 |h_K(z)|\ge\tfrac12e^{-A_\epsilon\ell^3-C_{\rm fs}}
                 \quad(z\in\Omega_\epsilon,\ K\ge K_\epsilon).
\]
Thus \(h\) has no zero there.
The cumulant bound passes to the limit; Cauchy's formula
on slightly larger compact sets gives convergence of all
mass derivatives. Local uniform convergence and continuity
give convergence at moving masses. In particular a continuum
limit already obtained at \(m_j\to m\in I\) is preserved
when this further extraction is made.

For reflection positivity no convergence of one-point
functions is required. At a real positive mass, center each
lattice insertion by its own mean. Every centered moment
is the sum, over partitions into blocks of size at least two,
of products of cumulants. It therefore converges by the
preceding result. On a positive-time polynomial the finite
lattice reflection form is nonnegative, where each monomial
has internally separated supports. Approximate a
coordinate reflection plane by a bond plane at distance at
most \(a\); the reflected tests differ by a translation
of at most \(2a\). For smooth tests supported away from the
plane, their norms converge and all separations in the
reflection form stay bounded below. The uniform cumulant
bound and multilinearity pass the inequality to the limit.
This proves reflection positivity of the centered functional
at every real mass in \(\Omega_\epsilon\).

The lattice permutation and internal identities hold for
the same finite integrals at these masses and pass to the
limits. Lattice translations approximate all torus translations
in the test norms, as in \cref{C-412-55}; cubic covariance
passes directly. Equivalently, the finite-cutoff identities
are identities of rational functions of the mass and hold
throughout their common domain of definition.
Finally the identity theorem, applied after clearing
denominators, proves the last assertion.
\end{proof}

\subsection{Correlation estimates at the continued masses}

The continuation preserves the two leading short-distance
kernels. We first estimate the change of the source energy.
The estimate uses Schwarz's lemma in the constant mass source;
its scale factor is needed before the sum over levels.

Use the test pairs of \eqref{eq:qcdm2-current-test-pair}, with
\(0<d\le1/200\), and put
\[
 \mathcal K_q(f,h)=\iint f(x)h(y)|x-y|^{-q}\,dx\,dy,\qquad
 H_d=1+\log_L(1/d).
\]
Thus
\begin{equation}\label{eq:qcdm2-continued-template-floor}
 \mathcal K_q(f,h)\ge
       10^{-q}\vartheta^2d^{6-q}\|f\|_\infty\|h\|_\infty.
\end{equation}
Only action-density and scalar-bilinear indices are used below.

\begin{lemma}[Changing the mass in the mixed energy]
\label{C-404-74}
Fix a centre \(x_0\in I\) and \(|\zeta|<2\varrho_J\), and set
\(q[\zeta]=(0,-\zeta1)\).
Let \((F_\nu,H_\nu)_{\nu=1}^N\) be translates and coordinate
mirror images of one test pair, including the induced action
on its labels. Suppose the pairs have mutual distance at least
\(d\) and \(N\ge(\ell/(96d))^3\). Write
\[
 Q_K(\zeta)=\frac1N\sum_{\nu=1}^N
 \left.\partial_{t_1}\partial_{t_2}
 E_K^{(x_0)}(q[\zeta]+t_1F_\nu+t_2H_\nu)\right|_{t=0}.
\]
There are constants \(A_B,A_0\), independent of \(K,\ell,b,x_0\),
such that
\[
 |Q_K(\zeta)-Q_K(0)|
 \le A_B|\zeta|d^3H_d\|f\|_\infty\|h\|_\infty
\]
at two bilinear indices, and
\[
 |Q_K(\zeta)-Q_K(0)|
 \le A_0\left(\frac{40M}{\vartheta\varrho}\right)^2
 |\zeta|dH_d^{1+2p_0}
                     \|f\|_\infty\|h\|_\infty
\]
at two action-density indices. Here \(p_0\) is the
exponent in the RG weight
\(\phi(t)=t^{-1/2}(1+\log t)^{p_0}\).
\end{lemma}

\begin{proof}
Write \(E=E_K^{(x_0)}\) and
\[
 e_j^\zeta(X;s)=e_j(X;q[\zeta]+s)-e_j(X;s).
\]
The desired difference is the sum over \(j,X\) of the mean
mixed derivatives of \(e_j^\zeta\).
For a fixed level let \(f_1=f,f_2=h\), and take source radii
\[
 r_i=\frac{r_J}{2\lambda_{j-1}\|f_i\|_\infty}
       \quad\hbox{at a bilinear index},\qquad
 r_i=\frac{2\varrho}{\|f_i\|_{\sharp,j-1}}
       \quad\hbox{at the action index}.
\]
For sources \(F,H\) with the disjoint supports and norms of
the given pair, \(s(t)=t_1F+t_2H\) lies in the level source
domain when \(|t_i|<r_i\). Indeed disjoint supports make the
norm of the sum the maximum of the two norms. Adding
\(q[\zeta']\), with
\(|\zeta'|<r_J/(2\lambda_{j-1})\), still gives
\(\lambda_{j-1}\|J\|_\infty<r_J\) and
\(\|w\|_{\sharp,j-1}<2\varrho\).

If \(s,s'\) are two such sources, the function
\[
 \zeta'\longmapsto e_j^{\zeta'}(X;s)-e_j^{\zeta'}(X;s')
\]
vanishes at zero. The local source-tail estimate bounds it
on this mass disc by
\[
 2C_Ee^{-\kappa d_j(X)}
   \exp\left[-\kappa_\sigma
             \operatorname{dist}(X,\operatorname{supp}(s-s'))/L^j\right].
\]
Schwarz's lemma therefore multiplies this bound by
\(2\lambda_{j-1}|\zeta|/r_J\).
Apply the estimate with \(s'\) obtained by setting either
\(t_1\) or \(t_2\) equal to zero. The finite-difference
formula and Cauchy's estimate on the two source circles give
\begin{equation}\label{eq:qcdm2-shifted-local-derivative}
 \begin{split}
 |\partial_{12}e_j^\zeta(X;F,H)|
 &\le \frac{8C_E\lambda_{j-1}|\zeta|}{r_Jr_1r_2}
             e^{-\kappa d_j(X)}\\
 &\quad{}\times
 \exp\left[-\frac{\kappa_\sigma}{L^j}
     \max\{\operatorname{dist}(X,\operatorname{supp}F),
            \operatorname{dist}(X,\operatorname{supp}H)\}\right].
 \end{split}
\end{equation}
The distances in this local display are in the original
lattice units. Their physical scale is \(\lambda_j\).

The energy-count estimate of \cite[Prop. 10.16 and its proof]{QCDM1} sums the last two
factors over polymers. At \(\lambda_j\le d\), each support
lies in a radius-\(d\) ball and
\(\#^+_j(\operatorname{supp}F_\nu)\le1728(d/\lambda_j)^3\).
Separation retains \(e^{-\kappa_Ed/\lambda_j}\), where
$\kappa_E>0$ is the separation rate in the cited counting
estimate; it is chosen at most
$\min\{\kappa_\sigma/8,\kappa/(4c_{\mathcal D}M)\}$.
Consequently the level contribution to the mean is at most
\[
 \frac{8\cdot1728C_{\mathscr C}\lambda_{j-1}|\zeta|}
                 {r_Jr_1r_2}
       (d/\lambda_j)^3e^{-\kappa_Ed/\lambda_j}.
\]
At \(\lambda_j>d\), give the pairs independent signs
\(\epsilon_\nu\). The two sums
\(F_\epsilon=\sum_\nu\epsilon_\nu F_\nu\) and
\(H_\epsilon=\sum_\nu\epsilon_\nu H_\nu\) retain the norms
and disjointness required in
\eqref{eq:qcdm2-shifted-local-derivative}.
The sign average of any bilinear form at these two sums is
its sum over the diagonal pairs. The total number of
level cubes is bounded by
\[
 (\ell/\lambda_j)^3
       \le96^3N(d/\lambda_j)^3.
\]
The mean contribution of this level is therefore at most
\[
 \frac{8\cdot96^3C_{\mathscr C}\lambda_{j-1}|\zeta|}
                 {r_Jr_1r_2}(d/\lambda_j)^3.
\]
This use of the total cube count avoids multiplying the
coarse-scale bound by \(N\).

For bilinears,
\[
 (r_1r_2)^{-1}
   =4\lambda_{j-1}^2r_J^{-2}
                       \|f\|_\infty\|h\|_\infty,\qquad
 \lambda_{j-1}^3(d/\lambda_j)^3=d^3/L^3.
\]
There are at most \(H_d\) coarse levels. The fine levels
have distinct integers \(p\ge0\) with
\(d/\lambda_j\in[L^p,L^{p+1})\); their exponentials have a
bounded sum. For example one may take
\[
 C'_\kappa=\sum_{p\ge0}L^{3p}e^{-\kappa_E L^p},\qquad
 A_B=\frac{2^{25}C_{\mathscr C}(1+C'_\kappa)}{L^3r_J^3}.
\]

For action sources let
\(\pi_{j-1}=(1+(K-j+1)/n_\pi)^{p_0}\), the weight
in the level source norm. The test-pair Lipschitz bound gives
\[
 \|f\|_{\sharp,j-1}\le
 \frac{40M\pi_{j-1}}{\vartheta}
       \max\{1,\lambda_{j-1}/d\}\|f\|_\infty.
\]
At the coarse levels,
\[
 \pi_{j-1}\le(3+\log_L(1/d))^{p_0},\qquad
 \lambda_{j-1}(d/\lambda_j)^3
        \max\{1,\lambda_{j-1}/d\}^2\le d/L.
\]
At the fine level indexed by \(p\), the corresponding bounds
are \(((3+\log_L(1/d))(1+p))^{p_0}\) and
\(dL^{2p+1}\). Thus the fine-scale sum is paid by
\[
 C''_\kappa=\sum_{p\ge0}
           L^{2p}(1+p)^{2p_0}e^{-\kappa_E L^p}.
\]
One may take
\[
 A_0=\frac{2^{21}9^{p_0}C_{\mathscr C}
                         (1+LC''_\kappa)}{r_J}.
\]
Substitution proves the two estimates.
\end{proof}

The next estimate keeps the partition function normalized at
the shifted real mass. Its loss is proportional to the volume
of an occupied cell, rather than to the volume of the torus.

\begin{lemma}[A cell at a shifted real mass]\label{C-404-75}
Let \(x_0\in I\), \(|x-x_0|<2\varrho_J\), and \(x\) be real.
Put \(q=q_x^{(x_0)}\), \(E=E_K^{(x_0)}\), and
\[
 u(x)=-\ell^{-3}\log\bigl(Z_K(x)e^{-E(q)}\bigr).
\]
Consider a reflection board with \(\mathsf N\) cells of
physical volume \(|\square|\). Suppose \(s=(w,J)\) is
supported in one cell, at distance at least \(\delta>0\)
from its walls, and
\[
 \|w\|_\sharp<2\varrho,\qquad
 \|J\|_\infty\le2\varrho_J-|x-x_0|.
\]
If \(r_{\rm lin}\) denotes the linear part of each board
reflection on sources, set
\[
 \lambda_x(s)=\frac1{\mathsf N}\sum_r
                    \{E(q+r_{\rm lin}s)-E(q)\}.
\]
Then
\[
 \left|\frac{Z[q+s]}{Z[q]}\right|
 \le\exp\left[\operatorname{Re}\lambda_x(s)
       +(C+2+u(x))|\square|
       +\mathscr C_K(\operatorname{supp}s;\delta)\right].
\]
For a sum \(s=\sum_\square s_\square\) satisfying these
conditions in its respective cells, the right side is the
product of the displayed factors for the occupied cells.
\end{lemma}

\begin{proof}
The constant scalar source \(q\) is real and is fixed by
reflection. Its exponential is the product of the even
factors \(e^{\Phi(q1_\square)}\) in the cell algebras.
Apply the chessboard inequality at the original positive
mass \(x_0\), using
\(e^{\Phi(q1_\square)}e^{O(w)+\Phi(J)}\) in the occupied
cell and \(e^{\Phi(q1_\square)}\) elsewhere.
Folding an unoccupied cell gives \(e^{\Phi(q)}\).
Folding the occupied cell gives the source \(q+s^{\rm ch}\),
where the images of \(s\) have conjugated coefficients after
an odd number of reflections. Hence
\[
 \left|\frac{Z[q+s]}{Z[q]}\right|
       \le\left(\frac{Z[q+s^{\rm ch}]}{Z[q]}\right)^{1/\mathsf N}.
\]
The folded numerator is nonnegative. Its source lies in
\(\mathcal P\): folding preserves the norm of \(w\), and
the matrix norm is at most \(|x-x_0|+\|J\|_\infty\).
The history bound and base-source almost additivity give
\[
 \begin{split}
 Z[q+s^{\rm ch}]
 &\le \exp\bigl[\operatorname{Re}E(q+s^{\rm ch})
                                      +(C+2)\ell^3\bigr],\\
 \left|E(q+s^{\rm ch})-E(q)
       -\sum_r\{E(q+r_\theta s)-E(q)\}\right|
 &\le\mathsf N\mathscr C_K(\operatorname{supp}s;\delta).
 \end{split}
\]
The reflected supports have mutual distance at least
\(2\delta\). The reflected cube count costs at most \(27\);
this factor is included in \(\mathscr C_K\), as in the
one-cell estimate of \cite[Prop. 10.17]{QCDM1}.
Since \(q^\star=q\) and
\(E(v^\star)=\overline{E(v)}\), the real part of the sum
is \(\mathsf N\operatorname{Re}\lambda_x(s)\).
Use \(Z[q]=e^{E(q)-u(x)\ell^3}\) and
\(\ell^3/\mathsf N=|\square|\), then take the
\(\mathsf N\)-th root.

For several occupied cells use their respective factors in
the same chessboard inequality. Each cell contributes
\((Z[q]/Z[0])^{1/\mathsf N}\); their product is
\(Z[q]/Z[0]\). For an occupied cell its remaining factor
is bounded by the single-cell computation. Dividing out
the common product proves the multiple-cell bound.
\end{proof}

\begin{proposition}[The two kernels at the continued masses]
\label{C-404-76}\label{C-404-77}
Let \(c_2=9/(16\pi^2)\), \(c_B=1/(2\pi^2)\) and fix
\(\epsilon,\vartheta\) as above. Require the test pair
$(f,h)$ to be of class $C^1$. There are functions
\(\epsilon_0(d),\epsilon_B(d)\), tending to zero as
\(d\downarrow0\), and constants \(D_0,D_B\) such that
every family in \cref{C-404-60} satisfies, on
\(\Omega_\epsilon\),
\[
 \begin{split}
 \mathcal S_2(z;(0,f),(0,h))
       &=c_2b^{-2}\mathcal K_6(f,h)(1+\eta_0),\\
 \mathcal S_2(z;(\Gamma,f),(\Gamma',h))
       &=c_B\mathcal K_4(f,h)(\delta_{\Gamma\Gamma'}+\eta_B),
 \end{split}
\]
where
\[
 \begin{split}
 |\eta_0|&\le\epsilon_0(d)+
       D_0\{d+\operatorname{dist}(z,I)dH_d^{1+2p_0}\},\\
 |\eta_B|&\le\epsilon_B(d)+
       D_B\{d+\operatorname{dist}(z,I)dH_d\}.
 \end{split}
\]
The functions are uniform in the admitted side and mass
centre; the displayed constants may depend on the side.
At real \(z\in\Omega_\epsilon\), the same conclusion holds
with constants independent of the side.
In particular the short-distance coefficient of the action
two-point function determines \(b^{-2}\) throughout
\(\Omega_\epsilon\).
\end{proposition}

\begin{proof}
We prove uniform finite-cutoff bounds and then pass to the
subsequence. Choose the centre \(x_0\in I\) nearest to \(z\)
and put \(\zeta=z-x_0\). On \(\Omega_\epsilon\),
\(|\zeta|<R_\epsilon\).
For a source associated to one entry define
\[
 \|F\|_\circ^\epsilon=
 \begin{cases}
    \|f\|_\sharp/\varrho,&\alpha=0,\\
    \|f\|_\infty/\epsilon,&\alpha\in\mathfrak B.
 \end{cases}
\]
For the test pair,
\(\|f\|_\sharp\le40M\|f\|_\infty/(\vartheta d)\).
On \(|t_i|<2/\|F_i\|_\circ^\epsilon\), the source
\(q+t_1F+t_2H\) lies in \(\mathcal P\):
the extra bilinear norm is less than \(2\epsilon\), while
\(|\zeta|+2\epsilon<2\varrho_J\).

Write \(E=E_K^{(x_0)}\). At a nonzero partition function,
\[
 S^T_{2,K}(z;F,H)
   =\partial_{12}E(q;F,H)+\mathcal B_\zeta(F,H),\qquad
 \mathcal B_\zeta(F,H)
   =\partial_{12}\log W_K[q+t_1F+t_2H]\big|_0 .
\]
Both terms are bilinear in \(F,H\).
The zero-free estimate, the finite-size normalization and
the history upper bound imply
\[
 \left|\frac{W_K[q+t_1F+t_2H]}
                    {W_K[q]}\right|
       \le B_c,\qquad
 B_c=2e^{(C+2+A_\epsilon)\ell^3+2C_{\rm fs}}.
\]
For a holomorphic function \(\psi\) on the indicated bidisc
with \(\psi(0)=1\) and \(|\psi|\le B\), Schwarz's lemma
on complex lines gives \(|\psi-1|\le1/2\) on the bidisc
smaller by \(1/(4B)\). Its logarithm and Cauchy's formula
therefore give
\[
 |\partial_{12}\log\psi(0)|
        \le4B^2\|F\|_\circ^\epsilon\|H\|_\circ^\epsilon.
\]
This applies to \(\mathcal B_\zeta\) with \(B=B_c\).

Take a torus grid with mesh at least \(12d\), rounded to
the lattice. For \(8a\le d\), the radius-\(5d\) balls
containing its translated pairs are at distance at least
\(d\). There are \(N\ge(\ell/(24d))^3\) such pairs.
Translation invariance holds at complex mass: the finite
integrals and their insertion derivatives are polynomials
in \(z\), so the real-mass identity extends as an identity
of rational functions. The sign average used in the proof
of \cref{C-404-74} now yields
\begin{equation}\label{eq:qcdm2-mass-packing-complex}
 S^T_{2,K}(z;F,H)=Q_K(\zeta)+\mathcal R_c,\qquad
 |\mathcal R_c|\le4B_c^2(24d/\ell)^3
                   \|F\|_\circ^\epsilon\|H\|_\circ^\epsilon .
\end{equation}

For real \(x\) with \(\operatorname{dist}(x,I)\le R_\epsilon\)
one can replace the volume loss by a cell loss.
The lower bound of \cref{C-404-70} gives \(u(x)\le A_\epsilon\).
Translate the pair into the radius-\(1/4\) ball of a unit
cell. The cell packing of \cite[Lem. 11.2 and its proof]{QCDM1} has at least \((96d)^{-3}\) translates
inside this ball. For their two sums with signs,
\cref{C-404-75} applies with margin \(1/4\) and gives
\[
 \left|e^{-\lambda_x(t_1F+t_2H)}
          \frac{Z[q+t_1F+t_2H]}{Z[q]}\right|
       \le B_r,\qquad
 B_r=e^{C+2+A_\epsilon+\mathscr C_*},
\]
where \(\mathscr C_*\) is the fixed upper bound for
\(\mathscr C_K(S;1/4)\) when \(S\) lies in that ball.
The logarithmic derivative subtracts the board average
of \(\partial_{12}E\). Averaging the signs retains its
diagonal pairs, now including their images in all
\(\ell^3\) cells. Thus
\begin{equation}\label{eq:qcdm2-mass-packing-real}
 S^T_{2,K}(x;F,H)=Q_K(x-x_0)+\mathcal R_r,\qquad
 |\mathcal R_r|\le4B_r^2(96d)^3
                   \|F\|_\circ^\epsilon\|H\|_\circ^\epsilon.
\end{equation}
Here \(N\ge(\ell/(96d))^3\), so \cref{C-404-74}
applies to both packings.

For clarity, the estimate of \(Q_K(0)\) uses the free-template and interaction-error estimates in
\cite[proofs of Lems. 11.5 and 11.6]{QCDM1}, before adding
the partition-function remainder in its two-point theorem.
Their proofs apply to either packing just used: at fine
levels they use the separated pair alone, and at coarse
levels they use precisely
\((\ell/\lambda_j)^3\le96^3N(d/\lambda_j)^3\).
The fourth-order source term has zero mixed derivative at
these indices. The free terms give \(c_2b^{-2}\mathcal K_6\)
and \(c_B\delta_{\Gamma\Gamma'}\mathcal K_4\).
After division by
\eqref{eq:qcdm2-continued-template-floor}, the interaction
errors are bounded by constants times
\(\phi(b/d)(1+\log(b/d))^{2p_0}\) and
\(\phi(b/d)\), respectively. The free remainders contain
powers of \(d\), \(d/\ell\), and the positive mass bounds,
and hence tend to zero uniformly for admitted \(\ell\).
Their mesh terms tend to zero for each fixed test pair
as \(K\to\infty\). Denote the resulting cutoff-independent
remainders by \(\epsilon_0(d),\epsilon_B(d)\).

Add the change \(Q_K(\zeta)-Q_K(0)\) from
\cref{C-404-74} and the appropriate packing remainder.
For the action term the denominator is at least
\[
 c_2\,10^{-6}\vartheta^2X_g^{-2}
                      \|f\|_\infty\|h\|_\infty;
\]
for the bilinear term it is at least
\(c_B10^{-4}\vartheta^2d^2\|f\|_\infty\|h\|_\infty\).
The mass-shift terms therefore have the displayed orders
\(|\zeta|dH_d^{1+2p_0}\) and \(|\zeta|dH_d\).
Each packing remainder has relative order \(d\).
Its constant depends on \(\ell\) in
\eqref{eq:qcdm2-mass-packing-complex}, and is independent
of \(\ell\) in \eqref{eq:qcdm2-mass-packing-real}.
Absorb these explicit factors into \(D_0,D_B\).

Finally take the cutoff subsequence defining the holomorphic
family. The tests are separated and of class \(C^1\), so the
cumulants converge; all mesh remainders vanish. The bounds
survive and their errors tend to zero as \(d\downarrow0\).
The coefficient conclusion follows after dividing the
first formula by \(c_2\mathcal K_6\).
\end{proof}

\begin{proposition}[Uniform real-mass cumulant bounds]
\label{C-404-78}
Let \(n\ge2\), and let the real Lipschitz tests \(f_i\)
have pairwise distance at least \(d>0\), each support lying
in a ball of radius \(R\). For
\[
 K\ge K_\epsilon,\qquad L^K\ge16n,\qquad
 x\in\mathbb R,\quad\operatorname{dist}(x,I)\le R_\epsilon,
\]
one has
\[
 |S^T_{n,K}(x;(\alpha_i,f_i)_i)|
 \le(1+d^{-1})^3 B_\epsilon(R)^n(n!)^{5+p/2}
                                      \prod_i\|f_i\|_{\alpha_i},
\]
where
\[
 B_\epsilon(R)=
 \frac{C_Be^{C+2+A_\epsilon}(1+R)^3}{\min\{\varrho,\epsilon\}},
 \qquad p=\frac{432C_{\mathscr C}}{\log L}.
\]
The constants \(C_B,C_{\mathscr C}\) are those of the
separated-cumulant bound in \cite[Thm. F]{QCDM1}.
This bound is independent of the side, cutoff, coupling
within \(I_g\), and mass in the displayed interval.
For tests in the continuum class specified at the start of this
section, it passes to every continued family at its real masses.
\end{proposition}

\begin{proof}
Choose \(x_0\in I\) nearest to \(x\), and use
\(q=q_x^{(x_0)}\). Then \(Z[q]>0\) and \(u(x)\le A_\epsilon\).
Set \(\eta=(16n)^{-1}\). The product hat partition of unity
at mesh \(\eta\) splits each \(f_i\) into at most
\((2R/\eta+3)^3\) pieces \(f_i\chi_a\).
Each piece is supported in a cube of side \(2\eta\), with
insertion norm at most
\((1+80M/\eta)\|f_i\|_{\alpha_i}\).

Fix one piece from every insertion. The board selection in
the proof of \cite[Thm. F]{QCDM1} leaves distance at least \(\eta\) between these
supports and the walls of unit cells. To recall why the
choice is available, in each coordinate a support forbids
at most \(4\eta L^K+1\) half-integral wall positions.
The \(n\) supports therefore forbid fewer than \(L^K\)
positions because \(L^K\ge16n\).
The resulting board has an even number of slabs on every
admitted torus.

Write \(s_i\) for the source of the chosen piece and choose
radii
\[
 \rho_i=
 \begin{cases}
 2\varrho/\|f_i\chi_{a_i}\|_\sharp,&\alpha_i=0,\\
 2\epsilon/\|f_i\chi_{a_i}\|_\infty,&\alpha_i\in\mathfrak B.
 \end{cases}
\]
Group \(s(z)=\sum_i z_is_i\) by its cells, and put
\[
 \Lambda_x(z)=\sum_\square\lambda_x(s_\square(z)),\qquad
 G(z)=e^{-\Lambda_x(z)}Z[q+s(z)]/Z[q].
\]
Disjoint supports ensure that every group has action norm
less than \(2\varrho\) and matrix norm less than
\(2\epsilon\le2\varrho_J-|x-x_0|\).
For a coordinate face indexed by \(B\), at most \(|B|\)
cells are occupied. \Cref{C-404-75} and subadditivity of
the support count give
\[
 |G(z_B)|\le
 \exp\left[|B|(C+2+A_\epsilon)
               +\sum_{i\in B}\mathscr C_K(\operatorname{supp}s_i;\eta)\right].
\]
The refined cube count from \cite[proof of Thm. F]{QCDM1} gives
\[
 \mathscr C_K(\operatorname{supp}s_i;\eta)
       \le216C_{\mathscr C}(1+\log_L(1/\eta))+C.
\]
Cauchy's formula on each face consequently bounds
\(\partial_BG(0)\) by this exponential times
\(\prod_{i\in B}\rho_i^{-1}\).

Differentiate \(\log Z=\Lambda_x+\log G+\log Z[q]\):
\[
 S^T_{n,K}
   =\partial_{[n]}\Lambda_x+
     \sum_\pi(-1)^{|\pi|-1}(|\pi|-1)!
                           \prod_{B\in\pi}\partial_BG(0).
\]
The sum of the absolute partition coefficients is at most
\(2^nn!\). Thus its contribution is bounded by
\[
 C^ne^{n(C+2+A_\epsilon)}
        (16n)^{np/2}2^nn!\prod_i\rho_i^{-1}.
\]
The full mixed derivative of \(\Lambda_x\) is zero unless
all the pieces occupy one cell. In that case its reflected
energy derivatives obey \cite[Prop. 10.18(i)]{QCDM1}, since every partial source remains in
\(\mathcal P\). It bounds this term by
\[
 2^{n-1}\prod_i\rho_i^{-1}
       \sup_Q\mathscr C_K(Q;d)
 \le C^n(1+d^{-1})^3\prod_i\rho_i^{-1},
\]
where \(Q\) ranges over cubes of side \(2\eta\).
The reflected count, including its factor \(27\), is the
one already used in \cref{C-404-75}.

Sum over the choices of pieces. Their number and norm
enlargement cost at most
\[
 (32nR+3)^{3n}(1+1280Mn)^n.
\]
The inequality \(n^n\le e^nn!\) turns these four powers
of \(n\) into \((n!)^4\). The partition sum contributes
one further factorial and the logarithmic count contributes
\((n!)^{p/2}\). Absorbing fixed constants gives the stated
\(B_\epsilon(R)\). No estimate used the torus side.
For the stated continuum tests, passing along the cutoff
subsequence proves the last assertion; no one-point limit is involved.
\end{proof}

\appendix
\section{Index of notation and conventions}
\label{app:qcdm2-notation}

The table lists the principal notation in the order in which it
is used. A cited result includes the definitions immediately
preceding it.

\begingroup
\small
\renewcommand{\arraystretch}{1.1}
\begin{longtable}{@{}p{0.24\textwidth}@{\hspace{0.025\textwidth}}p{0.735\textwidth}@{}}
\textbf{Notation} & \textbf{Meaning and location}\\
\hline
\endfirsthead
\textbf{Notation} & \textbf{Meaning and location (continued)}\\
\hline
\endhead
\endfoot
\endlastfoot
$a=L^{-K}$, $\lambda_j$, $\beta$, $m_0$
& The fine spacing, the level spacing $\lambda_j=L^{j-K}$,
and the bare parameters $\beta=\mathsf b/a$, $m_0=a\mathsf m$;
\cref{sec:qcdm2-setting}.\\
$I_g$, $\mathfrak L(g)$, $T_\ell$
& The admitted interval for $\mathsf b$, the admitted torus
sides, and $(\mathbb R/\ell\mathbb Z)^3$;
\cref{sec:qcdm2-setting}.\\
$D[U]$, $D_\pm$, $e_\pm$, $\mathsf G_\pm$
& The massless chamber operator, its two internal blocks,
their projections $(1\pm\tau_3)/2$, and the physical
resolvents $a(D_\pm+m_0)^{-1}$;
\eqref{eq:qcdm2-Dirac}.\\
$\mathfrak B$, $O_K(w)$, $\Phi_K(J)$
& The reflection-real meson basis, the action insertion,
and the site bilinear insertion;
\eqref{eq:qcdm2-action}--\eqref{eq:qcdm2-meson}.
The label $0$ denotes the action density.\\
$\|f\|_\alpha$, $\|f\|_\sharp$
& $\|f\|_0=\|f\|_\sharp
=\max(\|f\|_\infty,40M\operatorname{Lip}f)$ and
$\|f\|_\Gamma=\|f\|_\infty$ for $\Gamma\in\mathfrak B$;
\cref{sec:qcdm2-setting}.\\
$\nu_K$, $\mathcal G^{\rm tor}_{K,\beta,m_0}$
& The finite-torus gauge law and the set of its infinite-lattice
torus limits; \cref{sec:qcdm2-circle}.\\
$\mathsf L(f)$, $\Xi_\nu^f(z)$
& The loop $\operatorname{Tr}(\mathsf G_+\hat f)$ and its
determinant moment; \eqref{eq:qcdm2-Xi}.
The hat denotes multiplication by the test function.\\
$D^{\rm mix}$, $E_K^{\rm mix}$, $h_K$
& Mixed extraction by inclusion--exclusion, the resulting
energy in \eqref{C-415-2}, and the holomorphic
normalization in the circle bound
\eqref{eq:qcdm2-pure-energy}.\\
$F_s$, $Z_K(s)$, $\nu_{K,s}$, $\mathfrak l_s$
& The determinant at physical mass $s$, its partition function,
its gauge law, and
$\mathfrak l_s=\partial_s\log|F_s|^2$;
\eqref{eq:qcdm2-mass-partition}--\eqref{eq:qcdm2-mass-loop}.\\
$R_b$, $R^{\rm mix}$, $P_{\rm col}$
& Deletion of species $b$, the mixed projection
$(1-R_1)(1-R_2)$, and independent constant colour averaging
in the two species at the trivial pair; \cref{C-415-6}.\\
$\mathsf t_j$, $x_j$, $\phi_j$, $\delta_j$, $\mathfrak l_j$
& The calibration coefficient $\lambda_j^{-1}\mathsf b$,
the running coefficient, $\phi(x_j)$, and the gauge fluctuation
radius; \cref{sec:qcdm2-mixed-proof}. The additional scale weight is
$\mathfrak l_j=\mathsf t_j^{-1/2}(\log\mathsf t_j)^{-p_0}$;
\cref{C-415-5}.\\
$\mathfrak c$, $\mathscr D_k^\circ$, $\sigma_i$
& The free coordinates $\mathfrak c=(0,\mathsf n,(\sigma_i)_{i\le k})$,
their domain at zero gauge source,
and the extracted mass-source corrections;
\eqref{eq:qcdm2-zero-source-domain}.\\
$H_k$, $\bar H_k$, $\Gamma_k^{-1}$, $W$
& The unbarred background map, its row twin, the fermion
fluctuation covariance and the fermion fluctuation;
\cref{sec:qcdm2-mixed-proof}. The gauge fluctuation is $B$.\\
$\epsilon_F$, $\|m_c^{\rm tor}\|_{\rm nat}$
& The weighted norm of the chosen local decomposition of $F$,
and the norm of its grouped cell zeroth moments;
\cref{C-415-9}.\\
$\mathscr S$, $M^*$, $\mathscr S_{\mathbb R}$
& The enlarged bilinear label space, the involution
$M^*=\tau_3M^\dagger\tau_3$, and its real form;
\eqref{C-401-2}. The insertion with label $M$ is $\Phi[M](f)$.\\
$G_{\mathbb C}$, $G$, $V(S)$
& $\mathrm{Sp}(2N_f,\mathbb C)^2$, its compact form
$\mathrm{USp}(2N_f)^2$, and the unitary action of $S\in G$;
\eqref{C-401-3} and \cref{C-401-46}.\\
$S^{\rm ext}$, $\sigma^{\rm ext}$, $\mathcal H^{\rm ext}$,
$H$, $P_1,P_2$
& The extended truncated functions, their centered functional,
the reflection Hilbert space, and its time and spatial
translation generators; \cref{C-401-49}.\\
$n$, $Q$, $\mathcal H_n^{\rm ext}$, $\mathcal H_\pi^{\rm ext}$
& A species weight, the quark-number operator $Q$ with eigenvalue
$q=\sum_fn_f$,
the corresponding weight space, and the $\pi$-isotypic space.
Here $|\pi|$ is the largest $\ell^1$ norm of a weight of $\pi$;
\cref{C-401-46,C-400-34}.\\
$\mathsf j^T$, $j_\mu^T$, $J^T(f)$
& The chamber-link current, its sum across a site bond,
and its smeared physical insertion, including its factor $i$;
\eqref{C-402-3}.\\
$\mathcal M$, $\mathsf n$, $\mathscr O_k^{\mathcal M}$,
$\mathcal P^{\mathcal M}$, $\mathcal O_j^{{\rm lev},\mathcal M}$
& The matrices acting on spin, internal and flavour indices
and trivially on colour; the bare source $\mathsf n=-aJ$;
and the chain, physical and level source domains;
the definitions preceding \cref{C-417-7}.\\
$\mathcal D_j$, $d_j(Y)$, $E_K(s)$, $e_j(Y;s)$, $W_K[s]$
& The polymers and their block-tree lengths, the extracted
energy, its local increments, and the remaining factor in
$Z_K[s]=e^{E_K(s)}W_K[s]$;
\eqref{eq:qcdm2-energy-decomposition}.\\
$\mathcal P$, $\rho_b$, $\rho_m$, $W_{\mathbb C}$
& The physical source ball and the complex coupling--mass
neighbourhood of the real parameter window;
\cref{sec:analytic-labels} and
\eqref{eq:qcdm2-analytic-radii}.\\
$d_\cap(S_1,\ldots,S_n)$
& $\min_{x\in T_\ell}\max_i\operatorname{dist}(x,S_i)$.
It is positive exactly when the closed supports have empty
total intersection; the definition preceding \cref{C-412-53}.\\
$I$, $U_r$, $\mathcal E_K(z)$, $h_K(z)$
& The mass interval, its complex $r$-neighbourhood, the common
vacuum normalization, and $h_K(z)=Z_K(z)e^{-\mathcal E_K(z)}$;
\cref{sec:qcdm2-mass-continuation,C-404-67}.\\
$I_\epsilon$, $\rho_\ell$, $w_\ell$, $\Omega_\epsilon$
& The longer real mass interval, the radius used in its
zero-free estimate, and its complex neighbourhood of width
$w_\ell=\rho_\ell/[2(e+1)]$;
\eqref{eq:qcdm2-mass-width} and \cref{C-404-72}.\\
$h$, $\mathcal Z$, $\mathcal S_n(z;X)$
& A subsequential limit of the normalized partition functions,
its zero set, and the common meromorphic continuation of the
separated cumulants; \cref{C-404-60}.\\
\end{longtable}
\endgroup

\subsection*{Conventions}

The barred and unbarred variables are independent Grassmann
generators. Haar measure has mass one, and the Berezin integral
of $e^{-\bar\psi A\psi}$ is $\det A$. The traces defining
$\mathsf L$ and $\mathfrak l_s$ include position, spin and colour
but no flavour; when a source acts on flavour, the propagator
is extended by the identity on that factor.

Operator formulas use lattice coordinates; supports and distances
are physical unless stated otherwise. Links on $T_\ell$ are
periodic and spinors are antiperiodic. Infinite-volume results
take volume to infinity before removing the cutoff. The complex
parameter results keep $T_\ell$ fixed.

The symbol $M^*$ on $\mathscr S$ denotes the displayed reflection
involution. Elsewhere $*$ on an operator denotes its Hilbert-space
adjoint. The functions called $h_K$ in the circle bound and in
mass continuation have the separate definitions listed above.
The former subtracts a one-species energy; the latter is a
normalized partition function.
Likewise $\mathfrak l_s$ in \cref{sec:qcdm2-mass-comparison}
is a random loop at physical mass $s$, whereas $\mathfrak l_j$
in the scale estimates is the deterministic weight listed above.
In \cref{sec:qcdm2-currents},
$G_U=(D_+[U]+m_0)^{-1}=a^{-1}\mathsf G_+$ denotes the
lattice propagator; $\mathsf G_+$ includes the physical factor $a$.

\begin{thebibliography}{QCDM1}
\bibitem[DM22]{DM22}
P.~Dupuis and Y.~Mao,
\emph{Formulation and properties of a divergence used to compare
probability measures without absolute continuity},
ESAIM Control Optim. Calc. Var. \textbf{28} (2022), article 10.
\href{https://doi.org/10.1051/cocv/2022002}{doi:10.1051/cocv/2022002}.

\bibitem[QCDM1]{QCDM1}
\emph{Massive SU(2) gauge theory in three dimensions},
companion manuscript, revised edition of 1 November 2026.
The accompanying \href{QCD3_M1.pdf}{QCD3\_M1.pdf} is the edition cited here.

\bibitem[YMM1]{YMM1}
\emph{Scalar fields from three-dimensional lattice Yang--Mills theory},
companion manuscript, corrected edition of 26 October 2026.
\href{https://github.com/ethanwsussman/Yang--Mills-LLM-experiment/blob/main/YM3/Revisions/YMM1-2026-10-26.pdf}{YMM1.pdf}.
\end{thebibliography}
\end{document}
